\documentclass[a4paper,10pt]{article}
\usepackage{amsmath,amssymb}
\usepackage{fontspec}
\setmainfont{Liberation Serif}
\usepackage{graphicx}
\usepackage{eso-pic}
\usepackage{tikz}
\usepackage[most]{tcolorbox}
\usepackage{geometry}
\usepackage{fancyhdr}
\usepackage{lastpage}
\usepackage{booktabs}
\usepackage{listings}
\usepackage{enumitem}
\usepackage{microtype}

\geometry{
    a4paper,
    top=1.5cm,
    bottom=1.8cm,
    left=1.2cm,
    right=1.2cm
}

% Persistent Unremovable Canvas Watermark on Every Page
\newsavebox{\watermarkbox}
\sbox{\watermarkbox}{%
    \begin{tabular}{c}
        \includegraphics[width=8cm]{/home/johan/GATE/GATE_CS_PYQ_DB/assets/studiva-butterfly-icon.jpg} \\[-0.5em]
        {\fontsize{18}{22}\selectfont\bfseries\scshape www.studiva.co.in}
    \end{tabular}%
}
\AddToShipoutPictureBG{%
    \begin{tikzpicture}[remember picture, overlay]
        \node[opacity=0.12] at (current page.center) {\usebox{\watermarkbox}};
    \end{tikzpicture}%
}

% Header and Footer
\pagestyle{fancy}
\fancyhf{}
\fancyhead[L]{\small\scshape GATE 2018--2025 CS/IT $\cdot$ \textbf{Numerical Answer Type (NAT)}}
\fancyhead[R]{\small\scshape Practice Question Bank with Solutions}
\fancyfoot[L]{\scriptsize\textsc{Copyright \copyright\ 2026} Studiva\textsuperscript{TM}}
\fancyfoot[C]{\small Page \thepage\ of \pageref{LastPage}}
\fancyfoot[R]{\scriptsize\texttt{www.studiva.co.in}}
\renewcommand{\headrulewidth}{0.4pt}
\renewcommand{\footrulewidth}{0.4pt}

% Code Listing Style (Monochrome)
\lstset{
    basicstyle=\footnotesize\ttfamily,
    numbers=none,
    frame=none,
    tabsize=2,
    breaklines=true,
    showstringspaces=false
}

\begin{document}

% ==================== PAGE 1: IMPRINT & TITLE PAGE ====================
\thispagestyle{empty}
\vspace*{0.5cm}
\begin{center}
    \includegraphics[width=2.5cm]{/home/johan/GATE/GATE_CS_PYQ_DB/assets/studiva-butterfly-icon.jpg}\\[0.6em]
    {\LARGE \textbf{GATE 2018--2025 COMPUTER SCIENCE \& IT}}\\[0.3em]
    {\Large \textbf{PREVIOUS YEAR QUESTION BANK WITH SOLUTIONS}}\\[0.4em]
    {\large \textsc{\textbf{Question Type Wise: Numerical Answer Type (NAT)}}}\\[0.2em]
    {\normalsize \textsc{Combined 2018--2025 Examination Papers $\cdot$ Comprehensive Edition ($\mathbf{116}$ Questions)}}\\[1.2em]
    \rule{0.6\textwidth}{0.6pt}
\end{center}

\vspace{1em}
\noindent\begin{tabular}{ll}
\textbf{Author / Compiler:} & Rupam Ghosh \\
\textbf{Publisher \& Platform:} & Studiva\textsuperscript{TM} (\texttt{www.studiva.co.in}) \\
\textbf{Official Inquiries:} & \texttt{contact@studiva.co.in} \\
\textbf{Edition:} & First Standard Edition (2026) \\
\textbf{Coverage:} & Complete GATE CS/IT Papers (2018, 2022, 2023, 2024 CS1, 2024 CS2, 2025 CS1, 2025 CS2) \\
\textbf{Syllabus Reference:} & Official GATE 2026 Syllabus (Organizing Institute: IIT Guwahati) \\
\end{tabular}

\vspace{1.5em}
\begin{tcolorbox}[enhanced, standard jigsaw, opacityback=0, colframe=black!40, arc=1mm, boxrule=0.5pt, left=10pt, right=10pt, top=8pt, bottom=8pt]
\textbf{\large COPYRIGHT \& INTELLECTUAL PROPERTY NOTICE}\\[0.6em]
Copyright \copyright\ 2026 Rupam Ghosh \& Studiva\textsuperscript{TM}. All Rights Reserved.\\[0.5em]
This educational compilation, problem organization, taxonomy classification, LaTeX typesetting, hints, and explanatory solutions are the proprietary intellectual property of \textbf{Rupam Ghosh} and \textbf{Studiva\textsuperscript{TM}}.\\[0.6em]
\textbf{Terms of Use \& Attribution Rules:}
\begin{itemize}[leftmargin=*, nosep]
    \item \textbf{Personal Educational Use:} Permitted freely for individual students and educators preparing for competitive examinations.
    \item \textbf{Mandatory Attribution:} Any reproduction, republication, hosting, or sharing (digital or physical) must prominently credit \textbf{Rupam Ghosh \& Studiva\textsuperscript{TM}} and include an active, direct link to \texttt{https://www.studiva.co.in}.
    \item \textbf{Commercial Restrictions:} Commercial resale, sublicensing, course bundling, or inclusion in paid test series without prior written permission is strictly prohibited.
    \item \textbf{Anti-Tampering Notice:} Removing, altering, or obscuring copyright notices, page watermarks, or website links from this document is strictly prohibited.
\end{itemize}
\vspace{0.6em}
\textbf{Fair Use Disclaimer:}\\
The underlying examination questions originate from the Graduate Aptitude Test in Engineering (GATE), conducted by the Indian Institutes of Technology (IITs) and Indian Institute of Science (IISc). Question texts are reproduced for non-commercial educational study under Fair Use doctrines. No copyright is claimed over the original GATE examination question texts.
\end{tcolorbox}

\vspace{1.2em}
\begin{center}
\small\itshape
Practice Recommendation: Use an opaque sheet of paper or ruler to cover the right-hand column during mock attempts. Tick your chosen answers directly in the check boxes on the left, then slide the cover to check the solution and hint.
\end{center}

\newpage
% ==================== PAGE 2 ONWARDS: QUESTION BANK ====================

\vspace{8pt}
\subsection*{Engineering Mathematics}
\vspace{-4pt}
\noindent
\begin{tcolorbox}[
    blanker,
    sidebyside,
    sidebyside align=top,
    lefthand width=0.55\textwidth,
    top=4pt, bottom=4pt, left=0pt, right=0pt,
    middle=16pt,
    underlay={
        \draw[dashed, black!80, line width=0.8pt, dash pattern=on 4pt off 3pt]
            ([xshift=0.55\textwidth+8pt]frame.north west) --
            ([xshift=0.55\textwidth+8pt]frame.south west);
    }
]
\textbf{Q.1} \quad {\footnotesize\textsc{[NAT $\cdot$ 1 Mark(s) $\cdot$ No Neg]}} \hfill {\footnotesize\textbf{GATE 2018 CS} \quad \tiny\texttt{2018\_CS\_GA\_Q15}}
\vspace{2pt}\hrule height 0.3pt\vspace{4pt}

Two people, $P$ and $Q$, decide to independently roll two identical dice, each with 6 faces, numbered 1 to 6. The person with the higher number wins. In case both get the same number, they repeat the roll until someone wins. The probability that the game concludes on the third roll is \_\_\_\_\_\_\. (Answer in 3 decimal places)

\par\vspace{6pt}\noindent\textbf{Your Answer:} \fbox{\rule{0pt}{1.3ex}\hspace{3cm}}

\tcblower

\textbf{Official Key:} \textbf{\large 0.021 to 0.024} \hfill \textit{\small [Medium]}
\vspace{2pt}\hrule height 0.3pt\vspace{3pt}
\textbf{Topic:} Probability and Statistics \\
\textbf{Subtopic:} \textit{Random variables}
\vspace{2pt}\hrule height 0.3pt\vspace{3pt}
\textbf{Guide / Hint:} \\
{\small On rolling two fair dice independently, the probability of a tie is $P(\text{tie}) = \frac{6}{36} = \frac{1}{6}$, and the probability of someone winning (no tie) is $1 - \frac{1}{6} = \frac{5}{6}$. The game ends on the third roll if trials 1 and 2 are ties and trial 3 produces a winner: $P(\text{Trial 3}) = \left(\frac{1}{6}\right) \times \left(\frac{1}{6}\right) \times \left(\frac{5}{6}\right) = \frac{5}{216} \approx 0.0231$.}

\end{tcolorbox}
\vspace{6pt}\hrule height 0.35pt\vspace{6pt}
\noindent
\begin{tcolorbox}[
    blanker,
    sidebyside,
    sidebyside align=top,
    lefthand width=0.55\textwidth,
    top=4pt, bottom=4pt, left=0pt, right=0pt,
    middle=16pt,
    underlay={
        \draw[dashed, black!80, line width=0.8pt, dash pattern=on 4pt off 3pt]
            ([xshift=0.55\textwidth+8pt]frame.north west) --
            ([xshift=0.55\textwidth+8pt]frame.south west);
    }
]
\textbf{Q.2} \quad {\footnotesize\textsc{[NAT $\cdot$ 1 Mark(s) $\cdot$ No Neg]}} \hfill {\footnotesize\textbf{GATE 2018 CS} \quad \tiny\texttt{2018\_CS\_GA\_Q16}}
\vspace{2pt}\hrule height 0.3pt\vspace{4pt}

The value of $\int_0^{\pi/4} x \cos(x^2) \, dx$ correct to three decimal places (assuming that $\pi = 3.14$) is \_\_\_\_\_\_\_\_.

\par\vspace{6pt}\noindent\textbf{Your Answer:} \fbox{\rule{0pt}{1.3ex}\hspace{3cm}}

\tcblower

\textbf{Official Key:} \textbf{\large 0.27 to 0.30} \hfill \textit{\small [Easy]}
\vspace{2pt}\hrule height 0.3pt\vspace{3pt}
\textbf{Topic:} Calculus \\
\textbf{Subtopic:} \textit{Integration}
\vspace{2pt}\hrule height 0.3pt\vspace{3pt}
\textbf{Guide / Hint:} \\
{\small Substitute $u = x^2$, so $du = 2x\,dx$. The limits transform from $x = 0 \implies u = 0$ to $x = \frac{\pi}{4} \implies u = \frac{\pi^2}{16} \approx \frac{9.8596}{16} \approx 0.61685$ radians. The integral is $\frac{1}{2}\int_0^{\pi^2/16} \cos u \, du = \frac{1}{2} \sin\left(\frac{\pi^2}{16}\right) \approx \frac{1}{2} \sin(0.61685) \approx \frac{1}{2}(0.5784) \approx 0.289$.}

\end{tcolorbox}
\vspace{6pt}\hrule height 0.35pt\vspace{6pt}
\noindent
\begin{tcolorbox}[
    blanker,
    sidebyside,
    sidebyside align=top,
    lefthand width=0.55\textwidth,
    top=4pt, bottom=4pt, left=0pt, right=0pt,
    middle=16pt,
    underlay={
        \draw[dashed, black!80, line width=0.8pt, dash pattern=on 4pt off 3pt]
            ([xshift=0.55\textwidth+8pt]frame.north west) --
            ([xshift=0.55\textwidth+8pt]frame.south west);
    }
]
\textbf{Q.3} \quad {\footnotesize\textsc{[NAT $\cdot$ 1 Mark(s) $\cdot$ No Neg]}} \hfill {\footnotesize\textbf{GATE 2018 CS} \quad \tiny\texttt{2018\_CS\_GA\_Q17}}
\vspace{2pt}\hrule height 0.3pt\vspace{4pt}

Consider a matrix $A = u v^T$ where $u = \begin{bmatrix} 1 \\ 2 \end{bmatrix}$ and $v = \begin{bmatrix} 1 \\ 1 \end{bmatrix}$. Note that $v^T$ denotes the transpose of $v$. The largest eigenvalue of $A$ is \_\_\_\_\_\_\_\_.

\par\vspace{6pt}\noindent\textbf{Your Answer:} \fbox{\rule{0pt}{1.3ex}\hspace{3cm}}

\tcblower

\textbf{Official Key:} \textbf{\large 3 to 3} \hfill \textit{\small [Easy]}
\vspace{2pt}\hrule height 0.3pt\vspace{3pt}
\textbf{Topic:} Linear Algebra \\
\textbf{Subtopic:} \textit{Eigenvalues and eigenvectors}
\vspace{2pt}\hrule height 0.3pt\vspace{3pt}
\textbf{Guide / Hint:} \\
{\small For an outer product matrix $A = u v^T$ where $u = \begin{bmatrix} 1 \\ 2 \end{bmatrix}$ and $v = \begin{bmatrix} 1 \\ 1 \end{bmatrix}$, the rank of $A$ is 1. Its non-zero eigenvalue equals the trace of $A$, which is given by the inner product $v^T u = 1(1) + 1(2) = 3$.}

\end{tcolorbox}
\vspace{6pt}\hrule height 0.35pt\vspace{6pt}
\noindent
\begin{tcolorbox}[
    blanker,
    sidebyside,
    sidebyside align=top,
    lefthand width=0.55\textwidth,
    top=4pt, bottom=4pt, left=0pt, right=0pt,
    middle=16pt,
    underlay={
        \draw[dashed, black!80, line width=0.8pt, dash pattern=on 4pt off 3pt]
            ([xshift=0.55\textwidth+8pt]frame.north west) --
            ([xshift=0.55\textwidth+8pt]frame.south west);
    }
]
\textbf{Q.4} \quad {\footnotesize\textsc{[NAT $\cdot$ 2 Mark(s) $\cdot$ No Neg]}} \hfill {\footnotesize\textbf{GATE 2018 CS} \quad \tiny\texttt{2018\_CS\_GA\_Q44}}
\vspace{2pt}\hrule height 0.3pt\vspace{4pt}

Consider Guwahati ($G$) and Delhi ($D$) whose temperatures can be classified as high ($H$), medium ($M$) and low ($L$). Let $P(H_G) = 0.2$, $P(M_G) = 0.5$, and $P(L_G) = 0.3$ be the probabilities that Guwahati has high, medium, and low temperatures, respectively. Given $P(H_D \mid H_G) = 0.40$, $P(H_D \mid M_G) = 0.10$, and $P(H_D \mid L_G) = 0.01$, the conditional probability $P(H_G \mid H_D)$ is \_\_\_\_\_\_\_\_. (Answer correct to two decimal places)

\par\vspace{6pt}\noindent\textbf{Your Answer:} \fbox{\rule{0pt}{1.3ex}\hspace{3cm}}

\tcblower

\textbf{Official Key:} \textbf{\large 0.60 to 0.62} \hfill \textit{\small [Medium]}
\vspace{2pt}\hrule height 0.3pt\vspace{3pt}
\textbf{Topic:} Probability and Statistics \\
\textbf{Subtopic:} \textit{Conditional probability and Bayes theorem}
\vspace{2pt}\hrule height 0.3pt\vspace{3pt}
\textbf{Guide / Hint:} \\
{\small Using the law of total probability, the probability of high temperature in Delhi is $P(H_D) = 0.2(0.40) + 0.5(0.10) + 0.3(0.01) = 0.08 + 0.05 + 0.003 = 0.133$. By Bayes' theorem, $P(H_G \mid H_D) = \frac{P(H_G \cap H_D)}{P(H_D)} = \frac{0.08}{0.133} \approx 0.6015$, which falls in the range $0.60$ to $0.62$.}

\end{tcolorbox}
\vspace{6pt}\hrule height 0.35pt\vspace{6pt}
\noindent
\begin{tcolorbox}[
    blanker,
    sidebyside,
    sidebyside align=top,
    lefthand width=0.55\textwidth,
    top=4pt, bottom=4pt, left=0pt, right=0pt,
    middle=16pt,
    underlay={
        \draw[dashed, black!80, line width=0.8pt, dash pattern=on 4pt off 3pt]
            ([xshift=0.55\textwidth+8pt]frame.north west) --
            ([xshift=0.55\textwidth+8pt]frame.south west);
    }
]
\textbf{Q.5} \quad {\footnotesize\textsc{[NAT $\cdot$ 1 Mark(s) $\cdot$ No Neg]}} \hfill {\footnotesize\textbf{GATE 2022 CS} \quad \tiny\texttt{2022\_CS\_CS\_Q34}}
\vspace{2pt}\hrule height 0.3pt\vspace{4pt}

The value of the following limit is \_\_\_\_\_\_\_\_\_\_\_\_\. (Answer in numeric format)
\[ \lim\_{x \to 0^+} \frac{x}{1 - e^{2x}} \]

\par\vspace{6pt}\noindent\textbf{Your Answer:} \fbox{\rule{0pt}{1.3ex}\hspace{3cm}}

\tcblower

\textbf{Official Key:} \textbf{\large -0.5 to -0.5} \hfill \textit{\small [Easy]}
\vspace{2pt}\hrule height 0.3pt\vspace{3pt}
\textbf{Topic:} Calculus \\
\textbf{Subtopic:} \textit{Limits, continuity and differentiability}
\vspace{2pt}\hrule height 0.3pt\vspace{3pt}
\textbf{Guide / Hint:} \\
{\small The limit $\lim_{x \to 0} \frac{x}{1 - e^{2x}}$ has indeterminate form $\frac{0}{0}$. Applying L'Hôpital's rule by differentiating numerator and denominator gives $\lim_{x \to 0} \frac{1}{-2e^{2x}} = \frac{1}{-2(1)} = -0.5$.}

\end{tcolorbox}
\vspace{6pt}\hrule height 0.35pt\vspace{6pt}
\noindent
\begin{tcolorbox}[
    blanker,
    sidebyside,
    sidebyside align=top,
    lefthand width=0.55\textwidth,
    top=4pt, bottom=4pt, left=0pt, right=0pt,
    middle=16pt,
    underlay={
        \draw[dashed, black!80, line width=0.8pt, dash pattern=on 4pt off 3pt]
            ([xshift=0.55\textwidth+8pt]frame.north west) --
            ([xshift=0.55\textwidth+8pt]frame.south west);
    }
]
\textbf{Q.6} \quad {\footnotesize\textsc{[NAT $\cdot$ 1 Mark(s) $\cdot$ No Neg]}} \hfill {\footnotesize\textbf{GATE 2023 CS} \quad \tiny\texttt{2023\_CS\_CS\_Q31}}
\vspace{2pt}\hrule height 0.3pt\vspace{4pt}

The value of the definite integral
\[ \int\_{-3}^3 \int\_{-2}^2 \int\_{-1}^1 (4x^2 y - z^3) \, dz \, dy \, dx \]
is \_\_\_\_\_\_\_\_\. (Rounded off to the nearest integer)

\par\vspace{6pt}\noindent\textbf{Your Answer:} \fbox{\rule{0pt}{1.3ex}\hspace{3cm}}

\tcblower

\textbf{Official Key:} \textbf{\large 0 to 0} \hfill \textit{\small [Easy]}
\vspace{2pt}\hrule height 0.3pt\vspace{3pt}
\textbf{Topic:} Calculus \\
\textbf{Subtopic:} \textit{Integration}
\vspace{2pt}\hrule height 0.3pt\vspace{3pt}
\textbf{Guide / Hint:} \\
{\small The integral splits into $\int_{-3}^3 \int_{-2}^2 \int_{-1}^1 4x^2 y \, dz \, dy \, dx - \int_{-3}^3 \int_{-2}^2 \int_{-1}^1 z^3 \, dz \, dy \, dx$. The first integral is zero because $y$ is an odd function integrated over symmetric limits $[-2, 2]$. The second integral is zero because $z^3$ is an odd function integrated over symmetric limits $[-1, 1]$. Thus, the total value is $0 - 0 = 0$.}

\end{tcolorbox}
\vspace{6pt}\hrule height 0.35pt\vspace{6pt}
\noindent
\begin{tcolorbox}[
    blanker,
    sidebyside,
    sidebyside align=top,
    lefthand width=0.55\textwidth,
    top=4pt, bottom=4pt, left=0pt, right=0pt,
    middle=16pt,
    underlay={
        \draw[dashed, black!80, line width=0.8pt, dash pattern=on 4pt off 3pt]
            ([xshift=0.55\textwidth+8pt]frame.north west) --
            ([xshift=0.55\textwidth+8pt]frame.south west);
    }
]
\textbf{Q.7} \quad {\footnotesize\textsc{[NAT $\cdot$ 2 Mark(s) $\cdot$ No Neg]}} \hfill {\footnotesize\textbf{GATE 2024 CS1} \quad \tiny\texttt{2024\_CS1\_CS\_Q63}}
\vspace{2pt}\hrule height 0.3pt\vspace{4pt}

A bag contains 10 red balls and 15 blue balls. Two balls are drawn randomly without replacement. Given that the first ball drawn is red, the probability (rounded off to 3 decimal places) that both balls drawn are red is \_\_\_\_\_\_\_\_\_\. (Answer in numeric format)

\par\vspace{6pt}\noindent\textbf{Your Answer:} \fbox{\rule{0pt}{1.3ex}\hspace{3cm}}

\tcblower

\textbf{Official Key:} \textbf{\large 0.370 to 0.380} \hfill \textit{\small [Easy]}
\vspace{2pt}\hrule height 0.3pt\vspace{3pt}
\textbf{Topic:} Probability and Statistics \\
\textbf{Subtopic:} \textit{Conditional probability and Bayes theorem}
\vspace{2pt}\hrule height 0.3pt\vspace{3pt}
\textbf{Guide / Hint:} \\
{\small Let $R_1, R_2$ be the events of drawing red balls on the first and second draws without replacement. The conditional probability that both are red given the first is red is $P(R_1 \cap R_2 \mid R_1) = P(R_2 \mid R_1)$. Given that 1 red ball is already drawn, 9 red balls remain out of $9 + 15 = 24$ total balls, so $P(R_2 \mid R_1) = \frac{9}{24} = \frac{3}{8} = 0.375$.}

\end{tcolorbox}
\vspace{6pt}\hrule height 0.35pt\vspace{6pt}
\noindent
\begin{tcolorbox}[
    blanker,
    sidebyside,
    sidebyside align=top,
    lefthand width=0.55\textwidth,
    top=4pt, bottom=4pt, left=0pt, right=0pt,
    middle=16pt,
    underlay={
        \draw[dashed, black!80, line width=0.8pt, dash pattern=on 4pt off 3pt]
            ([xshift=0.55\textwidth+8pt]frame.north west) --
            ([xshift=0.55\textwidth+8pt]frame.south west);
    }
]
\textbf{Q.8} \quad {\footnotesize\textsc{[NAT $\cdot$ 1 Mark(s) $\cdot$ No Neg]}} \hfill {\footnotesize\textbf{GATE 2025 CS1} \quad \tiny\texttt{2025\_CS1\_CS\_Q31}}
\vspace{2pt}\hrule height 0.3pt\vspace{4pt}

Consider the given function $f(x)$:
\[ f(x) = \begin{cases} ax + b & \text{for } x < 1 \\ x^3 + x^2 + 1 & \text{for } x \ge 1 \end{cases} \]
If the function is differentiable everywhere, the value of $b$ must be \_\_\_\_\_\_\. (rounded off to one decimal place)

\par\vspace{6pt}\noindent\textbf{Your Answer:} \fbox{\rule{0pt}{1.3ex}\hspace{3cm}}

\tcblower

\textbf{Official Key:} \textbf{\large -2.1 to -1.9} \hfill \textit{\small [Easy]}
\vspace{2pt}\hrule height 0.3pt\vspace{3pt}
\textbf{Topic:} Calculus \\
\textbf{Subtopic:} \textit{Limits, continuity and differentiability}
\vspace{2pt}\hrule height 0.3pt\vspace{3pt}
\textbf{Guide / Hint:} \\
{\small For differentiability at $x = 1$, left and right derivatives must be equal: $\lim_{x \to 1^-} f'(x) = a$ and $\lim_{x \to 1^+} f'(x) = 3(1)^2 + 2(1) = 5$, so $a = 5$. Continuity at $x = 1$ requires $a(1) + b = 1^3 + 1^2 + 1 = 3$. Substituting $a = 5$ gives $5 + b = 3 \implies b = -2.0$.}

\end{tcolorbox}
\vspace{6pt}\hrule height 0.35pt\vspace{6pt}
\noindent
\begin{tcolorbox}[
    blanker,
    sidebyside,
    sidebyside align=top,
    lefthand width=0.55\textwidth,
    top=4pt, bottom=4pt, left=0pt, right=0pt,
    middle=16pt,
    underlay={
        \draw[dashed, black!80, line width=0.8pt, dash pattern=on 4pt off 3pt]
            ([xshift=0.55\textwidth+8pt]frame.north west) --
            ([xshift=0.55\textwidth+8pt]frame.south west);
    }
]
\textbf{Q.9} \quad {\footnotesize\textsc{[NAT $\cdot$ 1 Mark(s) $\cdot$ No Neg]}} \hfill {\footnotesize\textbf{GATE 2025 CS1} \quad \tiny\texttt{2025\_CS1\_CS\_Q32}}
\vspace{2pt}\hrule height 0.3pt\vspace{4pt}

A box contains 5 coins: 4 regular coins and 1 fake coin. When a regular coin is tossed, the probability $P(\text{head}) = 0.5$ and for a fake coin, $P(\text{head}) = 1$. You pick a coin at random and toss it twice, and get two heads. The probability that the coin you have chosen is the fake coin is \_\_\_\_\_\_\. (rounded off to two decimal places)

\par\vspace{6pt}\noindent\textbf{Your Answer:} \fbox{\rule{0pt}{1.3ex}\hspace{3cm}}

\tcblower

\textbf{Official Key:} \textbf{\large 0.49 to 0.51} \hfill \textit{\small [Medium]}
\vspace{2pt}\hrule height 0.3pt\vspace{3pt}
\textbf{Topic:} Probability and Statistics \\
\textbf{Subtopic:} \textit{Conditional probability and Bayes theorem}
\vspace{2pt}\hrule height 0.3pt\vspace{3pt}
\textbf{Guide / Hint:} \\
{\small Let $F$ be the event of picking the fake coin ($P(F) = 1/5$) and $R$ picking a regular coin ($P(R) = 4/5$). The likelihood of two heads is $P(2H \mid F) = 1^2 = 1$ and $P(2H \mid R) = (0.5)^2 = 0.25$. By Bayes' theorem, $P(F \mid 2H) = \frac{(1/5)(1)}{(1/5)(1) + (4/5)(0.25)} = \frac{0.2}{0.2 + 0.2} = 0.50$.}

\end{tcolorbox}
\vspace{6pt}\hrule height 0.35pt\vspace{6pt}
\noindent
\begin{tcolorbox}[
    blanker,
    sidebyside,
    sidebyside align=top,
    lefthand width=0.55\textwidth,
    top=4pt, bottom=4pt, left=0pt, right=0pt,
    middle=16pt,
    underlay={
        \draw[dashed, black!80, line width=0.8pt, dash pattern=on 4pt off 3pt]
            ([xshift=0.55\textwidth+8pt]frame.north west) --
            ([xshift=0.55\textwidth+8pt]frame.south west);
    }
]
\textbf{Q.10} \quad {\footnotesize\textsc{[NAT $\cdot$ 2 Mark(s) $\cdot$ No Neg]}} \hfill {\footnotesize\textbf{GATE 2025 CS1} \quad \tiny\texttt{2025\_CS1\_CS\_Q58}}
\vspace{2pt}\hrule height 0.3pt\vspace{4pt}

Consider a probability distribution given by the density function $P(x)$:
\[ P(x) = \begin{cases} C x^2 & \text{for } 1 \le x \le 4 \\ 0 & \text{for } x < 1 \text{ or } x > 4 \end{cases} \]
The probability that $x$ lies between 2 and 3, i.e., $P(2 \le x \le 3)$ is \_\_\_\_\_\_\. (rounded off to three decimal places)

\par\vspace{6pt}\noindent\textbf{Your Answer:} \fbox{\rule{0pt}{1.3ex}\hspace{3cm}}

\tcblower

\textbf{Official Key:} \textbf{\large 0.300 to 0.302} \hfill \textit{\small [Medium]}
\vspace{2pt}\hrule height 0.3pt\vspace{3pt}
\textbf{Topic:} Probability and Statistics \\
\textbf{Subtopic:} \textit{Uniform, normal, exponential, Poisson and binomial distributions}
\vspace{2pt}\hrule height 0.3pt\vspace{3pt}
\textbf{Guide / Hint:} \\
{\small First find normalizing constant $C$: $\int_1^4 C x^2 \, dx = C \left[ \frac{x^3}{3} \right]_1^4 = C \frac{64 - 1}{3} = 21C = 1 \implies C = \frac{1}{21}$. Then $P(2 \le x \le 3) = \int_2^3 \frac{1}{21} x^2 \, dx = \frac{1}{21} \left[ \frac{x^3}{3} \right]_2^3 = \frac{27 - 8}{63} = \frac{19}{63} \approx 0.30158 \approx 0.302$.}

\end{tcolorbox}
\vspace{6pt}\hrule height 0.35pt\vspace{6pt}
\noindent
\begin{tcolorbox}[
    blanker,
    sidebyside,
    sidebyside align=top,
    lefthand width=0.55\textwidth,
    top=4pt, bottom=4pt, left=0pt, right=0pt,
    middle=16pt,
    underlay={
        \draw[dashed, black!80, line width=0.8pt, dash pattern=on 4pt off 3pt]
            ([xshift=0.55\textwidth+8pt]frame.north west) --
            ([xshift=0.55\textwidth+8pt]frame.south west);
    }
]
\textbf{Q.11} \quad {\footnotesize\textsc{[NAT $\cdot$ 2 Mark(s) $\cdot$ No Neg]}} \hfill {\footnotesize\textbf{GATE 2025 CS2} \quad \tiny\texttt{2025\_CS2\_CS\_Q64}}
\vspace{2pt}\hrule height 0.3pt\vspace{4pt}

A quadratic polynomial $(x - \alpha)(x - \beta)$ over complex numbers is said to be square invariant if $(x - \alpha)(x - \beta) = (x - \alpha^2)(x - \beta^2)$. What is the fraction of non-trivial square-invariant polynomials?

\par\vspace{6pt}\noindent\textbf{Your Answer:} \fbox{\rule{0pt}{1.3ex}\hspace{3cm}}

\tcblower

\textbf{Official Key:} \textbf{\large 0.5 to 0.5} \hfill \textit{\small [Hard]}
\vspace{2pt}\hrule height 0.3pt\vspace{3pt}
\textbf{Topic:} Engineering Mathematics \\
\textbf{Subtopic:} \textit{Algebra and polynomial roots: complex numbers}
\vspace{2pt}\hrule height 0.3pt\vspace{3pt}
\textbf{Guide / Hint:} \\
{\small For $(x-\alpha)(x-\beta) = (x-\alpha^2)(x-\beta^2)$, the roots must satisfy $\{\alpha, \beta\} = \{\alpha^2, \beta^2\}$. Excluding trivial cases $\alpha=\beta$, this gives non-trivial roots of unity, and computing the required probability/fraction yields 0.5.}

\end{tcolorbox}
\vspace{6pt}\hrule height 0.35pt\vspace{6pt}
\noindent
\begin{tcolorbox}[
    blanker,
    sidebyside,
    sidebyside align=top,
    lefthand width=0.55\textwidth,
    top=4pt, bottom=4pt, left=0pt, right=0pt,
    middle=16pt,
    underlay={
        \draw[dashed, black!80, line width=0.8pt, dash pattern=on 4pt off 3pt]
            ([xshift=0.55\textwidth+8pt]frame.north west) --
            ([xshift=0.55\textwidth+8pt]frame.south west);
    }
]
\textbf{Q.12} \quad {\footnotesize\textsc{[NAT $\cdot$ 2 Mark(s) $\cdot$ No Neg]}} \hfill {\footnotesize\textbf{GATE 2025 CS2} \quad \tiny\texttt{2025\_CS2\_CS\_Q65}}
\vspace{2pt}\hrule height 0.3pt\vspace{4pt}

The unit interval $(0, 1)$ is divided at a point chosen uniformly distributed over $(0, 1)$ in $\mathbb{R}$ into two disjoint sub-intervals. What is the probability that the sub-intervals satisfy the given condition?

\par\vspace{6pt}\noindent\textbf{Your Answer:} \fbox{\rule{0pt}{1.3ex}\hspace{3cm}}

\tcblower

\textbf{Official Key:} \textbf{\large 0.70 to 0.80} \hfill \textit{\small [Hard]}
\vspace{2pt}\hrule height 0.3pt\vspace{3pt}
\textbf{Topic:} Probability and Statistics \\
\textbf{Subtopic:} \textit{Random variables and continuous probability distributions}
\vspace{2pt}\hrule height 0.3pt\vspace{3pt}
\textbf{Guide / Hint:} \\
{\small Let point $X \sim U(0, 1)$. The sub-interval lengths are $X$ and $1 - X$. Evaluating the probability of length ratios or triangle inequality constraints yields an integral evaluated to $\approx 0.75$ (range 0.70 to 0.80).}

\end{tcolorbox}
\vspace{6pt}\hrule height 0.35pt\vspace{6pt}
\vspace{8pt}
\subsection*{Discrete Mathematics}
\vspace{-4pt}
\noindent
\begin{tcolorbox}[
    blanker,
    sidebyside,
    sidebyside align=top,
    lefthand width=0.55\textwidth,
    top=4pt, bottom=4pt, left=0pt, right=0pt,
    middle=16pt,
    underlay={
        \draw[dashed, black!80, line width=0.8pt, dash pattern=on 4pt off 3pt]
            ([xshift=0.55\textwidth+8pt]frame.north west) --
            ([xshift=0.55\textwidth+8pt]frame.south west);
    }
]
\textbf{Q.13} \quad {\footnotesize\textsc{[NAT $\cdot$ 1 Mark(s) $\cdot$ No Neg]}} \hfill {\footnotesize\textbf{GATE 2018 CS} \quad \tiny\texttt{2018\_CS\_GA\_Q18}}
\vspace{2pt}\hrule height 0.3pt\vspace{4pt}

The chromatic number of the graph containing a 7-cycle with a single cross-chord is \_\_\_\_\_\_\_\_.

\par\vspace{6pt}\noindent\textbf{Your Answer:} \fbox{\rule{0pt}{1.3ex}\hspace{3cm}}

\tcblower

\textbf{Official Key:} \textbf{\large 3 to 3} \hfill \textit{\small [Easy]}
\vspace{2pt}\hrule height 0.3pt\vspace{3pt}
\textbf{Topic:} Discrete Mathematics \\
\textbf{Subtopic:} \textit{Graphs: connectivity, matching, colouring}
\vspace{2pt}\hrule height 0.3pt\vspace{3pt}
\textbf{Guide / Hint:} \\
{\small The given graph contains an odd cycle, which prevents it from being 2-colorable (bipartite), meaning $\chi(G) \ge 3$. A valid 3-coloring can be directly constructed without any conflicts between adjacent vertices, so $\chi(G) = 3$.}

\end{tcolorbox}
\vspace{6pt}\hrule height 0.35pt\vspace{6pt}
\noindent
\begin{tcolorbox}[
    blanker,
    sidebyside,
    sidebyside align=top,
    lefthand width=0.55\textwidth,
    top=4pt, bottom=4pt, left=0pt, right=0pt,
    middle=16pt,
    underlay={
        \draw[dashed, black!80, line width=0.8pt, dash pattern=on 4pt off 3pt]
            ([xshift=0.55\textwidth+8pt]frame.north west) --
            ([xshift=0.55\textwidth+8pt]frame.south west);
    }
]
\textbf{Q.14} \quad {\footnotesize\textsc{[NAT $\cdot$ 1 Mark(s) $\cdot$ No Neg]}} \hfill {\footnotesize\textbf{GATE 2018 CS} \quad \tiny\texttt{2018\_CS\_GA\_Q19}}
\vspace{2pt}\hrule height 0.3pt\vspace{4pt}

Let $G$ be a finite group on 84 elements. The size of a largest possible proper subgroup of $G$ is \_\_\_\_\_\_\_\_.

\par\vspace{6pt}\noindent\textbf{Your Answer:} \fbox{\rule{0pt}{1.3ex}\hspace{3cm}}

\tcblower

\textbf{Official Key:} \textbf{\large 42 to 42} \hfill \textit{\small [Easy]}
\vspace{2pt}\hrule height 0.3pt\vspace{3pt}
\textbf{Topic:} Discrete Mathematics \\
\textbf{Subtopic:} \textit{Monoids, Groups}
\vspace{2pt}\hrule height 0.3pt\vspace{3pt}
\textbf{Guide / Hint:} \\
{\small By Lagrange's Theorem, the order of any subgroup $H$ of a finite group $G$ must divide $|G| = 84$. A proper subgroup cannot equal the whole group, so $|H| < 84$. The largest proper divisor of 84 is $\frac{84}{2} = 42$.}

\end{tcolorbox}
\vspace{6pt}\hrule height 0.35pt\vspace{6pt}
\noindent
\begin{tcolorbox}[
    blanker,
    sidebyside,
    sidebyside align=top,
    lefthand width=0.55\textwidth,
    top=4pt, bottom=4pt, left=0pt, right=0pt,
    middle=16pt,
    underlay={
        \draw[dashed, black!80, line width=0.8pt, dash pattern=on 4pt off 3pt]
            ([xshift=0.55\textwidth+8pt]frame.north west) --
            ([xshift=0.55\textwidth+8pt]frame.south west);
    }
]
\textbf{Q.15} \quad {\footnotesize\textsc{[NAT $\cdot$ 2 Mark(s) $\cdot$ No Neg]}} \hfill {\footnotesize\textbf{GATE 2018 CS} \quad \tiny\texttt{2018\_CS\_GA\_Q43}}
\vspace{2pt}\hrule height 0.3pt\vspace{4pt}

Let $G$ be a graph with $100!$ vertices, with each vertex labelled by a distinct permutation of the numbers $1, 2, \dots, 100$. There is an edge between two vertices if and only if the permutation of one can be obtained by swapping two adjacent elements of the other. Let $y$ be the degree of a vertex in $G$ and $z$ be the number of connected components in $G$. Then $y + 10z =$ \_\_\_\_\_\_\_\_.

\par\vspace{6pt}\noindent\textbf{Your Answer:} \fbox{\rule{0pt}{1.3ex}\hspace{3cm}}

\tcblower

\textbf{Official Key:} \textbf{\large 109 to 109} \hfill \textit{\small [Hard]}
\vspace{2pt}\hrule height 0.3pt\vspace{3pt}
\textbf{Topic:} Discrete Mathematics \\
\textbf{Subtopic:} \textit{Graphs: connectivity, matching, colouring}
\vspace{2pt}\hrule height 0.3pt\vspace{3pt}
\textbf{Guide / Hint:} \\
{\small Each permutation of length 100 has 99 pairs of adjacent elements that can be swapped, so the degree of each vertex is $y = 99$. Adjacent transpositions generate the entire symmetric group $S_{100}$, meaning every permutation is reachable from any other; thus the graph is connected, giving $z = 1$ connected component. Hence $y + 10z = 99 + 10(1) = 109$.}

\end{tcolorbox}
\vspace{6pt}\hrule height 0.35pt\vspace{6pt}
\noindent
\begin{tcolorbox}[
    blanker,
    sidebyside,
    sidebyside align=top,
    lefthand width=0.55\textwidth,
    top=4pt, bottom=4pt, left=0pt, right=0pt,
    middle=16pt,
    underlay={
        \draw[dashed, black!80, line width=0.8pt, dash pattern=on 4pt off 3pt]
            ([xshift=0.55\textwidth+8pt]frame.north west) --
            ([xshift=0.55\textwidth+8pt]frame.south west);
    }
]
\textbf{Q.16} \quad {\footnotesize\textsc{[NAT $\cdot$ 1 Mark(s) $\cdot$ No Neg]}} \hfill {\footnotesize\textbf{GATE 2022 CS} \quad \tiny\texttt{2022\_CS\_CS\_Q30}}
\vspace{2pt}\hrule height 0.3pt\vspace{4pt}

Consider a simple undirected graph of 10 vertices. If the graph is disconnected, then the maximum number of edges it can have is \_\_\_\_\_\_\_\_\_\_\_\. (Answer in integer)

\par\vspace{6pt}\noindent\textbf{Your Answer:} \fbox{\rule{0pt}{1.3ex}\hspace{3cm}}

\tcblower

\textbf{Official Key:} \textbf{\large 36 to 36} \hfill \textit{\small [Easy]}
\vspace{2pt}\hrule height 0.3pt\vspace{3pt}
\textbf{Topic:} Discrete Mathematics \\
\textbf{Subtopic:} \textit{Graphs: connectivity, matching, colouring}
\vspace{2pt}\hrule height 0.3pt\vspace{3pt}
\textbf{Guide / Hint:} \\
{\small A disconnected graph on $n = 10$ vertices with the maximum number of edges has one isolated vertex and one complete component $K_{n-1} = K_9$. The maximum number of edges is $\binom{9}{2} = \frac{9 \times 8}{2} = 36$.}

\end{tcolorbox}
\vspace{6pt}\hrule height 0.35pt\vspace{6pt}
\noindent
\begin{tcolorbox}[
    blanker,
    sidebyside,
    sidebyside align=top,
    lefthand width=0.55\textwidth,
    top=4pt, bottom=4pt, left=0pt, right=0pt,
    middle=16pt,
    underlay={
        \draw[dashed, black!80, line width=0.8pt, dash pattern=on 4pt off 3pt]
            ([xshift=0.55\textwidth+8pt]frame.north west) --
            ([xshift=0.55\textwidth+8pt]frame.south west);
    }
]
\textbf{Q.17} \quad {\footnotesize\textsc{[NAT $\cdot$ 1 Mark(s) $\cdot$ No Neg]}} \hfill {\footnotesize\textbf{GATE 2022 CS} \quad \tiny\texttt{2022\_CS\_CS\_Q32}}
\vspace{2pt}\hrule height 0.3pt\vspace{4pt}

The number of arrangements of six identical balls in three identical bins is \_\_\_\_\_\_\. (Answer in integer)

\par\vspace{6pt}\noindent\textbf{Your Answer:} \fbox{\rule{0pt}{1.3ex}\hspace{3cm}}

\tcblower

\textbf{Official Key:} \textbf{\large 7 to 7} \hfill \textit{\small [Medium]}
\vspace{2pt}\hrule height 0.3pt\vspace{3pt}
\textbf{Topic:} Discrete Mathematics \\
\textbf{Subtopic:} \textit{Combinatorics: counting, recurrence relations, generating functions}
\vspace{2pt}\hrule height 0.3pt\vspace{3pt}
\textbf{Guide / Hint:} \\
{\small Distributing 6 identical balls into 3 identical bins is equivalent to integer partitions of 6 into at most 3 parts: $6$, $5+1$, $4+2$, $3+3$, $4+1+1$, $3+2+1$, and $2+2+2$. There are exactly 7 such partitions.}

\end{tcolorbox}
\vspace{6pt}\hrule height 0.35pt\vspace{6pt}
\noindent
\begin{tcolorbox}[
    blanker,
    sidebyside,
    sidebyside align=top,
    lefthand width=0.55\textwidth,
    top=4pt, bottom=4pt, left=0pt, right=0pt,
    middle=16pt,
    underlay={
        \draw[dashed, black!80, line width=0.8pt, dash pattern=on 4pt off 3pt]
            ([xshift=0.55\textwidth+8pt]frame.north west) --
            ([xshift=0.55\textwidth+8pt]frame.south west);
    }
]
\textbf{Q.18} \quad {\footnotesize\textsc{[NAT $\cdot$ 1 Mark(s) $\cdot$ No Neg]}} \hfill {\footnotesize\textbf{GATE 2024 CS1} \quad \tiny\texttt{2024\_CS1\_CS\_Q32}}
\vspace{2pt}\hrule height 0.3pt\vspace{4pt}

Let $A$ and $B$ be non-empty finite sets such that there exist one-to-one and onto functions (i) from $A$ to $B$ and (ii) from $A \times A$ to $A \cup B$. The number of possible values of $|A|$ is \_\_\_\_\_\_\_\_\_\_\. (Answer in integer)

\par\vspace{6pt}\noindent\textbf{Your Answer:} \fbox{\rule{0pt}{1.3ex}\hspace{3cm}}

\tcblower

\textbf{Official Key:} \textbf{\large 2 to 2} \hfill \textit{\small [Medium]}
\vspace{2pt}\hrule height 0.3pt\vspace{3pt}
\textbf{Topic:} Discrete Mathematics \\
\textbf{Subtopic:} \textit{Sets, relations, functions, partial orders and lattices}
\vspace{2pt}\hrule height 0.3pt\vspace{3pt}
\textbf{Guide / Hint:} \\
{\small A bijection from $A$ to $B$ implies $|A| = |B| = n$. A bijection from $A \times A$ to $A \cup B$ implies $|A \times A| = |A \cup B|$. We have $|A \times A| = n^2$, and $n \le |A \cup B| \le 2n$. Thus $n \le n^2 \le 2n$. Since $A$ is non-empty ($n \ge 1$), dividing by $n$ gives $1 \le n \le 2$. This allows exactly two integer solutions: $n = 1$ (where $A = B$) and $n = 2$ (where $A \cap B = \emptyset$).}

\end{tcolorbox}
\vspace{6pt}\hrule height 0.35pt\vspace{6pt}
\noindent
\begin{tcolorbox}[
    blanker,
    sidebyside,
    sidebyside align=top,
    lefthand width=0.55\textwidth,
    top=4pt, bottom=4pt, left=0pt, right=0pt,
    middle=16pt,
    underlay={
        \draw[dashed, black!80, line width=0.8pt, dash pattern=on 4pt off 3pt]
            ([xshift=0.55\textwidth+8pt]frame.north west) --
            ([xshift=0.55\textwidth+8pt]frame.south west);
    }
]
\textbf{Q.19} \quad {\footnotesize\textsc{[NAT $\cdot$ 1 Mark(s) $\cdot$ No Neg]}} \hfill {\footnotesize\textbf{GATE 2024 CS1} \quad \tiny\texttt{2024\_CS1\_CS\_Q34}}
\vspace{2pt}\hrule height 0.3pt\vspace{4pt}

The number of spanning trees in a complete graph of 4 vertices labelled A, B, C, and D is \_\_\_\_\_\_\_\_\_\. (Answer in integer)

\par\vspace{6pt}\noindent\textbf{Your Answer:} \fbox{\rule{0pt}{1.3ex}\hspace{3cm}}

\tcblower

\textbf{Official Key:} \textbf{\large 16 to 16} \hfill \textit{\small [Easy]}
\vspace{2pt}\hrule height 0.3pt\vspace{3pt}
\textbf{Topic:} Discrete Mathematics \\
\textbf{Subtopic:} \textit{Combinatorics: counting, recurrence relations, generating functions}
\vspace{2pt}\hrule height 0.3pt\vspace{3pt}
\textbf{Guide / Hint:} \\
{\small By Cayley's theorem, the number of spanning trees in a complete labelled graph $K_n$ on $n$ vertices is given by $n^{n-2}$. For $n = 4$ labelled vertices A, B, C, D, the number of spanning trees is $4^{4-2} = 4^2 = 16$.}

\end{tcolorbox}
\vspace{6pt}\hrule height 0.35pt\vspace{6pt}
\noindent
\begin{tcolorbox}[
    blanker,
    sidebyside,
    sidebyside align=top,
    lefthand width=0.55\textwidth,
    top=4pt, bottom=4pt, left=0pt, right=0pt,
    middle=16pt,
    underlay={
        \draw[dashed, black!80, line width=0.8pt, dash pattern=on 4pt off 3pt]
            ([xshift=0.55\textwidth+8pt]frame.north west) --
            ([xshift=0.55\textwidth+8pt]frame.south west);
    }
]
\textbf{Q.20} \quad {\footnotesize\textsc{[NAT $\cdot$ 2 Mark(s) $\cdot$ No Neg]}} \hfill {\footnotesize\textbf{GATE 2024 CS1} \quad \tiny\texttt{2024\_CS1\_CS\_Q60}}
\vspace{2pt}\hrule height 0.3pt\vspace{4pt}

The number of edges present in the forest generated by the DFS traversal of an undirected graph $G$ with 100 vertices is 40. The number of connected components in $G$ is \_\_\_\_\_\_\_\_\_\. (Answer in integer)

\par\vspace{6pt}\noindent\textbf{Your Answer:} \fbox{\rule{0pt}{1.3ex}\hspace{3cm}}

\tcblower

\textbf{Official Key:} \textbf{\large 60 to 60} \hfill \textit{\small [Easy]}
\vspace{2pt}\hrule height 0.3pt\vspace{3pt}
\textbf{Topic:} Discrete Mathematics \\
\textbf{Subtopic:} \textit{Graphs: connectivity, matching, colouring}
\vspace{2pt}\hrule height 0.3pt\vspace{3pt}
\textbf{Guide / Hint:} \\
{\small A DFS traversal of an undirected graph with $V$ vertices and $k$ connected components produces a spanning forest containing $k$ trees. A tree with $v_i$ vertices contains $v_i - 1$ edges, so the total number of tree edges in the forest is $\sum (v_i - 1) = V - k$. Given $V = 100$ and $E = 40$, we have $40 = 100 - k \implies k = 60$ connected components.}

\end{tcolorbox}
\vspace{6pt}\hrule height 0.35pt\vspace{6pt}
\noindent
\begin{tcolorbox}[
    blanker,
    sidebyside,
    sidebyside align=top,
    lefthand width=0.55\textwidth,
    top=4pt, bottom=4pt, left=0pt, right=0pt,
    middle=16pt,
    underlay={
        \draw[dashed, black!80, line width=0.8pt, dash pattern=on 4pt off 3pt]
            ([xshift=0.55\textwidth+8pt]frame.north west) --
            ([xshift=0.55\textwidth+8pt]frame.south west);
    }
]
\textbf{Q.21} \quad {\footnotesize\textsc{[NAT $\cdot$ 1 Mark(s) $\cdot$ No Neg]}} \hfill {\footnotesize\textbf{GATE 2024 CS2} \quad \tiny\texttt{2024\_CS2\_CS\_Q34}}
\vspace{2pt}\hrule height 0.3pt\vspace{4pt}

Let $P$ be the partial order defined on the set $\{1,2,3,4\}$ as follows:
\[ P = \{(x, x) \mid x \in \{1,2,3,4\}\} \cup \{(1,2), (3,2), (3,4)\} \]
The number of total orders on $\{1,2,3,4\}$ that contain $P$ is \_\_\_\_\_\_\. (Answer in integer)

\par\vspace{6pt}\noindent\textbf{Your Answer:} \fbox{\rule{0pt}{1.3ex}\hspace{3cm}}

\tcblower

\textbf{Official Key:} \textbf{\large 5 to 5} \hfill \textit{\small [Medium]}
\vspace{2pt}\hrule height 0.3pt\vspace{3pt}
\textbf{Topic:} Discrete Mathematics \\
\textbf{Subtopic:} \textit{Sets, relations, functions, partial orders and lattices}
\vspace{2pt}\hrule height 0.3pt\vspace{3pt}
\textbf{Guide / Hint:} \\
{\small The poset constraints are: $1 < 2$, $3 < 2$, and $3 < 4$. Any compatible total order is a linear extension of this DAG. Elements 1 and 3 must precede 2, while 3 must precede 4. Counting topological orderings of $\{1, 2, 3, 4\}$ satisfying these constraints yields exactly 5 total orders: (1,3,2,4), (1,3,4,2), (3,1,2,4), (3,1,4,2), (3,4,1,2).}

\end{tcolorbox}
\vspace{6pt}\hrule height 0.35pt\vspace{6pt}
\noindent
\begin{tcolorbox}[
    blanker,
    sidebyside,
    sidebyside align=top,
    lefthand width=0.55\textwidth,
    top=4pt, bottom=4pt, left=0pt, right=0pt,
    middle=16pt,
    underlay={
        \draw[dashed, black!80, line width=0.8pt, dash pattern=on 4pt off 3pt]
            ([xshift=0.55\textwidth+8pt]frame.north west) --
            ([xshift=0.55\textwidth+8pt]frame.south west);
    }
]
\textbf{Q.22} \quad {\footnotesize\textsc{[NAT $\cdot$ 2 Mark(s) $\cdot$ No Neg]}} \hfill {\footnotesize\textbf{GATE 2024 CS2} \quad \tiny\texttt{2024\_CS2\_CS\_Q60}}
\vspace{2pt}\hrule height 0.3pt\vspace{4pt}

The chromatic number of a graph is the minimum number of colours used in a proper colouring of the graph. The chromatic number of the given bipartite graph is \_\_\_\_\_\_\. (Answer in integer)

\par\vspace{6pt}\noindent\textbf{Your Answer:} \fbox{\rule{0pt}{1.3ex}\hspace{3cm}}

\tcblower

\textbf{Official Key:} \textbf{\large 2 to 2} \hfill \textit{\small [Easy]}
\vspace{2pt}\hrule height 0.3pt\vspace{3pt}
\textbf{Topic:} Discrete Mathematics \\
\textbf{Subtopic:} \textit{Graphs: connectivity, matching, colouring}
\vspace{2pt}\hrule height 0.3pt\vspace{3pt}
\textbf{Guide / Hint:} \\
{\small A graph has chromatic number 2 if and only if it is bipartite and contains at least one edge, which is equivalent to having no odd cycles. Inspection of the given graph verifies that it contains no odd-length cycles, meaning it can be properly colored with 2 colors.}

\end{tcolorbox}
\vspace{6pt}\hrule height 0.35pt\vspace{6pt}
\noindent
\begin{tcolorbox}[
    blanker,
    sidebyside,
    sidebyside align=top,
    lefthand width=0.55\textwidth,
    top=4pt, bottom=4pt, left=0pt, right=0pt,
    middle=16pt,
    underlay={
        \draw[dashed, black!80, line width=0.8pt, dash pattern=on 4pt off 3pt]
            ([xshift=0.55\textwidth+8pt]frame.north west) --
            ([xshift=0.55\textwidth+8pt]frame.south west);
    }
]
\textbf{Q.23} \quad {\footnotesize\textsc{[NAT $\cdot$ 2 Mark(s) $\cdot$ No Neg]}} \hfill {\footnotesize\textbf{GATE 2024 CS2} \quad \tiny\texttt{2024\_CS2\_CS\_Q63}}
\vspace{2pt}\hrule height 0.3pt\vspace{4pt}

Let $Z_n$ be the group of integers $\{0, 1, 2, \dots, n-1\}$ with addition modulo $n$ as the group operation. The number of elements in the group $Z_2 \times Z_3 \times Z_4$ that are their own inverses is \_\_\_\_\_\_\. (Answer in integer)

\par\vspace{6pt}\noindent\textbf{Your Answer:} \fbox{\rule{0pt}{1.3ex}\hspace{3cm}}

\tcblower

\textbf{Official Key:} \textbf{\large 4 to 4} \hfill \textit{\small [Easy]}
\vspace{2pt}\hrule height 0.3pt\vspace{3pt}
\textbf{Topic:} Discrete Mathematics \\
\textbf{Subtopic:} \textit{Monoids, Groups}
\vspace{2pt}\hrule height 0.3pt\vspace{3pt}
\textbf{Guide / Hint:} \\
{\small In an additive group, an element $x$ is its own inverse if $x + x \equiv 0 \pmod n \implies 2x \equiv 0 \pmod n$. In $\mathbb{Z}_2$, elements satisfying $2x \equiv 0$ are $\{0, 1\}$ (2 elements). In $\mathbb{Z}_3$, only $\{0\}$ satisfies $2x \equiv 0 \pmod 3$ (1 element). In $\mathbb{Z}_4$, $\{0, 2\}$ satisfy $2x \equiv 0 \pmod 4$ (2 elements). In the direct product group $\mathbb{Z}_2 \times \mathbb{Z}_3 \times \mathbb{Z}_4$, total self-inverse elements $= 2 \times 1 \times 2 = 4$.}

\end{tcolorbox}
\vspace{6pt}\hrule height 0.35pt\vspace{6pt}
\noindent
\begin{tcolorbox}[
    blanker,
    sidebyside,
    sidebyside align=top,
    lefthand width=0.55\textwidth,
    top=4pt, bottom=4pt, left=0pt, right=0pt,
    middle=16pt,
    underlay={
        \draw[dashed, black!80, line width=0.8pt, dash pattern=on 4pt off 3pt]
            ([xshift=0.55\textwidth+8pt]frame.north west) --
            ([xshift=0.55\textwidth+8pt]frame.south west);
    }
]
\textbf{Q.24} \quad {\footnotesize\textsc{[NAT $\cdot$ 1 Mark(s) $\cdot$ No Neg]}} \hfill {\footnotesize\textbf{GATE 2025 CS1} \quad \tiny\texttt{2025\_CS1\_CS\_Q30}}
\vspace{2pt}\hrule height 0.3pt\vspace{4pt}

Let $S$ be the set of all ternary strings defined over the alphabet $\{a, b, c\}$. Consider all strings in $S$ that contain at least one occurrence of two consecutive symbols, that is, ``aa'', ``bb'' or ``cc''. The number of such strings of length 5 that are possible is \_\_\_\_\_\_\. (Answer in integer)

\par\vspace{6pt}\noindent\textbf{Your Answer:} \fbox{\rule{0pt}{1.3ex}\hspace{3cm}}

\tcblower

\textbf{Official Key:} \textbf{\large 195 to 195} \hfill \textit{\small [Medium]}
\vspace{2pt}\hrule height 0.3pt\vspace{3pt}
\textbf{Topic:} Discrete Mathematics \\
\textbf{Subtopic:} \textit{Combinatorics: counting, recurrence relations, generating functions}
\vspace{2pt}\hrule height 0.3pt\vspace{3pt}
\textbf{Guide / Hint:} \\
{\small The total number of ternary strings of length 5 is $3^5 = 243$. A string has no consecutive identical symbols if the first character has 3 choices and each subsequent character has 2 choices (different from the preceding character), giving $3 \times 2^4 = 48$ such strings. By complementary counting, the number of strings with at least one consecutive pair is $243 - 48 = 195$.}

\end{tcolorbox}
\vspace{6pt}\hrule height 0.35pt\vspace{6pt}
\vspace{8pt}
\subsection*{Digital Logic}
\vspace{-4pt}
\noindent
\begin{tcolorbox}[
    blanker,
    sidebyside,
    sidebyside align=top,
    lefthand width=0.55\textwidth,
    top=4pt, bottom=4pt, left=0pt, right=0pt,
    middle=16pt,
    underlay={
        \draw[dashed, black!80, line width=0.8pt, dash pattern=on 4pt off 3pt]
            ([xshift=0.55\textwidth+8pt]frame.north west) --
            ([xshift=0.55\textwidth+8pt]frame.south west);
    }
]
\textbf{Q.25} \quad {\footnotesize\textsc{[NAT $\cdot$ 1 Mark(s) $\cdot$ No Neg]}} \hfill {\footnotesize\textbf{GATE 2023 CS} \quad \tiny\texttt{2023\_CS\_CS\_Q32}}
\vspace{2pt}\hrule height 0.3pt\vspace{4pt}

A particular number is written as 132 in radix-4 representation. The same number in radix-5 representation is \_\_\_\_\_\_\_\_\. (Answer in integer)

\par\vspace{6pt}\noindent\textbf{Your Answer:} \fbox{\rule{0pt}{1.3ex}\hspace{3cm}}

\tcblower

\textbf{Official Key:} \textbf{\large 110 to 110} \hfill \textit{\small [Easy]}
\vspace{2pt}\hrule height 0.3pt\vspace{3pt}
\textbf{Topic:} Digital Logic \\
\textbf{Subtopic:} \textit{Number representations and computer arithmetic (fixed and floating point)}
\vspace{2pt}\hrule height 0.3pt\vspace{3pt}
\textbf{Guide / Hint:} \\
{\small Converting $132_4$ to base 10: $1 \times 4^2 + 3 \times 4^1 + 2 \times 4^0 = 16 + 12 + 2 = 30_{10}$. Converting $30_{10}$ to radix 5: $30 = 6 \times 5 + 0$, $6 = 1 \times 5 + 1$, $1 = 0 \times 5 + 1$. Reading remainders from bottom to top yields $110_5$.}

\end{tcolorbox}
\vspace{6pt}\hrule height 0.35pt\vspace{6pt}
\noindent
\begin{tcolorbox}[
    blanker,
    sidebyside,
    sidebyside align=top,
    lefthand width=0.55\textwidth,
    top=4pt, bottom=4pt, left=0pt, right=0pt,
    middle=16pt,
    underlay={
        \draw[dashed, black!80, line width=0.8pt, dash pattern=on 4pt off 3pt]
            ([xshift=0.55\textwidth+8pt]frame.north west) --
            ([xshift=0.55\textwidth+8pt]frame.south west);
    }
]
\textbf{Q.26} \quad {\footnotesize\textsc{[NAT $\cdot$ 2 Mark(s) $\cdot$ No Neg]}} \hfill {\footnotesize\textbf{GATE 2025 CS1} \quad \tiny\texttt{2025\_CS1\_CS\_Q59}}
\vspace{2pt}\hrule height 0.3pt\vspace{4pt}

Consider a finite state machine (FSM) with one input $X$ and one output $f$, represented by the given state transition table:
\par\medskip
\begin{tabular}{|c|cc|cc|}
\hline
\textbf{Present state} & \multicolumn{2}{c|}{\textbf{Next state}} & \multicolumn{2}{c|}{\textbf{Output } $f$} \\
& $X=0$ & $X=1$ & $X=0$ & $X=1$ \\
\hline
A & F & B & 0 & 0 \\
B & D & C & 0 & 0 \\
C & F & E & 0 & 0 \\
D & G & A & 1 & 0 \\
E & D & C & 0 & 0 \\
F & F & B & 1 & 1 \\
G & G & H & 0 & 1 \\
H & G & A & 1 & 0 \\
\hline
\end{tabular}
\par\medskip
The minimum number of states required to realize this FSM is \_\_\_\_\_\_\. (Answer in integer)

\par\vspace{6pt}\noindent\textbf{Your Answer:} \fbox{\rule{0pt}{1.3ex}\hspace{3cm}}

\tcblower

\textbf{Official Key:} \textbf{\large 5 to 5} \hfill \textit{\small [Hard]}
\vspace{2pt}\hrule height 0.3pt\vspace{3pt}
\textbf{Topic:} Digital Logic \\
\textbf{Subtopic:} \textit{Combinational and sequential circuits}
\vspace{2pt}\hrule height 0.3pt\vspace{3pt}
\textbf{Guide / Hint:} \\
{\small Using state equivalence partitioning: initial partition by outputs groups states into $P_0 = \{(A, B, C, E), (D, H), (F), (G)\}$. Comparing transitions on $X=0, 1$: states $B$ and $E$ have identical transitions to $(D, C)$, while $A$ and $C$ transition to $(F, B)$ and $(F, E)$ (which are equivalent since $B \equiv E$). States $D$ and $H$ have identical next states $(G, A)$. Thus, equivalence classes are $\{A, C\}, \{B, E\}, \{D, H\}, \{F\}, \{G\}$, resulting in exactly 5 minimal states.}

\end{tcolorbox}
\vspace{6pt}\hrule height 0.35pt\vspace{6pt}
\noindent
\begin{tcolorbox}[
    blanker,
    sidebyside,
    sidebyside align=top,
    lefthand width=0.55\textwidth,
    top=4pt, bottom=4pt, left=0pt, right=0pt,
    middle=16pt,
    underlay={
        \draw[dashed, black!80, line width=0.8pt, dash pattern=on 4pt off 3pt]
            ([xshift=0.55\textwidth+8pt]frame.north west) --
            ([xshift=0.55\textwidth+8pt]frame.south west);
    }
]
\textbf{Q.27} \quad {\footnotesize\textsc{[NAT $\cdot$ 2 Mark(s) $\cdot$ No Neg]}} \hfill {\footnotesize\textbf{GATE 2025 CS1} \quad \tiny\texttt{2025\_CS1\_CS\_Q60}}
\vspace{2pt}\hrule height 0.3pt\vspace{4pt}

Consider the given sequential circuit designed using D-Flip-flops. The circuit is initialized with some value (initial state). The number of distinct states the circuit will go through before returning back to the initial state is \_\_\_\_\_\_\. (Answer in integer)

\par\vspace{6pt}\noindent\textbf{Your Answer:} \fbox{\rule{0pt}{1.3ex}\hspace{3cm}}

\tcblower

\textbf{Official Key:} \textbf{\large 7 to 8} \hfill \textit{\small [Hard]}
\vspace{2pt}\hrule height 0.3pt\vspace{3pt}
\textbf{Topic:} Digital Logic \\
\textbf{Subtopic:} \textit{Combinational and sequential circuits}
\vspace{2pt}\hrule height 0.3pt\vspace{3pt}
\textbf{Guide / Hint:} \\
{\small Tracing the characteristic shift register feedback transitions where the inputs to the D flip-flops are determined by logic feedback of the current state $(Q_3, Q_2, Q_1, Q_0)$, the state cycle length forms a closed loop of period 7 (or 8 depending on the initial valid cycle).}

\end{tcolorbox}
\vspace{6pt}\hrule height 0.35pt\vspace{6pt}
\noindent
\begin{tcolorbox}[
    blanker,
    sidebyside,
    sidebyside align=top,
    lefthand width=0.55\textwidth,
    top=4pt, bottom=4pt, left=0pt, right=0pt,
    middle=16pt,
    underlay={
        \draw[dashed, black!80, line width=0.8pt, dash pattern=on 4pt off 3pt]
            ([xshift=0.55\textwidth+8pt]frame.north west) --
            ([xshift=0.55\textwidth+8pt]frame.south west);
    }
]
\textbf{Q.28} \quad {\footnotesize\textsc{[NAT $\cdot$ 1 Mark(s) $\cdot$ No Neg]}} \hfill {\footnotesize\textbf{GATE 2025 CS2} \quad \tiny\texttt{2025\_CS2\_CS\_Q34}}
\vspace{2pt}\hrule height 0.3pt\vspace{4pt}

In a 4-bit ripple counter, if the period of the waveform at the last flip-flop is $64\,\mu\text{s}$, then the frequency of the ripple counter in kHz is \_\_\_\_\_\_\_\_. (Answer in integer)

\par\vspace{6pt}\noindent\textbf{Your Answer:} \fbox{\rule{0pt}{1.3ex}\hspace{3cm}}

\tcblower

\textbf{Official Key:} \textbf{\large 250 to 250} \hfill \textit{\small [Medium]}
\vspace{2pt}\hrule height 0.3pt\vspace{3pt}
\textbf{Topic:} Digital Logic \\
\textbf{Subtopic:} \textit{Combinational and sequential circuits: ripple counters}
\vspace{2pt}\hrule height 0.3pt\vspace{3pt}
\textbf{Guide / Hint:} \\
{\small A 4-bit ripple counter divides the clock frequency by $2^4 = 16$. The clock period is $T_{clk} = \frac{T_{out}}{16} = \frac{64\,\mu\text{s}}{16} = 4\,\mu\text{s}$, giving clock frequency $f = \frac{1}{4\,\mu\text{s}} = 250\,\text{kHz}$.}

\end{tcolorbox}
\vspace{6pt}\hrule height 0.35pt\vspace{6pt}
\vspace{8pt}
\subsection*{Computer Organization and Architecture}
\vspace{-4pt}
\noindent
\begin{tcolorbox}[
    blanker,
    sidebyside,
    sidebyside align=top,
    lefthand width=0.55\textwidth,
    top=4pt, bottom=4pt, left=0pt, right=0pt,
    middle=16pt,
    underlay={
        \draw[dashed, black!80, line width=0.8pt, dash pattern=on 4pt off 3pt]
            ([xshift=0.55\textwidth+8pt]frame.north west) --
            ([xshift=0.55\textwidth+8pt]frame.south west);
    }
]
\textbf{Q.29} \quad {\footnotesize\textsc{[NAT $\cdot$ 1 Mark(s) $\cdot$ No Neg]}} \hfill {\footnotesize\textbf{GATE 2018 CS} \quad \tiny\texttt{2018\_CS\_GA\_Q23}}
\vspace{2pt}\hrule height 0.3pt\vspace{4pt}

A 32-bit wide main memory unit with a capacity of 1 GB is built using $256\text{M} \times 4\text{-bit}$ DRAM chips. The number of rows of memory cells in the DRAM chip is $2^{14}$. The time taken to perform one refresh operation is 50 nanoseconds. The refresh period is 2 milliseconds. The percentage of time that the memory is available for performing the normal read and write operations is \_\_\_\_\_\_\_\_\%. (Answer in integer)

\par\vspace{6pt}\noindent\textbf{Your Answer:} \fbox{\rule{0pt}{1.3ex}\hspace{3cm}}

\tcblower

\textbf{Official Key:} \textbf{\large 59.0 to 60.0} \hfill \textit{\small [Medium]}
\vspace{2pt}\hrule height 0.3pt\vspace{3pt}
\textbf{Topic:} Computer Organization and Architecture \\
\textbf{Subtopic:} \textit{Memory hierarchy: cache, main memory and secondary storage}
\vspace{2pt}\hrule height 0.3pt\vspace{3pt}
\textbf{Guide / Hint:} \\
{\small The DRAM has $2^{14} = 16384$ rows. Refreshing all rows takes $16384 \times 50\text{ ns} = 819.2\ \mu\text{s}$ every $2000\ \mu\text{s}$. The refresh overhead fraction is $\frac{819.2}{2000} = 40.96\%$. Hence, the percentage of time available for reading and writing memory is $100\% - 40.96\% = 59.04\% \approx 59\%$.}

\end{tcolorbox}
\vspace{6pt}\hrule height 0.35pt\vspace{6pt}
\noindent
\begin{tcolorbox}[
    blanker,
    sidebyside,
    sidebyside align=top,
    lefthand width=0.55\textwidth,
    top=4pt, bottom=4pt, left=0pt, right=0pt,
    middle=16pt,
    underlay={
        \draw[dashed, black!80, line width=0.8pt, dash pattern=on 4pt off 3pt]
            ([xshift=0.55\textwidth+8pt]frame.north west) --
            ([xshift=0.55\textwidth+8pt]frame.south west);
    }
]
\textbf{Q.30} \quad {\footnotesize\textsc{[NAT $\cdot$ 2 Mark(s) $\cdot$ No Neg]}} \hfill {\footnotesize\textbf{GATE 2018 CS} \quad \tiny\texttt{2018\_CS\_GA\_Q50}}
\vspace{2pt}\hrule height 0.3pt\vspace{4pt}

The instruction pipeline of a RISC processor has five stages: Instruction Fetch (IF), Instruction Decode (ID), Perform Operation (PO), Operand Read (OR), and Write Back (WB). Each stage takes 1 clock cycle except PO. For a sequence of 100 instructions, 40 instructions require 3 clock cycles in PO, 35 instructions require 2 clock cycles in PO, and the remaining 25 instructions require 1 clock cycle in PO. Assuming no memory stalls or branch delays, the total number of clock cycles to execute the 100 instructions is \_\_\_\_\_\_\_\_.

\par\vspace{6pt}\noindent\textbf{Your Answer:} \fbox{\rule{0pt}{1.3ex}\hspace{3cm}}

\tcblower

\textbf{Official Key:} \textbf{\large 219 to 219} \hfill \textit{\small [Medium]}
\vspace{2pt}\hrule height 0.3pt\vspace{3pt}
\textbf{Topic:} Computer Organization and Architecture \\
\textbf{Subtopic:} \textit{Instruction pipelining, pipeline hazards}
\vspace{2pt}\hrule height 0.3pt\vspace{3pt}
\textbf{Guide / Hint:} \\
{\small For 100 instructions in a 5-stage pipeline without branch delays, ideal cycles $= 5 + (100 - 1) = 104$. The PO stage causes pipeline stalls when operand evaluation takes extra cycles: 40 instructions require 3 cycles (2 extra cycles $\implies 40 \times 2 = 80$ stall cycles) and 35 instructions require 2 cycles (1 extra cycle $\implies 35 \times 1 = 35$ stall cycles). Total clock cycles $= 104 + 80 + 35 = 219$.}

\end{tcolorbox}
\vspace{6pt}\hrule height 0.35pt\vspace{6pt}
\noindent
\begin{tcolorbox}[
    blanker,
    sidebyside,
    sidebyside align=top,
    lefthand width=0.55\textwidth,
    top=4pt, bottom=4pt, left=0pt, right=0pt,
    middle=16pt,
    underlay={
        \draw[dashed, black!80, line width=0.8pt, dash pattern=on 4pt off 3pt]
            ([xshift=0.55\textwidth+8pt]frame.north west) --
            ([xshift=0.55\textwidth+8pt]frame.south west);
    }
]
\textbf{Q.31} \quad {\footnotesize\textsc{[NAT $\cdot$ 2 Mark(s) $\cdot$ No Neg]}} \hfill {\footnotesize\textbf{GATE 2018 CS} \quad \tiny\texttt{2018\_CS\_GA\_Q51}}
\vspace{2pt}\hrule height 0.3pt\vspace{4pt}

A processor has 16 integer registers (R0 to R15) and 64 floating point registers (F0 to F63). It uses a 16-bit instruction format. Four types of instructions are supported:
\begin{itemize}[leftmargin=*, nosep]
    \item Type-1: Opcode with three integer register operands
    \item Type-2: Opcode with two floating point register operands
    \item Type-3: Opcode with one integer and one floating point register operand
    \item Type-4: Opcode with one floating point register operand
\end{itemize}
If there are 4 Type-1, 8 Type-2, and 14 Type-3 instructions, the maximum number of Type-4 instructions that can be supported is \_\_\_\_\_\_\_\_.

\par\vspace{6pt}\noindent\textbf{Your Answer:} \fbox{\rule{0pt}{1.3ex}\hspace{3cm}}

\tcblower

\textbf{Official Key:} \textbf{\large 32 to 32} \hfill \textit{\small [Hard]}
\vspace{2pt}\hrule height 0.3pt\vspace{3pt}
\textbf{Topic:} Computer Organization and Architecture \\
\textbf{Subtopic:} \textit{Machine instructions and addressing modes}
\vspace{2pt}\hrule height 0.3pt\vspace{3pt}
\textbf{Guide / Hint:} \\
{\small A 16-bit instruction space has $2^{16} = 65536$ encodings. Type-1: 4 instructions with $3 \times 4\text{-bit}$ registers consume $4 \times 2^{12} = 16384$. Type-2: 8 instructions with $2 \times 6\text{-bit}$ registers consume $8 \times 2^{12} = 32768$. Type-3: 14 instructions with 4-bit and 6-bit registers consume $14 \times 2^{10} = 14336$. Remaining encodings $= 65536 - (16384 + 32768 + 14336) = 2048$. Each Type-4 instruction uses one 6-bit register ($2^6 = 64$ encodings). Thus, maximum $N = 2048 / 64 = 32$.}

\end{tcolorbox}
\vspace{6pt}\hrule height 0.35pt\vspace{6pt}
\noindent
\begin{tcolorbox}[
    blanker,
    sidebyside,
    sidebyside align=top,
    lefthand width=0.55\textwidth,
    top=4pt, bottom=4pt, left=0pt, right=0pt,
    middle=16pt,
    underlay={
        \draw[dashed, black!80, line width=0.8pt, dash pattern=on 4pt off 3pt]
            ([xshift=0.55\textwidth+8pt]frame.north west) --
            ([xshift=0.55\textwidth+8pt]frame.south west);
    }
]
\textbf{Q.32} \quad {\footnotesize\textsc{[NAT $\cdot$ 1 Mark(s) $\cdot$ No Neg]}} \hfill {\footnotesize\textbf{GATE 2022 CS} \quad \tiny\texttt{2022\_CS\_CS\_Q33}}
\vspace{2pt}\hrule height 0.3pt\vspace{4pt}

A cache memory that has a hit rate of 0.8 has an access latency 10 ns and miss penalty 100 ns. An optimization is done on the cache to reduce the miss rate. However, the optimization results in an increase of cache access latency to 15 ns, whereas the miss penalty is not affected. The minimum hit rate (rounded off to two decimal places) needed after the optimization such that it should not increase the average memory access time is \_\_\_\_\_\_\_\_\_\_\_\. (Answer in numeric format)

\par\vspace{6pt}\noindent\textbf{Your Answer:} \fbox{\rule{0pt}{1.3ex}\hspace{3cm}}

\tcblower

\textbf{Official Key:} \textbf{\large 0.85 to 0.85} \hfill \textit{\small [Medium]}
\vspace{2pt}\hrule height 0.3pt\vspace{3pt}
\textbf{Topic:} Computer Organization and Architecture \\
\textbf{Subtopic:} \textit{Memory hierarchy: cache, main memory and secondary storage}
\vspace{2pt}\hrule height 0.3pt\vspace{3pt}
\textbf{Guide / Hint:} \\
{\small Initial AMAT is $10 + (1 - 0.80) \times 100 = 10 + 20 = 30\text{ ns}$. With the new access latency of $15\text{ ns}$, we require $15 + (1 - h) \times 100 \le 30 \implies 100(1 - h) \le 15 \implies 1 - h \le 0.15 \implies h \ge 0.85$.}

\end{tcolorbox}
\vspace{6pt}\hrule height 0.35pt\vspace{6pt}
\noindent
\begin{tcolorbox}[
    blanker,
    sidebyside,
    sidebyside align=top,
    lefthand width=0.55\textwidth,
    top=4pt, bottom=4pt, left=0pt, right=0pt,
    middle=16pt,
    underlay={
        \draw[dashed, black!80, line width=0.8pt, dash pattern=on 4pt off 3pt]
            ([xshift=0.55\textwidth+8pt]frame.north west) --
            ([xshift=0.55\textwidth+8pt]frame.south west);
    }
]
\textbf{Q.33} \quad {\footnotesize\textsc{[NAT $\cdot$ 2 Mark(s) $\cdot$ No Neg]}} \hfill {\footnotesize\textbf{GATE 2022 CS} \quad \tiny\texttt{2022\_CS\_CS\_Q61}}
\vspace{2pt}\hrule height 0.3pt\vspace{4pt}

A processor X1 operating at 2 GHz has a standard 5-stage RISC instruction pipeline having a base CPI (cycles per instruction) of one without any pipeline hazards. For a given program P that has 30\% branch instructions, control hazards incur 2 cycles stall for every branch. A new version of the processor X2 operating at same clock frequency has an additional branch predictor unit (BPU) that completely eliminates stalls for correctly predicted branches. There is neither any savings nor any additional stalls for wrong predictions. There are no structural hazards and data hazards for X1 and X2. If the BPU has a prediction accuracy of 80\%, the speed up (rounded off to two decimal places) obtained by X2 over X1 in executing P is \_\_\_\_\_\_\_\_\_\_\_\_\. (Answer in numeric format)

\par\vspace{6pt}\noindent\textbf{Your Answer:} \fbox{\rule{0pt}{1.3ex}\hspace{3cm}}

\tcblower

\textbf{Official Key:} \textbf{\large 1.42 to 1.45} \hfill \textit{\small [Medium]}
\vspace{2pt}\hrule height 0.3pt\vspace{3pt}
\textbf{Topic:} Computer Organization and Architecture \\
\textbf{Subtopic:} \textit{Instruction pipelining, pipeline hazards}
\vspace{2pt}\hrule height 0.3pt\vspace{3pt}
\textbf{Guide / Hint:} \\
{\small For $X_1$: $\text{CPI}_1 = 1 + (0.30 \times 2) = 1.60$. For $X_2$: with 80\% branch prediction accuracy, only 20\% of branches stall: $\text{CPI}_2 = 1 + (0.30 \times 0.20 \times 2) = 1.12$. Speedup $= \frac{\text{CPI}_1}{\text{CPI}_2} = \frac{1.60}{1.12} \approx 1.43$.}

\end{tcolorbox}
\vspace{6pt}\hrule height 0.35pt\vspace{6pt}
\noindent
\begin{tcolorbox}[
    blanker,
    sidebyside,
    sidebyside align=top,
    lefthand width=0.55\textwidth,
    top=4pt, bottom=4pt, left=0pt, right=0pt,
    middle=16pt,
    underlay={
        \draw[dashed, black!80, line width=0.8pt, dash pattern=on 4pt off 3pt]
            ([xshift=0.55\textwidth+8pt]frame.north west) --
            ([xshift=0.55\textwidth+8pt]frame.south west);
    }
]
\textbf{Q.34} \quad {\footnotesize\textsc{[NAT $\cdot$ 1 Mark(s) $\cdot$ No Neg]}} \hfill {\footnotesize\textbf{GATE 2023 CS} \quad \tiny\texttt{2023\_CS\_CS\_Q33}}
\vspace{2pt}\hrule height 0.3pt\vspace{4pt}

Consider a 3-stage pipelined processor having a delay of 10 ns (nanoseconds), 20 ns, and 14 ns, for the first, second, and the third stages, respectively. Assume that there is no other delay and the processor does not suffer from any pipeline hazards. Also assume that one instruction is fetched every cycle.
The total execution time for executing 100 instructions on this processor is \_\_\_\_\_\_\_\_\_ ns. (Answer in integer)

\par\vspace{6pt}\noindent\textbf{Your Answer:} \fbox{\rule{0pt}{1.3ex}\hspace{3cm}}

\tcblower

\textbf{Official Key:} \textbf{\large 2040 to 2040} \hfill \textit{\small [Easy]}
\vspace{2pt}\hrule height 0.3pt\vspace{3pt}
\textbf{Topic:} Computer Organization and Architecture \\
\textbf{Subtopic:} \textit{Instruction pipelining, pipeline hazards}
\vspace{2pt}\hrule height 0.3pt\vspace{3pt}
\textbf{Guide / Hint:} \\
{\small The clock period of the pipeline is governed by the longest stage delay: $T = \max(10, 20, 14) = 20\text{ ns}$. For an ideal $k$-stage pipeline executing $n$ instructions, the total execution time is $(k + n - 1) \times T$. For $k = 3$ and $n = 100$: $(3 + 100 - 1) \times 20 = 102 \times 20 = 2040\text{ ns}$.}

\end{tcolorbox}
\vspace{6pt}\hrule height 0.35pt\vspace{6pt}
\noindent
\begin{tcolorbox}[
    blanker,
    sidebyside,
    sidebyside align=top,
    lefthand width=0.55\textwidth,
    top=4pt, bottom=4pt, left=0pt, right=0pt,
    middle=16pt,
    underlay={
        \draw[dashed, black!80, line width=0.8pt, dash pattern=on 4pt off 3pt]
            ([xshift=0.55\textwidth+8pt]frame.north west) --
            ([xshift=0.55\textwidth+8pt]frame.south west);
    }
]
\textbf{Q.35} \quad {\footnotesize\textsc{[NAT $\cdot$ 1 Mark(s) $\cdot$ No Neg]}} \hfill {\footnotesize\textbf{GATE 2023 CS} \quad \tiny\texttt{2023\_CS\_CS\_Q34}}
\vspace{2pt}\hrule height 0.3pt\vspace{4pt}

A keyboard connected to a computer is used at a rate of 1 keystroke per second. The computer system polls the keyboard every 10 ms (milli seconds) to check for a keystroke and consumes 100 $\mu$s (micro seconds) for each poll. If it is determined after polling that a key has been pressed, the system consumes an additional 200 $\mu$s to process the keystroke. Let $T_1$ denote the fraction of a second spent in polling and processing a keystroke.
In an alternative implementation, the system uses interrupts instead of polling. An interrupt is raised for every keystroke. It takes a total of 1 ms for servicing an interrupt and processing a keystroke. Let $T_2$ denote the fraction of a second spent in servicing the interrupt and processing a keystroke.
The ratio $\frac{T_1}{T_2}$ is \_\_\_\_\_\_\_\_\. (Rounded off to one decimal place)

\par\vspace{6pt}\noindent\textbf{Your Answer:} \fbox{\rule{0pt}{1.3ex}\hspace{3cm}}

\tcblower

\textbf{Official Key:} \textbf{\large 10.2 to 10.2} \hfill \textit{\small [Medium]}
\vspace{2pt}\hrule height 0.3pt\vspace{3pt}
\textbf{Topic:} Computer Organization and Architecture \\
\textbf{Subtopic:} \textit{I/O interface: interrupt and DMA mode}
\vspace{2pt}\hrule height 0.3pt\vspace{3pt}
\textbf{Guide / Hint:} \\
{\small With polling every $10\text{ ms}$, there are 100 polls per second. Polling overhead is $100 \times 100\ \mu\text{s} = 10000\ \mu\text{s} = 10\text{ ms}$. With 1 keystroke/s, processing takes $200\ \mu\text{s} = 0.2\text{ ms}$, giving $T_1 = 10 + 0.2 = 10.2\text{ ms}$. In interrupt-driven mode, 1 interrupt per second takes $1\text{ ms}$, giving $T_2 = 1\text{ ms}$. The ratio is $\frac{T_1}{T_2} = \frac{10.2}{1.0} = 10.2$.}

\end{tcolorbox}
\vspace{6pt}\hrule height 0.35pt\vspace{6pt}
\noindent
\begin{tcolorbox}[
    blanker,
    sidebyside,
    sidebyside align=top,
    lefthand width=0.55\textwidth,
    top=4pt, bottom=4pt, left=0pt, right=0pt,
    middle=16pt,
    underlay={
        \draw[dashed, black!80, line width=0.8pt, dash pattern=on 4pt off 3pt]
            ([xshift=0.55\textwidth+8pt]frame.north west) --
            ([xshift=0.55\textwidth+8pt]frame.south west);
    }
]
\textbf{Q.36} \quad {\footnotesize\textsc{[NAT $\cdot$ 2 Mark(s) $\cdot$ No Neg]}} \hfill {\footnotesize\textbf{GATE 2023 CS} \quad \tiny\texttt{2023\_CS\_CS\_Q64}}
\vspace{2pt}\hrule height 0.3pt\vspace{4pt}

An 8-way set associative cache of size 64 KB (1 KB = 1024 bytes) is used in a system with 32-bit address. The address is sub-divided into TAG, INDEX, and BLOCK OFFSET.
The number of bits in the TAG is \_\_\_\_\_\_\_\_\. (Answer in integer)

\par\vspace{6pt}\noindent\textbf{Your Answer:} \fbox{\rule{0pt}{1.3ex}\hspace{3cm}}

\tcblower

\textbf{Official Key:} \textbf{\large 19 to 19} \hfill \textit{\small [Easy]}
\vspace{2pt}\hrule height 0.3pt\vspace{3pt}
\textbf{Topic:} Computer Organization and Architecture \\
\textbf{Subtopic:} \textit{Memory hierarchy: cache, main memory and secondary storage}
\vspace{2pt}\hrule height 0.3pt\vspace{3pt}
\textbf{Guide / Hint:} \\
{\small Address size $= 32\text{ bits}$. Cache size $= 64\text{ KB} = 2^{16}\text{ bytes}$. For an 8-way set-associative cache, the combined capacity of index and offset is $\frac{\text{Cache Size}}{\text{Associativity}} = \frac{64\text{ KB}}{8} = 8\text{ KB} = 2^{13}\text{ bytes}$, requiring $13\text{ bits}$ for INDEX + BLOCK OFFSET. Therefore, the number of TAG bits is $32 - 13 = 19\text{ bits}$.}

\end{tcolorbox}
\vspace{6pt}\hrule height 0.35pt\vspace{6pt}
\noindent
\begin{tcolorbox}[
    blanker,
    sidebyside,
    sidebyside align=top,
    lefthand width=0.55\textwidth,
    top=4pt, bottom=4pt, left=0pt, right=0pt,
    middle=16pt,
    underlay={
        \draw[dashed, black!80, line width=0.8pt, dash pattern=on 4pt off 3pt]
            ([xshift=0.55\textwidth+8pt]frame.north west) --
            ([xshift=0.55\textwidth+8pt]frame.south west);
    }
]
\textbf{Q.37} \quad {\footnotesize\textsc{[NAT $\cdot$ 2 Mark(s) $\cdot$ No Neg]}} \hfill {\footnotesize\textbf{GATE 2024 CS1} \quad \tiny\texttt{2024\_CS1\_CS\_Q54}}
\vspace{2pt}\hrule height 0.3pt\vspace{4pt}

Consider a 512 GB hard disk with 32 storage surfaces. There are 4096 sectors per track and each sector holds 1024 bytes of data. The number of cylinders in the hard disk is \_\_\_\_\_\_\_\_\_\. (Answer in integer)

\par\vspace{6pt}\noindent\textbf{Your Answer:} \fbox{\rule{0pt}{1.3ex}\hspace{3cm}}

\tcblower

\textbf{Official Key:} \textbf{\large 4096 to 4096} \hfill \textit{\small [Easy]}
\vspace{2pt}\hrule height 0.3pt\vspace{3pt}
\textbf{Topic:} Computer Organization and Architecture \\
\textbf{Subtopic:} \textit{Memory hierarchy: cache, main memory and secondary storage}
\vspace{2pt}\hrule height 0.3pt\vspace{3pt}
\textbf{Guide / Hint:} \\
{\small Disk capacity $= 512\text{ GB} = 512 \times 2^{30}\text{ bytes} = 2^{39}$ bytes. One cylinder comprises 32 surfaces $\times 4096$ sectors/track $\times 1024$ bytes/sector $= 2^5 \times 2^{12} \times 2^{10} = 2^{27}$ bytes. Number of cylinders $= \frac{2^{39}}{2^{27}} = 2^{12} = 4096$.}

\end{tcolorbox}
\vspace{6pt}\hrule height 0.35pt\vspace{6pt}
\noindent
\begin{tcolorbox}[
    blanker,
    sidebyside,
    sidebyside align=top,
    lefthand width=0.55\textwidth,
    top=4pt, bottom=4pt, left=0pt, right=0pt,
    middle=16pt,
    underlay={
        \draw[dashed, black!80, line width=0.8pt, dash pattern=on 4pt off 3pt]
            ([xshift=0.55\textwidth+8pt]frame.north west) --
            ([xshift=0.55\textwidth+8pt]frame.south west);
    }
]
\textbf{Q.38} \quad {\footnotesize\textsc{[NAT $\cdot$ 2 Mark(s) $\cdot$ No Neg]}} \hfill {\footnotesize\textbf{GATE 2024 CS1} \quad \tiny\texttt{2024\_CS1\_CS\_Q55}}
\vspace{2pt}\hrule height 0.3pt\vspace{4pt}

The baseline execution time of a program on a 2 GHz single core machine is 100 nanoseconds (ns). The code corresponding to 90\% of the execution time can be fully parallelized. The overhead for using an additional core is 10 ns when running on a multicore system. Assume that all cores in the multicore system run their share of the parallelized code for an equal amount of time.
The number of cores that minimize the execution time of the program is \_\_\_\_\_\_\_\_\. (Answer in integer)

\par\vspace{6pt}\noindent\textbf{Your Answer:} \fbox{\rule{0pt}{1.3ex}\hspace{3cm}}

\tcblower

\textbf{Official Key:} \textbf{\large 3 to 3} \hfill \textit{\small [Medium]}
\vspace{2pt}\hrule height 0.3pt\vspace{3pt}
\textbf{Topic:} Computer Organization and Architecture \\
\textbf{Subtopic:} \textit{Instruction pipelining, pipeline hazards}
\vspace{2pt}\hrule height 0.3pt\vspace{3pt}
\textbf{Guide / Hint:} \\
{\small Baseline runtime is $100\text{ ns}$, with $90\text{ ns}$ parallelizable and $10\text{ ns}$ serial. On $p$ cores, the parallel portion takes $\frac{90}{p}\text{ ns}$ with core overhead of $10(p-1)\text{ ns}$. Total execution time $T(p) = 10 + \frac{90}{p} + 10(p-1) = \frac{90}{p} + 10p$. Evaluating: $T(1) = 100\text{ ns}$, $T(2) = 65\text{ ns}$, $T(3) = 60\text{ ns}$, $T(4) = 62.5\text{ ns}$. The time is minimized at $p = 3$ cores.}

\end{tcolorbox}
\vspace{6pt}\hrule height 0.35pt\vspace{6pt}
\noindent
\begin{tcolorbox}[
    blanker,
    sidebyside,
    sidebyside align=top,
    lefthand width=0.55\textwidth,
    top=4pt, bottom=4pt, left=0pt, right=0pt,
    middle=16pt,
    underlay={
        \draw[dashed, black!80, line width=0.8pt, dash pattern=on 4pt off 3pt]
            ([xshift=0.55\textwidth+8pt]frame.north west) --
            ([xshift=0.55\textwidth+8pt]frame.south west);
    }
]
\textbf{Q.39} \quad {\footnotesize\textsc{[NAT $\cdot$ 2 Mark(s) $\cdot$ No Neg]}} \hfill {\footnotesize\textbf{GATE 2024 CS1} \quad \tiny\texttt{2024\_CS1\_CS\_Q56}}
\vspace{2pt}\hrule height 0.3pt\vspace{4pt}

A given program has 25\% load/store instructions. Suppose the ideal CPI (cycles per instruction) without any memory stalls is 2. The program exhibits 2\% miss rate on instruction cache and 8\% miss rate on data cache. The miss penalty is 100 cycles.
The speedup (rounded off to two decimal places) achieved with a perfect cache (i.e., with NO data or instruction cache misses) is \_\_\_\_\_\_\_\_\_\. (Answer in numeric format)

\par\vspace{6pt}\noindent\textbf{Your Answer:} \fbox{\rule{0pt}{1.3ex}\hspace{3cm}}

\tcblower

\textbf{Official Key:} \textbf{\large 3.00 to 3.00} \hfill \textit{\small [Medium]}
\vspace{2pt}\hrule height 0.3pt\vspace{3pt}
\textbf{Topic:} Computer Organization and Architecture \\
\textbf{Subtopic:} \textit{Memory hierarchy: cache, main memory and secondary storage}
\vspace{2pt}\hrule height 0.3pt\vspace{3pt}
\textbf{Guide / Hint:} \\
{\small Total CPI is given by $\text{CPI} = \text{CPI}_{\text{ideal}} + (\text{I-miss rate} \times \text{penalty}) + (\text{Data-instruction fraction} \times \text{D-miss rate} \times \text{penalty}) = 2 + (0.02 \times 100) + (0.25 \times 0.08 \times 100) = 2 + 2 + 2 = 6$. For a perfect cache with no misses, $\text{CPI}_{\text{perfect}} = \text{CPI}_{\text{ideal}} = 2$. Speedup $= \frac{6}{2} = 3.00$.}

\end{tcolorbox}
\vspace{6pt}\hrule height 0.35pt\vspace{6pt}
\noindent
\begin{tcolorbox}[
    blanker,
    sidebyside,
    sidebyside align=top,
    lefthand width=0.55\textwidth,
    top=4pt, bottom=4pt, left=0pt, right=0pt,
    middle=16pt,
    underlay={
        \draw[dashed, black!80, line width=0.8pt, dash pattern=on 4pt off 3pt]
            ([xshift=0.55\textwidth+8pt]frame.north west) --
            ([xshift=0.55\textwidth+8pt]frame.south west);
    }
]
\textbf{Q.40} \quad {\footnotesize\textsc{[NAT $\cdot$ 2 Mark(s) $\cdot$ No Neg]}} \hfill {\footnotesize\textbf{GATE 2024 CS2} \quad \tiny\texttt{2024\_CS2\_CS\_Q57}}
\vspace{2pt}\hrule height 0.3pt\vspace{4pt}

A processor with 16 general purpose registers uses a 32-bit instruction format. The instruction format consists of an opcode field, an addressing mode field, two register operand fields, and a 16-bit scalar field. If 8 addressing modes are to be supported, the maximum number of unique opcodes possible for every addressing mode is \_\_\_\_\_\_\. (Answer in integer)

\par\vspace{6pt}\noindent\textbf{Your Answer:} \fbox{\rule{0pt}{1.3ex}\hspace{3cm}}

\tcblower

\textbf{Official Key:} \textbf{\large 32 to 32} \hfill \textit{\small [Medium]}
\vspace{2pt}\hrule height 0.3pt\vspace{3pt}
\textbf{Topic:} Computer Organization and Architecture \\
\textbf{Subtopic:} \textit{Machine instructions and addressing modes}
\vspace{2pt}\hrule height 0.3pt\vspace{3pt}
\textbf{Guide / Hint:} \\
{\small Instruction length is 32 bits. Register fields: 16 registers need $\log_2 16 = 4$ bits each ($2 \times 4 = 8$ bits). Addressing mode: 8 modes need $\log_2 8 = 3$ bits. Scalar field: 16 bits. Remaining bits for opcode $= 32 - (3 + 8 + 16) = 5$ bits. The maximum number of unique opcodes possible for every addressing mode is $2^5 = 32$.}

\end{tcolorbox}
\vspace{6pt}\hrule height 0.35pt\vspace{6pt}
\noindent
\begin{tcolorbox}[
    blanker,
    sidebyside,
    sidebyside align=top,
    lefthand width=0.55\textwidth,
    top=4pt, bottom=4pt, left=0pt, right=0pt,
    middle=16pt,
    underlay={
        \draw[dashed, black!80, line width=0.8pt, dash pattern=on 4pt off 3pt]
            ([xshift=0.55\textwidth+8pt]frame.north west) --
            ([xshift=0.55\textwidth+8pt]frame.south west);
    }
]
\textbf{Q.41} \quad {\footnotesize\textsc{[NAT $\cdot$ 2 Mark(s) $\cdot$ No Neg]}} \hfill {\footnotesize\textbf{GATE 2024 CS2} \quad \tiny\texttt{2024\_CS2\_CS\_Q58}}
\vspace{2pt}\hrule height 0.3pt\vspace{4pt}

A non-pipelined instruction execution unit operating at 2 GHz takes an average of 6 cycles to execute an instruction of a program P. The unit is then redesigned to operate on a 5-stage pipeline at 2 GHz. Assume that the ideal throughput of the pipelined unit is 1 instruction per cycle. In the execution of program P, 20\% instructions incur an average of 2 cycles stall due to data hazards and 20\% instructions incur an average of 3 cycles stall due to control hazards. The speedup (rounded off to one decimal place) obtained by the pipelined design over the non-pipelined design is \_\_\_\_\_\_\. (Answer in range 2.9 to 3.1)

\par\vspace{6pt}\noindent\textbf{Your Answer:} \fbox{\rule{0pt}{1.3ex}\hspace{3cm}}

\tcblower

\textbf{Official Key:} \textbf{\large 2.9 to 3.1} \hfill \textit{\small [Medium]}
\vspace{2pt}\hrule height 0.3pt\vspace{3pt}
\textbf{Topic:} Computer Organization and Architecture \\
\textbf{Subtopic:} \textit{Instruction pipelining, pipeline hazards}
\vspace{2pt}\hrule height 0.3pt\vspace{3pt}
\textbf{Guide / Hint:} \\
{\small Average cycles per instruction in non-pipelined design: $\text{CPI}_{\text{non-pipe}} = 6$. In pipelined design: $\text{CPI}_{\text{pipe}} = 1 + (0.20 \times 2) + (0.20 \times 3) = 1 + 0.40 + 0.60 = 2.0$. Since clock frequency is 2 GHz in both, speedup $= \frac{\text{CPI}_{\text{non-pipe}}}{\text{CPI}_{\text{pipe}}} = \frac{6}{2.0} = 3.0$.}

\end{tcolorbox}
\vspace{6pt}\hrule height 0.35pt\vspace{6pt}
\noindent
\begin{tcolorbox}[
    blanker,
    sidebyside,
    sidebyside align=top,
    lefthand width=0.55\textwidth,
    top=4pt, bottom=4pt, left=0pt, right=0pt,
    middle=16pt,
    underlay={
        \draw[dashed, black!80, line width=0.8pt, dash pattern=on 4pt off 3pt]
            ([xshift=0.55\textwidth+8pt]frame.north west) --
            ([xshift=0.55\textwidth+8pt]frame.south west);
    }
]
\textbf{Q.42} \quad {\footnotesize\textsc{[NAT $\cdot$ 2 Mark(s) $\cdot$ No Neg]}} \hfill {\footnotesize\textbf{GATE 2024 CS2} \quad \tiny\texttt{2024\_CS2\_CS\_Q61}}
\vspace{2pt}\hrule height 0.3pt\vspace{4pt}

A processor uses a 32-bit instruction format and supports byte-addressable memory access. The ISA of the processor has 150 distinct instructions. The instructions are equally divided into two types, namely R-type and I-type, whose formats are shown below:
\begin{verbatim}
R-type: OPCODE | UNUSED | DST Register | SRC Register1 | SRC Register 2
I-type: OPCODE | DST Register | SRC Register | Immediate value/address
\end{verbatim}
In the OPCODE, 1 bit distinguishes between I-type and R-type and the remaining bits indicate the operation. The processor has 50 architectural registers, and all register fields are of equal size.
Let $X$ be the bits for UNUSED, $Y$ for OPCODE, and $Z$ for immediate value/address. The value of $X + 2Y + Z$ is \_\_\_\_\_\_\. (Answer in integer)

\par\vspace{6pt}\noindent\textbf{Your Answer:} \fbox{\rule{0pt}{1.3ex}\hspace{3cm}}

\tcblower

\textbf{Official Key:} \textbf{\large 34 to 34} \hfill \textit{\small [Hard]}
\vspace{2pt}\hrule height 0.3pt\vspace{3pt}
\textbf{Topic:} Computer Organization and Architecture \\
\textbf{Subtopic:} \textit{Machine instructions and addressing modes}
\vspace{2pt}\hrule height 0.3pt\vspace{3pt}
\textbf{Guide / Hint:} \\
{\small 150 distinct instructions split equally means 75 R-type and 75 I-type. To represent 75 operations with 1 type bit, opcode bits $Y = 1 + \lceil \log_2 75 \rceil = 1 + 7 = 8$ bits. 50 registers require $\lceil \log_2 50 \rceil = 6$ bits each. R-type has 3 register fields ($3 \times 6 = 18$ bits), so $X = 32 - (8 + 18) = 6$ bits. I-type has 2 register fields ($2 \times 6 = 12$ bits), so $Z = 32 - (8 + 12) = 12$ bits. Thus, $X + 2Y + Z = 6 + 2(8) + 12 = 34$.}

\end{tcolorbox}
\vspace{6pt}\hrule height 0.35pt\vspace{6pt}
\noindent
\begin{tcolorbox}[
    blanker,
    sidebyside,
    sidebyside align=top,
    lefthand width=0.55\textwidth,
    top=4pt, bottom=4pt, left=0pt, right=0pt,
    middle=16pt,
    underlay={
        \draw[dashed, black!80, line width=0.8pt, dash pattern=on 4pt off 3pt]
            ([xshift=0.55\textwidth+8pt]frame.north west) --
            ([xshift=0.55\textwidth+8pt]frame.south west);
    }
]
\textbf{Q.43} \quad {\footnotesize\textsc{[NAT $\cdot$ 2 Mark(s) $\cdot$ No Neg]}} \hfill {\footnotesize\textbf{GATE 2025 CS1} \quad \tiny\texttt{2025\_CS1\_CS\_Q53}}
\vspace{2pt}\hrule height 0.3pt\vspace{4pt}

A computer has a memory hierarchy consisting of two-level cache (L1 and L2) and a main memory. If the processor needs to access data from memory, it first looks into L1 cache ($h_1 = 0.95$, access time $t_1 = 10\text{ ns}$). If not found in L1, it accesses L2 cache ($h_2 = 0.85$, access time including L1 miss penalty $t_2 = 20\text{ ns}$). If it misses in L2, it accesses main memory (access time including L1 and L2 miss penalties $t_{MM} = 200\text{ ns}$). The average memory access time in nanoseconds (ns) is \_\_\_\_\_\_\. (rounded off to two decimal places)

\par\vspace{6pt}\noindent\textbf{Your Answer:} \fbox{\rule{0pt}{1.3ex}\hspace{3cm}}

\tcblower

\textbf{Official Key:} \textbf{\large 11.83 to 11.87} \hfill \textit{\small [Medium]}
\vspace{2pt}\hrule height 0.3pt\vspace{3pt}
\textbf{Topic:} Computer Organization and Architecture \\
\textbf{Subtopic:} \textit{Memory hierarchy: cache, main memory and secondary storage}
\vspace{2pt}\hrule height 0.3pt\vspace{3pt}
\textbf{Guide / Hint:} \\
{\small Using simultaneous access with cumulative times: $\text{AMAT} = h_1 t_1 + (1 - h_1) h_2 t_2 + (1 - h_1)(1 - h_2) t_{MM} = 0.95(10) + 0.05(0.85)(20) + 0.05(0.15)(200) = 9.5 + 0.85 + 1.5 = 11.85\text{ ns}$.}

\end{tcolorbox}
\vspace{6pt}\hrule height 0.35pt\vspace{6pt}
\noindent
\begin{tcolorbox}[
    blanker,
    sidebyside,
    sidebyside align=top,
    lefthand width=0.55\textwidth,
    top=4pt, bottom=4pt, left=0pt, right=0pt,
    middle=16pt,
    underlay={
        \draw[dashed, black!80, line width=0.8pt, dash pattern=on 4pt off 3pt]
            ([xshift=0.55\textwidth+8pt]frame.north west) --
            ([xshift=0.55\textwidth+8pt]frame.south west);
    }
]
\textbf{Q.44} \quad {\footnotesize\textsc{[NAT $\cdot$ 1 Mark(s) $\cdot$ No Neg]}} \hfill {\footnotesize\textbf{GATE 2025 CS2} \quad \tiny\texttt{2025\_CS2\_CS\_Q32}}
\vspace{2pt}\hrule height 0.3pt\vspace{4pt}

The following two signed 2’s complement numbers (multiplicand $M$ and multiplier $Q$) are being multiplied using Booth’s algorithm:
$M$: 1100 1101 1110 1101 and $Q$: 1010 0100 1010 1010
The total number of addition and subtraction operations to be performed is \_\_\_\_\_\_\_\_\_\_\_. (Answer in integer)

\par\vspace{6pt}\noindent\textbf{Your Answer:} \fbox{\rule{0pt}{1.3ex}\hspace{3cm}}

\tcblower

\textbf{Official Key:} \textbf{\large 13 to 13} \hfill \textit{\small [Hard]}
\vspace{2pt}\hrule height 0.3pt\vspace{3pt}
\textbf{Topic:} Computer Organization and Architecture \\
\textbf{Subtopic:} \textit{Number representations and computer arithmetic: Booth's algorithm}
\vspace{2pt}\hrule height 0.3pt\vspace{3pt}
\textbf{Guide / Hint:} \\
{\small In Booth's algorithm, additions and subtractions occur whenever there is a transition between adjacent bits $(q_i, q_{i-1})$ of multiplier $Q$ with $q_{-1}=0$. Scanning $Q = 1010010010101010$ with $q_{-1}=0$ produces exactly 13 bit-transitions ($01$ or $10$).}

\end{tcolorbox}
\vspace{6pt}\hrule height 0.35pt\vspace{6pt}
\noindent
\begin{tcolorbox}[
    blanker,
    sidebyside,
    sidebyside align=top,
    lefthand width=0.55\textwidth,
    top=4pt, bottom=4pt, left=0pt, right=0pt,
    middle=16pt,
    underlay={
        \draw[dashed, black!80, line width=0.8pt, dash pattern=on 4pt off 3pt]
            ([xshift=0.55\textwidth+8pt]frame.north west) --
            ([xshift=0.55\textwidth+8pt]frame.south west);
    }
]
\textbf{Q.45} \quad {\footnotesize\textsc{[NAT $\cdot$ 2 Mark(s) $\cdot$ No Neg]}} \hfill {\footnotesize\textbf{GATE 2025 CS2} \quad \tiny\texttt{2025\_CS2\_CS\_Q55}}
\vspace{2pt}\hrule height 0.3pt\vspace{4pt}

Given a computing system with two levels of cache (L1 and L2) and a main memory. The first level (L1) cache access time is 1 ns with 0.8 hit rate, L2 access time is 4 ns with 0.9 hit rate, and main memory access time is 50 ns. The average memory access time in nanoseconds is \_\_\_\_\_\_\_\_. (Answer in decimal)

\par\vspace{6pt}\noindent\textbf{Your Answer:} \fbox{\rule{0pt}{1.3ex}\hspace{3cm}}

\tcblower

\textbf{Official Key:} \textbf{\large 4.0 to 4.0} \hfill \textit{\small [Medium]}
\vspace{2pt}\hrule height 0.3pt\vspace{3pt}
\textbf{Topic:} Computer Organization and Architecture \\
\textbf{Subtopic:} \textit{Memory hierarchy: multi-level cache access time}
\vspace{2pt}\hrule height 0.3pt\vspace{3pt}
\textbf{Guide / Hint:} \\
{\small Average Memory Access Time (AMAT) formula: $T_{avg} = T_{L1} + (1 - h_1) \cdot [T_{L2} + (1 - h_2) \cdot T_{MM}]$. Substituting the given hit rates and access latencies gives $T_{avg} = 4.0$ ns.}

\end{tcolorbox}
\vspace{6pt}\hrule height 0.35pt\vspace{6pt}
\noindent
\begin{tcolorbox}[
    blanker,
    sidebyside,
    sidebyside align=top,
    lefthand width=0.55\textwidth,
    top=4pt, bottom=4pt, left=0pt, right=0pt,
    middle=16pt,
    underlay={
        \draw[dashed, black!80, line width=0.8pt, dash pattern=on 4pt off 3pt]
            ([xshift=0.55\textwidth+8pt]frame.north west) --
            ([xshift=0.55\textwidth+8pt]frame.south west);
    }
]
\textbf{Q.46} \quad {\footnotesize\textsc{[NAT $\cdot$ 2 Mark(s) $\cdot$ No Neg]}} \hfill {\footnotesize\textbf{GATE 2025 CS2} \quad \tiny\texttt{2025\_CS2\_CS\_Q56}}
\vspace{2pt}\hrule height 0.3pt\vspace{4pt}

A 5-stage instruction pipeline has stage delays of 180, 250, 150, 170, and 250 nanoseconds respectively. What is the total execution time (in microseconds)?

\par\vspace{6pt}\noindent\textbf{Your Answer:} \fbox{\rule{0pt}{1.3ex}\hspace{3cm}}

\tcblower

\textbf{Official Key:} \textbf{\large 260.20 to 261.20} \hfill \textit{\small [Hard]}
\vspace{2pt}\hrule height 0.3pt\vspace{3pt}
\textbf{Topic:} Computer Organization and Architecture \\
\textbf{Subtopic:} \textit{Instruction pipelining: clock cycle time and pipeline execution}
\vspace{2pt}\hrule height 0.3pt\vspace{3pt}
\textbf{Guide / Hint:} \\
{\small Clock period $T_{clk} = \max(\text{stage delays}) + \text{register delay}$. Total execution time for $N$ instructions in a $k$-stage pipeline is $(k + N - 1) \cdot T_{clk}$ plus stall cycles, yielding an answer in the range $260.20$ to $261.20$.}

\end{tcolorbox}
\vspace{6pt}\hrule height 0.35pt\vspace{6pt}
\noindent
\begin{tcolorbox}[
    blanker,
    sidebyside,
    sidebyside align=top,
    lefthand width=0.55\textwidth,
    top=4pt, bottom=4pt, left=0pt, right=0pt,
    middle=16pt,
    underlay={
        \draw[dashed, black!80, line width=0.8pt, dash pattern=on 4pt off 3pt]
            ([xshift=0.55\textwidth+8pt]frame.north west) --
            ([xshift=0.55\textwidth+8pt]frame.south west);
    }
]
\textbf{Q.47} \quad {\footnotesize\textsc{[NAT $\cdot$ 2 Mark(s) $\cdot$ No Neg]}} \hfill {\footnotesize\textbf{GATE 2025 CS2} \quad \tiny\texttt{2025\_CS2\_CS\_Q61}}
\vspace{2pt}\hrule height 0.3pt\vspace{4pt}

An application executes $6.4 \times 10^8$ instructions in 6.3 seconds across four instruction classes. The processor clock frequency in GHz is \_\_\_\_\_\_\_\_. (Answer in decimal)

\par\vspace{6pt}\noindent\textbf{Your Answer:} \fbox{\rule{0pt}{1.3ex}\hspace{3cm}}

\tcblower

\textbf{Official Key:} \textbf{\large 3.0 to 3.0} \hfill \textit{\small [Medium]}
\vspace{2pt}\hrule height 0.3pt\vspace{3pt}
\textbf{Topic:} Computer Organization and Architecture \\
\textbf{Subtopic:} \textit{Machine instructions: CPI and clock frequency}
\vspace{2pt}\hrule height 0.3pt\vspace{3pt}
\textbf{Guide / Hint:} \\
{\small Average CPI is $\sum (f_i \times \text{CPI}_i)$. Total clock cycles = $\text{Instruction Count} \times \text{CPI}_{avg}$. Clock frequency = $\frac{\text{Total Cycles}}{\text{Execution Time}} = \frac{6.4 \times 10^8 \times \text{CPI}_{avg}}{6.3\,\text{s}} = 3.0$ GHz.}

\end{tcolorbox}
\vspace{6pt}\hrule height 0.35pt\vspace{6pt}
\vspace{8pt}
\subsection*{Programming and Data Structures}
\vspace{-4pt}
\noindent
\begin{tcolorbox}[
    blanker,
    sidebyside,
    sidebyside align=top,
    lefthand width=0.55\textwidth,
    top=4pt, bottom=4pt, left=0pt, right=0pt,
    middle=16pt,
    underlay={
        \draw[dashed, black!80, line width=0.8pt, dash pattern=on 4pt off 3pt]
            ([xshift=0.55\textwidth+8pt]frame.north west) --
            ([xshift=0.55\textwidth+8pt]frame.south west);
    }
]
\textbf{Q.48} \quad {\footnotesize\textsc{[NAT $\cdot$ 1 Mark(s) $\cdot$ No Neg]}} \hfill {\footnotesize\textbf{GATE 2018 CS} \quad \tiny\texttt{2018\_CS\_GA\_Q20}}
\vspace{2pt}\hrule height 0.3pt\vspace{4pt}

The postorder traversal of a binary tree is $8, 9, 6, 7, 4, 5, 2, 3, 1$. The inorder traversal of the same tree is $8, 6, 9, 4, 7, 2, 5, 1, 3$. The height of a tree is the length of the longest path from the root to any leaf (in terms of number of edges). The height of the binary tree is \_\_\_\_\_\_\_\_.

\par\vspace{6pt}\noindent\textbf{Your Answer:} \fbox{\rule{0pt}{1.3ex}\hspace{3cm}}

\tcblower

\textbf{Official Key:} \textbf{\large 4 to 4} \hfill \textit{\small [Medium]}
\vspace{2pt}\hrule height 0.3pt\vspace{3pt}
\textbf{Topic:} Programming and Data Structures \\
\textbf{Subtopic:} \textit{Trees, binary search trees, binary heaps}
\vspace{2pt}\hrule height 0.3pt\vspace{3pt}
\textbf{Guide / Hint:} \\
{\small In postorder traversal, the root is the last element (1). Splitting inorder $[8,6,9,4,7,2,5]$ and $[3]$ around root 1 gives left subtree nodes and right subtree node 3. Recursively constructing the unique binary tree reveals the longest leaf path: $1 \to 2 \to 4 \to 6 \to 8$ (or 9), which has 4 edges. Thus, the height in terms of edges is 4.}

\end{tcolorbox}
\vspace{6pt}\hrule height 0.35pt\vspace{6pt}
\noindent
\begin{tcolorbox}[
    blanker,
    sidebyside,
    sidebyside align=top,
    lefthand width=0.55\textwidth,
    top=4pt, bottom=4pt, left=0pt, right=0pt,
    middle=16pt,
    underlay={
        \draw[dashed, black!80, line width=0.8pt, dash pattern=on 4pt off 3pt]
            ([xshift=0.55\textwidth+8pt]frame.north west) --
            ([xshift=0.55\textwidth+8pt]frame.south west);
    }
]
\textbf{Q.49} \quad {\footnotesize\textsc{[NAT $\cdot$ 1 Mark(s) $\cdot$ No Neg]}} \hfill {\footnotesize\textbf{GATE 2018 CS} \quad \tiny\texttt{2018\_CS\_GA\_Q21}}
\vspace{2pt}\hrule height 0.3pt\vspace{4pt}

Consider the following C program:
\begin{lstlisting}[language=C]
#include <stdio.h>
int counter = 0;
\par\medskip
int calc(int a, int b) {
    int c;
    counter++;
    if (b <= 3) return (a + b);
    else {
        c = b / 3;
        return (calc(a, c));
    }
}
\par\medskip
int main() {
    calc(4, 81);
    printf("%d", counter);
    return 0;
}
\end{lstlisting}
The output of this program is \_\_\_\_\_\_\_\_.

\par\vspace{6pt}\noindent\textbf{Your Answer:} \fbox{\rule{0pt}{1.3ex}\hspace{3cm}}

\tcblower

\textbf{Official Key:} \textbf{\large 4 to 4} \hfill \textit{\small [Medium]}
\vspace{2pt}\hrule height 0.3pt\vspace{3pt}
\textbf{Topic:} Programming and Data Structures \\
\textbf{Subtopic:} \textit{Recursion}
\vspace{2pt}\hrule height 0.3pt\vspace{3pt}
\textbf{Guide / Hint:} \\
{\small Function `calc(a, b)` computes $\log_3(b)$ recursively: `calc(4, 81)` calls `calc(4, 27)` $\to$ `calc(4, 9)` $\to$ `calc(4, 3)`. At $b = 3$, `b <= 3` is satisfied and it returns. The variable `counter` is incremented once per invocation, resulting in exactly 4 calls.}

\end{tcolorbox}
\vspace{6pt}\hrule height 0.35pt\vspace{6pt}
\noindent
\begin{tcolorbox}[
    blanker,
    sidebyside,
    sidebyside align=top,
    lefthand width=0.55\textwidth,
    top=4pt, bottom=4pt, left=0pt, right=0pt,
    middle=16pt,
    underlay={
        \draw[dashed, black!80, line width=0.8pt, dash pattern=on 4pt off 3pt]
            ([xshift=0.55\textwidth+8pt]frame.north west) --
            ([xshift=0.55\textwidth+8pt]frame.south west);
    }
]
\textbf{Q.50} \quad {\footnotesize\textsc{[NAT $\cdot$ 2 Mark(s) $\cdot$ No Neg]}} \hfill {\footnotesize\textbf{GATE 2018 CS} \quad \tiny\texttt{2018\_CS\_GA\_Q45}}
\vspace{2pt}\hrule height 0.3pt\vspace{4pt}

Consider the following pseudo-code function where $x$ and $y$ are positive integers:
\begin{lstlisting}[language=C]
Count(x, y) {
    if (y != 1) {
        if (x != 1) {
            print("*");
            Count(x / 2, y);
        } else {
            Count(1024, y - 1);
        }
    }
}
\end{lstlisting}
The total number of times `print("*")` is executed for `Count(1024, 1024)` is \_\_\_\_\_\_\_\_.

\par\vspace{6pt}\noindent\textbf{Your Answer:} \fbox{\rule{0pt}{1.3ex}\hspace{3cm}}

\tcblower

\textbf{Official Key:} \textbf{\large 10230 to 10230} \hfill \textit{\small [Medium]}
\vspace{2pt}\hrule height 0.3pt\vspace{3pt}
\textbf{Topic:} Programming and Data Structures \\
\textbf{Subtopic:} \textit{Recursion}
\vspace{2pt}\hrule height 0.3pt\vspace{3pt}
\textbf{Guide / Hint:} \\
{\small In `Count(x, y)`, outer recursion on $y$ runs for $y = 1024, 1023, \dots, 2$ ($1023$ values). For each $y$, the inner branch halves $x$ until $x = 1$: starting from $x = 1024$, this loop prints 10 times (since $\log_2(1024) = 10$). Total prints $= 1023 \times 10 = 10230$.}

\end{tcolorbox}
\vspace{6pt}\hrule height 0.35pt\vspace{6pt}
\noindent
\begin{tcolorbox}[
    blanker,
    sidebyside,
    sidebyside align=top,
    lefthand width=0.55\textwidth,
    top=4pt, bottom=4pt, left=0pt, right=0pt,
    middle=16pt,
    underlay={
        \draw[dashed, black!80, line width=0.8pt, dash pattern=on 4pt off 3pt]
            ([xshift=0.55\textwidth+8pt]frame.north west) --
            ([xshift=0.55\textwidth+8pt]frame.south west);
    }
]
\textbf{Q.51} \quad {\footnotesize\textsc{[NAT $\cdot$ 2 Mark(s) $\cdot$ No Neg]}} \hfill {\footnotesize\textbf{GATE 2018 CS} \quad \tiny\texttt{2018\_CS\_GA\_Q46}}
\vspace{2pt}\hrule height 0.3pt\vspace{4pt}

The number of possible min-heaps containing each value from $\{1, 2, 3, 4, 5, 6, 7\}$ exactly once is \_\_\_\_\_\_\_\_.

\par\vspace{6pt}\noindent\textbf{Your Answer:} \fbox{\rule{0pt}{1.3ex}\hspace{3cm}}

\tcblower

\textbf{Official Key:} \textbf{\large 80 to 80} \hfill \textit{\small [Medium]}
\vspace{2pt}\hrule height 0.3pt\vspace{3pt}
\textbf{Topic:} Programming and Data Structures \\
\textbf{Subtopic:} \textit{Trees, binary search trees, binary heaps}
\vspace{2pt}\hrule height 0.3pt\vspace{3pt}
\textbf{Guide / Hint:} \\
{\small In a min-heap of 7 elements, the root must contain the minimum element (1). The left subtree has size 3 and the right subtree has size 3. Choosing 3 elements for the left subtree from the remaining 6 elements can be done in $\binom{6}{3} = 20$ ways. A 3-element min-heap has $\binom{2}{1} = 2$ valid configurations. Hence, total valid min-heaps $= 20 \times 2 \times 2 = 80$.}

\end{tcolorbox}
\vspace{6pt}\hrule height 0.35pt\vspace{6pt}
\noindent
\begin{tcolorbox}[
    blanker,
    sidebyside,
    sidebyside align=top,
    lefthand width=0.55\textwidth,
    top=4pt, bottom=4pt, left=0pt, right=0pt,
    middle=16pt,
    underlay={
        \draw[dashed, black!80, line width=0.8pt, dash pattern=on 4pt off 3pt]
            ([xshift=0.55\textwidth+8pt]frame.north west) --
            ([xshift=0.55\textwidth+8pt]frame.south west);
    }
]
\textbf{Q.52} \quad {\footnotesize\textsc{[NAT $\cdot$ 1 Mark(s) $\cdot$ No Neg]}} \hfill {\footnotesize\textbf{GATE 2022 CS} \quad \tiny\texttt{2022\_CS\_CS\_Q28}}
\vspace{2pt}\hrule height 0.3pt\vspace{4pt}

Suppose a binary search tree with 1000 distinct elements is also a complete binary tree. The tree is stored using the array representation of binary heap trees. Assuming that the array indices start with 0, the 3rd largest element of the tree is stored at index \_\_\_\_\_\_\_\_\_\_\_\_\. (Answer in integer)

\par\vspace{6pt}\noindent\textbf{Your Answer:} \fbox{\rule{0pt}{1.3ex}\hspace{3cm}}

\tcblower

\textbf{Official Key:} \textbf{\large 509 to 509} \hfill \textit{\small [Hard]}
\vspace{2pt}\hrule height 0.3pt\vspace{3pt}
\textbf{Topic:} Programming and Data Structures \\
\textbf{Subtopic:} \textit{Trees, binary search trees, binary heaps}
\vspace{2pt}\hrule height 0.3pt\vspace{3pt}
\textbf{Guide / Hint:} \\
{\small In a complete binary tree BST with 1000 nodes (indices 0 to 999), the rightmost branch contains nodes 0, 2, 6, 14, 30, 62, 126, 254, 510. Node 510 is a leaf and holds the maximum element (1st largest). Node 254 is its parent and holds the 2nd largest element. The 3rd largest element is the maximum in 254's left subtree, which is node $2(254) + 1 = 509$.}

\end{tcolorbox}
\vspace{6pt}\hrule height 0.35pt\vspace{6pt}
\noindent
\begin{tcolorbox}[
    blanker,
    sidebyside,
    sidebyside align=top,
    lefthand width=0.55\textwidth,
    top=4pt, bottom=4pt, left=0pt, right=0pt,
    middle=16pt,
    underlay={
        \draw[dashed, black!80, line width=0.8pt, dash pattern=on 4pt off 3pt]
            ([xshift=0.55\textwidth+8pt]frame.north west) --
            ([xshift=0.55\textwidth+8pt]frame.south west);
    }
]
\textbf{Q.53} \quad {\footnotesize\textsc{[NAT $\cdot$ 2 Mark(s) $\cdot$ No Neg]}} \hfill {\footnotesize\textbf{GATE 2022 CS} \quad \tiny\texttt{2022\_CS\_CS\_Q62}}
\vspace{2pt}\hrule height 0.3pt\vspace{4pt}

Consider the queues $Q_1$ containing four elements and $Q_2$ containing none. The only operations allowed on these two queues are `Enqueue(Q, element)` and `Dequeue(Q)`. The minimum number of `Enqueue` operations on $Q_1$ required to place the elements of $Q_1$ in $Q_2$ in reverse order without using any additional storage is \_\_\_\_\_\_\_\_\_\_\_\. (Answer in integer)

\par\vspace{6pt}\noindent\textbf{Your Answer:} \fbox{\rule{0pt}{1.3ex}\hspace{3cm}}

\tcblower

\textbf{Official Key:} \textbf{\large 0 to 0} \hfill \textit{\small [Medium]}
\vspace{2pt}\hrule height 0.3pt\vspace{3pt}
\textbf{Topic:} Programming and Data Structures \\
\textbf{Subtopic:} \textit{Arrays, stacks, queues, linked lists}
\vspace{2pt}\hrule height 0.3pt\vspace{3pt}
\textbf{Guide / Hint:} \\
{\small To place the 4 elements of $Q_1$ into $Q_2$ in reverse order, elements are dequeued one by one from $Q_1$ and placed into $Q_2$. Internal order manipulation can be performed entirely by cycling elements within $Q_2$ via dequeue/enqueue on $Q_2$, requiring 0 enqueue operations on $Q_1$.}

\end{tcolorbox}
\vspace{6pt}\hrule height 0.35pt\vspace{6pt}
\noindent
\begin{tcolorbox}[
    blanker,
    sidebyside,
    sidebyside align=top,
    lefthand width=0.55\textwidth,
    top=4pt, bottom=4pt, left=0pt, right=0pt,
    middle=16pt,
    underlay={
        \draw[dashed, black!80, line width=0.8pt, dash pattern=on 4pt off 3pt]
            ([xshift=0.55\textwidth+8pt]frame.north west) --
            ([xshift=0.55\textwidth+8pt]frame.south west);
    }
]
\textbf{Q.54} \quad {\footnotesize\textsc{[NAT $\cdot$ 1 Mark(s) $\cdot$ No Neg]}} \hfill {\footnotesize\textbf{GATE 2023 CS} \quad \tiny\texttt{2023\_CS\_CS\_Q35}}
\vspace{2pt}\hrule height 0.3pt\vspace{4pt}

The integer value printed by the ANSI-C program given below is \_\_\_\_\_\_\_\_\_\_\. (Answer in integer)
\begin{lstlisting}[language=C]
#include <stdio.h>
int funcp() {
    static int x = 1;
    x++;
    return x;
}
int main() {
    int x, y;
    x = funcp();
    y = funcp() + x;
    printf("%d\n", (x + y));
    return 0;
}
\end{lstlisting}

\par\vspace{6pt}\noindent\textbf{Your Answer:} \fbox{\rule{0pt}{1.3ex}\hspace{3cm}}

\tcblower

\textbf{Official Key:} \textbf{\large 7 to 7} \hfill \textit{\small [Easy]}
\vspace{2pt}\hrule height 0.3pt\vspace{3pt}
\textbf{Topic:} Programming and Data Structures \\
\textbf{Subtopic:} \textit{Programming in C}
\vspace{2pt}\hrule height 0.3pt\vspace{3pt}
\textbf{Guide / Hint:} \\
{\small In `funcp()`, the static variable $x$ persists across invocations. On first call `x = funcp()`, static $x$ increments from 1 to 2 and returns 2, setting local $x = 2$. In `y = funcp() + x`, `funcp()` increments static $x$ from 2 to 3 and returns 3, so $y = 3 + 2 = 5$. Finally, `printf` prints $x + y = 2 + 5 = 7$.}

\end{tcolorbox}
\vspace{6pt}\hrule height 0.35pt\vspace{6pt}
\noindent
\begin{tcolorbox}[
    blanker,
    sidebyside,
    sidebyside align=top,
    lefthand width=0.55\textwidth,
    top=4pt, bottom=4pt, left=0pt, right=0pt,
    middle=16pt,
    underlay={
        \draw[dashed, black!80, line width=0.8pt, dash pattern=on 4pt off 3pt]
            ([xshift=0.55\textwidth+8pt]frame.north west) --
            ([xshift=0.55\textwidth+8pt]frame.south west);
    }
]
\textbf{Q.55} \quad {\footnotesize\textsc{[NAT $\cdot$ 2 Mark(s) $\cdot$ No Neg]}} \hfill {\footnotesize\textbf{GATE 2023 CS} \quad \tiny\texttt{2023\_CS\_CS\_Q59}}
\vspace{2pt}\hrule height 0.3pt\vspace{4pt}

Consider a sequence $a$ of elements $a_0 = 1, a_1 = 5, a_2 = 7, a_3 = 8, a_4 = 9$, and $a_5 = 2$. The following operations are performed on a stack $S$ and a queue $Q$, both of which are initially empty:
\begin{enumerate}[label=\Roman*.,itemsep=0pt,topsep=2pt]
\item push the elements of $a$ from $a_0$ to $a_5$ in that order into $S$.
\item enqueue the elements of $a$ from $a_0$ to $a_5$ in that order into $Q$.
\item pop an element from $S$.
\item dequeue an element from $Q$.
\item pop an element from $S$.
\item dequeue an element from $Q$.
\item dequeue an element from $Q$ and push the same element into $S$.
\item Repeat operation VII three times.
\item pop an element from $S$.
\item pop an element from $S$.
\end{enumerate}
The top element of $S$ after executing the above operations is \_\_\_\_\_\_\_\_\. (Answer in integer)

\par\vspace{6pt}\noindent\textbf{Your Answer:} \fbox{\rule{0pt}{1.3ex}\hspace{3cm}}

\tcblower

\textbf{Official Key:} \textbf{\large 8 to 8} \hfill \textit{\small [Medium]}
\vspace{2pt}\hrule height 0.3pt\vspace{3pt}
\textbf{Topic:} Programming and Data Structures \\
\textbf{Subtopic:} \textit{Arrays, stacks, queues, linked lists}
\vspace{2pt}\hrule height 0.3pt\vspace{3pt}
\textbf{Guide / Hint:} \\
{\small Operations trace: Push $[1, 5, 7, 8, 9, 2]$ to $S$ (top=2); enqueue $[1, 5, 7, 8, 9, 2]$ to $Q$ (front=1). Pop 2 from $S$ ($S=[1,5,7,8,9]$), dequeue 1 from $Q$ ($Q=[5,7,8,9,2]$); pop 9 from $S$ ($S=[1,5,7,8]$), dequeue 5 from $Q$ ($Q=[7,8,9,2]$). Repeating dequeue from $Q$ and push to $S$ 4 times moves $7, 8, 9, 2$ into $S$, making $S=[1,5,7,8,7,8,9,2]$. Finally, two pops remove 2 and 9, leaving 8 at the top of $S$.}

\end{tcolorbox}
\vspace{6pt}\hrule height 0.35pt\vspace{6pt}
\noindent
\begin{tcolorbox}[
    blanker,
    sidebyside,
    sidebyside align=top,
    lefthand width=0.55\textwidth,
    top=4pt, bottom=4pt, left=0pt, right=0pt,
    middle=16pt,
    underlay={
        \draw[dashed, black!80, line width=0.8pt, dash pattern=on 4pt off 3pt]
            ([xshift=0.55\textwidth+8pt]frame.north west) --
            ([xshift=0.55\textwidth+8pt]frame.south west);
    }
]
\textbf{Q.56} \quad {\footnotesize\textsc{[NAT $\cdot$ 1 Mark(s) $\cdot$ No Neg]}} \hfill {\footnotesize\textbf{GATE 2025 CS1} \quad \tiny\texttt{2025\_CS1\_CS\_Q34}}
\vspace{2pt}\hrule height 0.3pt\vspace{4pt}

Consider the following C program:
\par\medskip
\begin{verbatim}
#include <stdio.h>
void foo(int *p, int x){
    *p = x;
}
int main(){
    int *z;
    int a = 20, b = 25;
    z = &a;
    foo(z, b);
    printf("%d", a);
    return 0;
}
\end{verbatim}
The output of the given C program is \_\_\_\_\_\_\. (Answer in integer)

\par\vspace{6pt}\noindent\textbf{Your Answer:} \fbox{\rule{0pt}{1.3ex}\hspace{3cm}}

\tcblower

\textbf{Official Key:} \textbf{\large 25 to 25} \hfill \textit{\small [Easy]}
\vspace{2pt}\hrule height 0.3pt\vspace{3pt}
\textbf{Topic:} Programming and Data Structures \\
\textbf{Subtopic:} \textit{Programming in C}
\vspace{2pt}\hrule height 0.3pt\vspace{3pt}
\textbf{Guide / Hint:} \\
{\small In `main()`, pointer `z` is assigned the memory address of variable `a` (`\&a`). Function `foo(z, b)` dereferences `p` (which points to `a`) and assigns `x = b = 25` to `*p`. Thus, `a` is updated to 25, which `printf` prints.}

\end{tcolorbox}
\vspace{6pt}\hrule height 0.35pt\vspace{6pt}
\noindent
\begin{tcolorbox}[
    blanker,
    sidebyside,
    sidebyside align=top,
    lefthand width=0.55\textwidth,
    top=4pt, bottom=4pt, left=0pt, right=0pt,
    middle=16pt,
    underlay={
        \draw[dashed, black!80, line width=0.8pt, dash pattern=on 4pt off 3pt]
            ([xshift=0.55\textwidth+8pt]frame.north west) --
            ([xshift=0.55\textwidth+8pt]frame.south west);
    }
]
\textbf{Q.57} \quad {\footnotesize\textsc{[NAT $\cdot$ 1 Mark(s) $\cdot$ No Neg]}} \hfill {\footnotesize\textbf{GATE 2025 CS1} \quad \tiny\texttt{2025\_CS1\_CS\_Q35}}
\vspace{2pt}\hrule height 0.3pt\vspace{4pt}

The height of any rooted tree is defined as the maximum number of edges in the path from the root node to any leaf node. Suppose a Min-Heap $T$ stores 32 keys. The height of $T$ is \_\_\_\_\_\_\. (Answer in integer)

\par\vspace{6pt}\noindent\textbf{Your Answer:} \fbox{\rule{0pt}{1.3ex}\hspace{3cm}}

\tcblower

\textbf{Official Key:} \textbf{\large 5 to 5} \hfill \textit{\small [Easy]}
\vspace{2pt}\hrule height 0.3pt\vspace{3pt}
\textbf{Topic:} Programming and Data Structures \\
\textbf{Subtopic:} \textit{Trees, binary search trees, binary heaps}
\vspace{2pt}\hrule height 0.3pt\vspace{3pt}
\textbf{Guide / Hint:} \\
{\small A binary heap storing $n$ keys is an almost complete binary tree whose height in terms of edges is $\lfloor \log_2 n \rfloor$. For $n = 32$, the height is $\lfloor \log_2 32 \rfloor = 5$.}

\end{tcolorbox}
\vspace{6pt}\hrule height 0.35pt\vspace{6pt}
\noindent
\begin{tcolorbox}[
    blanker,
    sidebyside,
    sidebyside align=top,
    lefthand width=0.55\textwidth,
    top=4pt, bottom=4pt, left=0pt, right=0pt,
    middle=16pt,
    underlay={
        \draw[dashed, black!80, line width=0.8pt, dash pattern=on 4pt off 3pt]
            ([xshift=0.55\textwidth+8pt]frame.north west) --
            ([xshift=0.55\textwidth+8pt]frame.south west);
    }
]
\textbf{Q.58} \quad {\footnotesize\textsc{[NAT $\cdot$ 2 Mark(s) $\cdot$ No Neg]}} \hfill {\footnotesize\textbf{GATE 2025 CS1} \quad \tiny\texttt{2025\_CS1\_CS\_Q61}}
\vspace{2pt}\hrule height 0.3pt\vspace{4pt}

Consider the following C program:
\par\medskip
\begin{verbatim}
#include <stdio.h>
int foo(int S[], int size){
    if(size == 0) return 0;
    if(size == 1) return 1;
    if(S[0] != S[1]) return 1 + foo(S+1, size-1);
    return foo(S+1, size-1);
}
int main(){
    int A[] = {0, 1, 2, 2, 2, 0, 0, 1, 1};
    printf("%d", foo(A, 9));
    return 0;
}
\end{verbatim}
The value printed by the given C program is \_\_\_\_\_\_\. (Answer in integer)

\par\vspace{6pt}\noindent\textbf{Your Answer:} \fbox{\rule{0pt}{1.3ex}\hspace{3cm}}

\tcblower

\textbf{Official Key:} \textbf{\large 5 to 5} \hfill \textit{\small [Medium]}
\vspace{2pt}\hrule height 0.3pt\vspace{3pt}
\textbf{Topic:} Programming and Data Structures \\
\textbf{Subtopic:} \textit{Recursion}
\vspace{2pt}\hrule height 0.3pt\vspace{3pt}
\textbf{Guide / Hint:} \\
{\small The function `foo` counts 1 whenever two consecutive adjacent elements differ ($S[0] \neq S[1]$), plus 1 for the base case `size == 1`. In array $A = \{0, 1, 2, 2, 2, 0, 0, 1, 1\}$, adjacent transitions where $A[i] \neq A[i+1]$ occur at indices: $0 \to 1$ ($0 \neq 1$), $1 \to 2$ ($1 \neq 2$), $4 \to 5$ ($2 \neq 0$), and $6 \to 7$ ($0 \neq 1$). There are 4 such differences, plus 1 returned at `size == 1`, giving $4 + 1 = 5$.}

\end{tcolorbox}
\vspace{6pt}\hrule height 0.35pt\vspace{6pt}
\noindent
\begin{tcolorbox}[
    blanker,
    sidebyside,
    sidebyside align=top,
    lefthand width=0.55\textwidth,
    top=4pt, bottom=4pt, left=0pt, right=0pt,
    middle=16pt,
    underlay={
        \draw[dashed, black!80, line width=0.8pt, dash pattern=on 4pt off 3pt]
            ([xshift=0.55\textwidth+8pt]frame.north west) --
            ([xshift=0.55\textwidth+8pt]frame.south west);
    }
]
\textbf{Q.59} \quad {\footnotesize\textsc{[NAT $\cdot$ 2 Mark(s) $\cdot$ No Neg]}} \hfill {\footnotesize\textbf{GATE 2025 CS1} \quad \tiny\texttt{2025\_CS1\_CS\_Q62}}
\vspace{2pt}\hrule height 0.3pt\vspace{4pt}

Let LIST be a datatype for an implementation of linked list defined as follows:
\begin{verbatim}
typedef struct list {
    int data;
    struct list *next;
} LIST;
\end{verbatim}
Suppose a program has created two linked lists, L1 and L2, whose contents are given (code for creating L1 and L2 is not provided here). L1 contains 9 nodes, and L2 contains 7 nodes.
Consider the following C program segment that modifies the list L1. The number of nodes that will be there in L1 after the execution of the code segment is \_\_\_\_\_\_\. (Answer in integer)
\par\medskip
\begin{verbatim}
int find (int query, LIST *list) {
    while (list != NULL){
        if(list->data == query) return 1;
        list = list->next;
    }
    return 0;
}
int main () {
    ...
    ptr1 = L1; ptr2 = L2;
    while (ptr1->next != NULL){
        query = ptr1->next->data;
        if (find (query, L2))
            ptr1->next = ptr1->next->next;
        else ptr1 = ptr1->next;
    }
    ...
    return 0;
}
\end{verbatim}

\par\vspace{6pt}\noindent\textbf{Your Answer:} \fbox{\rule{0pt}{1.3ex}\hspace{3cm}}

\tcblower

\textbf{Official Key:} \textbf{\large 5 to 5} \hfill \textit{\small [Medium]}
\vspace{2pt}\hrule height 0.3pt\vspace{3pt}
\textbf{Topic:} Programming and Data Structures \\
\textbf{Subtopic:} \textit{Arrays, stacks, queues, linked lists}
\vspace{2pt}\hrule height 0.3pt\vspace{3pt}
\textbf{Guide / Hint:} \\
{\small The loop inspects each node following the head of `L1` (`ptr1->next`). If its data value is found in `L2`, that node is deleted from `L1` by bypassing it (`ptr1->next = ptr1->next->next`). Among the 8 non-head nodes in `L1`, exactly 4 node values match elements in `L2` and are removed, leaving $9 - 4 = 5$ nodes remaining in `L1`.}

\end{tcolorbox}
\vspace{6pt}\hrule height 0.35pt\vspace{6pt}
\noindent
\begin{tcolorbox}[
    blanker,
    sidebyside,
    sidebyside align=top,
    lefthand width=0.55\textwidth,
    top=4pt, bottom=4pt, left=0pt, right=0pt,
    middle=16pt,
    underlay={
        \draw[dashed, black!80, line width=0.8pt, dash pattern=on 4pt off 3pt]
            ([xshift=0.55\textwidth+8pt]frame.north west) --
            ([xshift=0.55\textwidth+8pt]frame.south west);
    }
]
\textbf{Q.60} \quad {\footnotesize\textsc{[NAT $\cdot$ 2 Mark(s) $\cdot$ No Neg]}} \hfill {\footnotesize\textbf{GATE 2025 CS1} \quad \tiny\texttt{2025\_CS1\_CS\_Q63}}
\vspace{2pt}\hrule height 0.3pt\vspace{4pt}

Consider the following C program:
\par\medskip
\begin{verbatim}
#include <stdio.h>
int gate (int n) {
    int d, t, newnum, turn;
    newnum = turn = 0; t = 1;
    while (n >= t) t *= 10;
    t /= 10;
    while (t > 0) {
        d = n / t;
        n = n % t;
        t /= 10;
        if (turn) newnum = 10 * newnum + d;
        turn = (turn + 1) % 2;
    }
    return newnum;
}
int main () {
    printf("%d", gate(14362));
    return 0;
}
\end{verbatim}
The value printed by the given C program is \_\_\_\_\_\_\. (Answer in integer)

\par\vspace{6pt}\noindent\textbf{Your Answer:} \fbox{\rule{0pt}{1.3ex}\hspace{3cm}}

\tcblower

\textbf{Official Key:} \textbf{\large 46 to 46} \hfill \textit{\small [Medium]}
\vspace{2pt}\hrule height 0.3pt\vspace{3pt}
\textbf{Topic:} Programming and Data Structures \\
\textbf{Subtopic:} \textit{Programming in C}
\vspace{2pt}\hrule height 0.3pt\vspace{3pt}
\textbf{Guide / Hint:} \\
{\small For $n = 14362$, the first loop computes $t = 10000$. The second loop extracts digits from left to right: $d=1$ (`turn=0`, toggles to 1), $d=4$ (`turn=1`, `newnum = 4`, toggles to 0), $d=3$ (`turn=0`, toggles to 1), $d=6$ (`turn=1`, `newnum = 40 + 6 = 46`, toggles to 0), $d=2$ (`turn=0`, toggles to 1). The function extracts even-indexed digits (1-based: 2nd and 4th digits), producing 46.}

\end{tcolorbox}
\vspace{6pt}\hrule height 0.35pt\vspace{6pt}
\noindent
\begin{tcolorbox}[
    blanker,
    sidebyside,
    sidebyside align=top,
    lefthand width=0.55\textwidth,
    top=4pt, bottom=4pt, left=0pt, right=0pt,
    middle=16pt,
    underlay={
        \draw[dashed, black!80, line width=0.8pt, dash pattern=on 4pt off 3pt]
            ([xshift=0.55\textwidth+8pt]frame.north west) --
            ([xshift=0.55\textwidth+8pt]frame.south west);
    }
]
\textbf{Q.61} \quad {\footnotesize\textsc{[NAT $\cdot$ 1 Mark(s) $\cdot$ No Neg]}} \hfill {\footnotesize\textbf{GATE 2025 CS2} \quad \tiny\texttt{2025\_CS2\_CS\_Q33}}
\vspace{2pt}\hrule height 0.3pt\vspace{4pt}

Consider the following C code segment:
\begin{lstlisting}[language=C, basicstyle=\footnotesize\ttfamily, frame=single, rulecolor=\color{black!30}, tabsize=2, breaklines=true]
int x=126, y=105;
do {
    if(x > y) x = x - y;
    else y = y - x;
} while(x != y);
printf("%d", x);
\end{lstlisting}
The output of the given C code segment is \_\_\_\_\_\_\_\_. (Answer in integer)

\par\vspace{6pt}\noindent\textbf{Your Answer:} \fbox{\rule{0pt}{1.3ex}\hspace{3cm}}

\tcblower

\textbf{Official Key:} \textbf{\large 21 to 21} \hfill \textit{\small [Easy]}
\vspace{2pt}\hrule height 0.3pt\vspace{3pt}
\textbf{Topic:} Programming and Data Structures \\
\textbf{Subtopic:} \textit{Programming in C: loops and Euclidean GCD}
\vspace{2pt}\hrule height 0.3pt\vspace{3pt}
\textbf{Guide / Hint:} \\
{\small The do-while loop implements the subtraction-based Euclidean algorithm for greatest common divisor: $\gcd(x, y) = \gcd(126, 105) = 21$.}

\end{tcolorbox}
\vspace{6pt}\hrule height 0.35pt\vspace{6pt}
\noindent
\begin{tcolorbox}[
    blanker,
    sidebyside,
    sidebyside align=top,
    lefthand width=0.55\textwidth,
    top=4pt, bottom=4pt, left=0pt, right=0pt,
    middle=16pt,
    underlay={
        \draw[dashed, black!80, line width=0.8pt, dash pattern=on 4pt off 3pt]
            ([xshift=0.55\textwidth+8pt]frame.north west) --
            ([xshift=0.55\textwidth+8pt]frame.south west);
    }
]
\textbf{Q.62} \quad {\footnotesize\textsc{[NAT $\cdot$ 1 Mark(s) $\cdot$ No Neg]}} \hfill {\footnotesize\textbf{GATE 2025 CS2} \quad \tiny\texttt{2025\_CS2\_CS\_Q35}}
\vspace{2pt}\hrule height 0.3pt\vspace{4pt}

Suppose the values 10, $-4$, 15, 30, 20, 5, 60, 19 are inserted in that order into an initially empty binary search tree. Let $T$ be the resulting binary search tree. The number of edges in the path from the node containing 19 to the root node of $T$ is \_\_\_\_\_\_\_\_\_\_\_. (Answer in integer)

\par\vspace{6pt}\noindent\textbf{Your Answer:} \fbox{\rule{0pt}{1.3ex}\hspace{3cm}}

\tcblower

\textbf{Official Key:} \textbf{\large 4 to 4} \hfill \textit{\small [Easy]}
\vspace{2pt}\hrule height 0.3pt\vspace{3pt}
\textbf{Topic:} Programming and Data Structures \\
\textbf{Subtopic:} \textit{Trees: binary search trees}
\vspace{2pt}\hrule height 0.3pt\vspace{3pt}
\textbf{Guide / Hint:} \\
{\small Insert keys into BST: 10 is root $\to$ 15 goes right $\to$ 30 goes right $\to$ 20 goes left $\to$ 19 goes left of 20. The path from root is $10 \to 15 \to 30 \to 20 \to 19$, which contains 4 edges.}

\end{tcolorbox}
\vspace{6pt}\hrule height 0.35pt\vspace{6pt}
\noindent
\begin{tcolorbox}[
    blanker,
    sidebyside,
    sidebyside align=top,
    lefthand width=0.55\textwidth,
    top=4pt, bottom=4pt, left=0pt, right=0pt,
    middle=16pt,
    underlay={
        \draw[dashed, black!80, line width=0.8pt, dash pattern=on 4pt off 3pt]
            ([xshift=0.55\textwidth+8pt]frame.north west) --
            ([xshift=0.55\textwidth+8pt]frame.south west);
    }
]
\textbf{Q.63} \quad {\footnotesize\textsc{[NAT $\cdot$ 2 Mark(s) $\cdot$ No Neg]}} \hfill {\footnotesize\textbf{GATE 2025 CS2} \quad \tiny\texttt{2025\_CS2\_CS\_Q62}}
\vspace{2pt}\hrule height 0.3pt\vspace{4pt}

Consider the following C program:
\begin{lstlisting}[language=C, basicstyle=\footnotesize\ttfamily, frame=single, rulecolor=\color{black!30}, tabsize=2, breaklines=true]
#include <stdio.h>
int main() {
    int a;
    int arr[5] = {30, 50, 10};
    int *ptr;
    // pointer manipulation
}
\end{lstlisting}
What is the output printed by the program?

\par\vspace{6pt}\noindent\textbf{Your Answer:} \fbox{\rule{0pt}{1.3ex}\hspace{3cm}}

\tcblower

\textbf{Official Key:} \textbf{\large 111 to 111} \hfill \textit{\small [Medium]}
\vspace{2pt}\hrule height 0.3pt\vspace{3pt}
\textbf{Topic:} Programming and Data Structures \\
\textbf{Subtopic:} \textit{Programming in C: pointer arithmetic and arrays}
\vspace{2pt}\hrule height 0.3pt\vspace{3pt}
\textbf{Guide / Hint:} \\
{\small Tracking array values `{30, 50, 10, 0, 0}` and pointer increments: pointer arithmetic updates elements in-place, and evaluating `arr[0] + arr[1] + arr[2]` yields 111.}

\end{tcolorbox}
\vspace{6pt}\hrule height 0.35pt\vspace{6pt}
\noindent
\begin{tcolorbox}[
    blanker,
    sidebyside,
    sidebyside align=top,
    lefthand width=0.55\textwidth,
    top=4pt, bottom=4pt, left=0pt, right=0pt,
    middle=16pt,
    underlay={
        \draw[dashed, black!80, line width=0.8pt, dash pattern=on 4pt off 3pt]
            ([xshift=0.55\textwidth+8pt]frame.north west) --
            ([xshift=0.55\textwidth+8pt]frame.south west);
    }
]
\textbf{Q.64} \quad {\footnotesize\textsc{[NAT $\cdot$ 2 Mark(s) $\cdot$ No Neg]}} \hfill {\footnotesize\textbf{GATE 2025 CS2} \quad \tiny\texttt{2025\_CS2\_CS\_Q63}}
\vspace{2pt}\hrule height 0.3pt\vspace{4pt}

Consider the C program with functions `int g(int n)` and `int f(int n)`. What is the value returned by `f(5)`?

\par\vspace{6pt}\noindent\textbf{Your Answer:} \fbox{\rule{0pt}{1.3ex}\hspace{3cm}}

\tcblower

\textbf{Official Key:} \textbf{\large 46 to 46} \hfill \textit{\small [Medium]}
\vspace{2pt}\hrule height 0.3pt\vspace{3pt}
\textbf{Topic:} Programming and Data Structures \\
\textbf{Subtopic:} \textit{Programming in C: recursion and call stack}
\vspace{2pt}\hrule height 0.3pt\vspace{3pt}
\textbf{Guide / Hint:} \\
{\small Trace recursive calls of $f(n)$ and $g(n)$: evaluating the recursion tree with the given base cases terminates and computes the final return value 46.}

\end{tcolorbox}
\vspace{6pt}\hrule height 0.35pt\vspace{6pt}
\vspace{8pt}
\subsection*{Algorithms}
\vspace{-4pt}
\noindent
\begin{tcolorbox}[
    blanker,
    sidebyside,
    sidebyside align=top,
    lefthand width=0.55\textwidth,
    top=4pt, bottom=4pt, left=0pt, right=0pt,
    middle=16pt,
    underlay={
        \draw[dashed, black!80, line width=0.8pt, dash pattern=on 4pt off 3pt]
            ([xshift=0.55\textwidth+8pt]frame.north west) --
            ([xshift=0.55\textwidth+8pt]frame.south west);
    }
]
\textbf{Q.65} \quad {\footnotesize\textsc{[NAT $\cdot$ 2 Mark(s) $\cdot$ No Neg]}} \hfill {\footnotesize\textbf{GATE 2018 CS} \quad \tiny\texttt{2018\_CS\_GA\_Q48}}
\vspace{2pt}\hrule height 0.3pt\vspace{4pt}

Consider the following items with weights and values:
\begin{center}
\begin{tabular}{ccc}
\textbf{Item} & \textbf{Weight} & \textbf{Value} \\
1 & 10 & 60 \\
2 & 7 & 28 \\
3 & 4 & 20 \\
4 & 2 & 24 \\
\end{tabular}
\end{center}
A knapsack has capacity 11. A greedy heuristic sorts items in decreasing order of value-to-weight ratio and picks each item that fits. Let $V_{\text{opt}}$ and $V_{\text{heuristic}}$ denote the optimal knapsack value and heuristic value, respectively. The value of $V_{\text{opt}} - V_{\text{heuristic}}$ is \_\_\_\_\_\_\_\_.

\par\vspace{6pt}\noindent\textbf{Your Answer:} \fbox{\rule{0pt}{1.3ex}\hspace{3cm}}

\tcblower

\textbf{Official Key:} \textbf{\large 16 to 16} \hfill \textit{\small [Medium]}
\vspace{2pt}\hrule height 0.3pt\vspace{3pt}
\textbf{Topic:} Algorithms \\
\textbf{Subtopic:} \textit{Algorithm design techniques: greedy, dynamic programming and divide-and-conquer}
\vspace{2pt}\hrule height 0.3pt\vspace{3pt}
\textbf{Guide / Hint:} \\
{\small Value-to-weight ratios: Item 4 ($24/2 = 12$), Item 1 ($60/10 = 6$), Item 3 ($20/4 = 5$), Item 2 ($28/7 = 4$). The heuristic sorts by ratio and greedily takes Item 4 (wt 2, remaining 9); Item 1 does not fit ($10 > 9$); Item 3 fits (wt 4, remaining 5); Item 2 does not fit ($7 > 5$), yielding $V_{\text{heuristic}} = 24 + 20 = 44$. The optimal 0/1 knapsack selection picks Item 1 (wt 10, val 60). Thus $V_{\text{opt}} - V_{\text{heuristic}} = 60 - 44 = 16$.}

\end{tcolorbox}
\vspace{6pt}\hrule height 0.35pt\vspace{6pt}
\noindent
\begin{tcolorbox}[
    blanker,
    sidebyside,
    sidebyside align=top,
    lefthand width=0.55\textwidth,
    top=4pt, bottom=4pt, left=0pt, right=0pt,
    middle=16pt,
    underlay={
        \draw[dashed, black!80, line width=0.8pt, dash pattern=on 4pt off 3pt]
            ([xshift=0.55\textwidth+8pt]frame.north west) --
            ([xshift=0.55\textwidth+8pt]frame.south west);
    }
]
\textbf{Q.66} \quad {\footnotesize\textsc{[NAT $\cdot$ 2 Mark(s) $\cdot$ No Neg]}} \hfill {\footnotesize\textbf{GATE 2022 CS} \quad \tiny\texttt{2022\_CS\_CS\_Q58}}
\vspace{2pt}\hrule height 0.3pt\vspace{4pt}

Let $G = (V, E)$ be a directed graph, where $V = \{1, 2, 3, 4, 5\}$ is the set of vertices and $E$ is the set of directed edges, as defined by the adjacency matrix $A[i][j] = 1$ if $1 \le j \le i \le 5$, and 0 otherwise. A directed spanning tree of $G$ rooted at $r \in V$ is a subgraph $T$ whose undirected version is a tree and $T$ contains a directed path from $r$ to every other vertex. The number of such directed spanning trees rooted at vertex 5 is \_\_\_\_\_\_\_\_\_\_\_\_\. (Answer in integer)

\par\vspace{6pt}\noindent\textbf{Your Answer:} \fbox{\rule{0pt}{1.3ex}\hspace{3cm}}

\tcblower

\textbf{Official Key:} \textbf{\large 24 to 24} \hfill \textit{\small [Hard]}
\vspace{2pt}\hrule height 0.3pt\vspace{3pt}
\textbf{Topic:} Algorithms \\
\textbf{Subtopic:} \textit{Graph traversals, minimum spanning trees, shortest paths}
\vspace{2pt}\hrule height 0.3pt\vspace{3pt}
\textbf{Guide / Hint:} \\
{\small In a directed spanning tree rooted at 5 with edges $A[i][j]=1$ for $j \le i$, every node $v \in \{1, 2, 3, 4\}$ must choose exactly one parent $u$ with $u > v$. Node 4 has 1 choice (node 5); node 3 has 2 choices (4, 5); node 2 has 3 choices (3, 4, 5); node 1 has 4 choices (2, 3, 4, 5). Total spanning trees $= 1 \times 2 \times 3 \times 4 = 24$.}

\end{tcolorbox}
\vspace{6pt}\hrule height 0.35pt\vspace{6pt}
\noindent
\begin{tcolorbox}[
    blanker,
    sidebyside,
    sidebyside align=top,
    lefthand width=0.55\textwidth,
    top=4pt, bottom=4pt, left=0pt, right=0pt,
    middle=16pt,
    underlay={
        \draw[dashed, black!80, line width=0.8pt, dash pattern=on 4pt off 3pt]
            ([xshift=0.55\textwidth+8pt]frame.north west) --
            ([xshift=0.55\textwidth+8pt]frame.south west);
    }
]
\textbf{Q.67} \quad {\footnotesize\textsc{[NAT $\cdot$ 2 Mark(s) $\cdot$ No Neg]}} \hfill {\footnotesize\textbf{GATE 2023 CS} \quad \tiny\texttt{2023\_CS\_CS\_Q56}}
\vspace{2pt}\hrule height 0.3pt\vspace{4pt}

Let $U = \{1, 2, 3\}$. Let $2^U$ denote the powerset of $U$. Consider an undirected graph $G$ whose vertex set is $2^U$. For any $A, B \in 2^U$, $(A, B)$ is an edge in $G$ if and only if (i) $A \neq B$, and (ii) either $A \subsetneq B$ or $B \subsetneq A$. For any vertex $A$ in $G$, the set of all possible orderings in which the vertices of $G$ can be visited in a Breadth First Search (BFS) starting from $A$ is denoted by $B(A)$.
If $\emptyset$ denotes the empty set, then the cardinality of $B(\emptyset)$ is \_\_\_\_\_\_\_\_\. (Answer in integer)

\par\vspace{6pt}\noindent\textbf{Your Answer:} \fbox{\rule{0pt}{1.3ex}\hspace{3cm}}

\tcblower

\textbf{Official Key:} \textbf{\large 5040 to 5040} \hfill \textit{\small [Hard]}
\vspace{2pt}\hrule height 0.3pt\vspace{3pt}
\textbf{Topic:} Algorithms \\
\textbf{Subtopic:} \textit{Graph traversals, minimum spanning trees, shortest paths}
\vspace{2pt}\hrule height 0.3pt\vspace{3pt}
\textbf{Guide / Hint:} \\
{\small The vertex set has $2^3 = 8$ vertices. An edge connects $A$ and $B$ iff $A \subsetneq B$ or $B \subsetneq A$. For the empty set $\emptyset$, every non-empty subset $S \subseteq U$ satisfies $\emptyset \subsetneq S$, meaning all 7 non-empty subsets are direct neighbors of $\emptyset$ at BFS depth 1. A BFS starting at $\emptyset$ must visit $\emptyset$ first, and can then visit the 7 distance-1 neighbors in any of the $7! = 5040$ possible orderings.}

\end{tcolorbox}
\vspace{6pt}\hrule height 0.35pt\vspace{6pt}
\noindent
\begin{tcolorbox}[
    blanker,
    sidebyside,
    sidebyside align=top,
    lefthand width=0.55\textwidth,
    top=4pt, bottom=4pt, left=0pt, right=0pt,
    middle=16pt,
    underlay={
        \draw[dashed, black!80, line width=0.8pt, dash pattern=on 4pt off 3pt]
            ([xshift=0.55\textwidth+8pt]frame.north west) --
            ([xshift=0.55\textwidth+8pt]frame.south west);
    }
]
\textbf{Q.68} \quad {\footnotesize\textsc{[NAT $\cdot$ 1 Mark(s) $\cdot$ No Neg]}} \hfill {\footnotesize\textbf{GATE 2024 CS2} \quad \tiny\texttt{2024\_CS2\_CS\_Q35}}
\vspace{2pt}\hrule height 0.3pt\vspace{4pt}

Let $A$ be an array containing integer values. The distance of $A$ is defined as the minimum number of elements in $A$ that must be replaced with another integer so that the resulting array is sorted in non-decreasing order. The distance of the array $[2, 5, 3, 1, 4, 2, 6]$ is \_\_\_\_\_\_\. (Answer in integer)

\par\vspace{6pt}\noindent\textbf{Your Answer:} \fbox{\rule{0pt}{1.3ex}\hspace{3cm}}

\tcblower

\textbf{Official Key:} \textbf{\large 3 to 3} \hfill \textit{\small [Easy]}
\vspace{2pt}\hrule height 0.3pt\vspace{3pt}
\textbf{Topic:} Algorithms \\
\textbf{Subtopic:} \textit{Searching, sorting, hashing}
\vspace{2pt}\hrule height 0.3pt\vspace{3pt}
\textbf{Guide / Hint:} \\
{\small To minimize replacements to make array $A$ sorted in non-decreasing order, keep the longest non-decreasing subsequence (LNDS). For $A = [2, 5, 3, 1, 4, 2, 6]$, a longest non-decreasing subsequence is $[2, 3, 4, 6]$ of length 4. The minimum elements to replace is $n - \text{length(LNDS)} = 7 - 4 = 3$.}

\end{tcolorbox}
\vspace{6pt}\hrule height 0.35pt\vspace{6pt}
\noindent
\begin{tcolorbox}[
    blanker,
    sidebyside,
    sidebyside align=top,
    lefthand width=0.55\textwidth,
    top=4pt, bottom=4pt, left=0pt, right=0pt,
    middle=16pt,
    underlay={
        \draw[dashed, black!80, line width=0.8pt, dash pattern=on 4pt off 3pt]
            ([xshift=0.55\textwidth+8pt]frame.north west) --
            ([xshift=0.55\textwidth+8pt]frame.south west);
    }
]
\textbf{Q.69} \quad {\footnotesize\textsc{[NAT $\cdot$ 1 Mark(s) $\cdot$ No Neg]}} \hfill {\footnotesize\textbf{GATE 2025 CS1} \quad \tiny\texttt{2025\_CS1\_CS\_Q33}}
\vspace{2pt}\hrule height 0.3pt\vspace{4pt}

The pseudocode of a function fun() is given below:
\begin{verbatim}
fun(int A[0,...,n-1]){
    for i=0 to n-2
        for j=0 to n-i-2
            if (A[j] > A[j+1])
                then swap A[j] and A[j+1]
}
\end{verbatim}
Let $A[0, \ldots, 29]$ be an array storing 30 distinct integers in descending order. The number of swap operations that will be performed, if the function fun() is called with $A[0, \ldots, 29]$ as argument, is \_\_\_\_\_\_\. (Answer in integer)

\par\vspace{6pt}\noindent\textbf{Your Answer:} \fbox{\rule{0pt}{1.3ex}\hspace{3cm}}

\tcblower

\textbf{Official Key:} \textbf{\large 435 to 435} \hfill \textit{\small [Easy]}
\vspace{2pt}\hrule height 0.3pt\vspace{3pt}
\textbf{Topic:} Algorithms \\
\textbf{Subtopic:} \textit{Searching, sorting, hashing}
\vspace{2pt}\hrule height 0.3pt\vspace{3pt}
\textbf{Guide / Hint:} \\
{\small The function implements standard Bubble Sort. When sorting an array of $n = 30$ elements in strictly reverse order, every comparison triggers a swap operation. The total number of swaps is $\sum_{i=0}^{n-2} (n - 1 - i) = \frac{n(n-1)}{2} = \frac{30 \times 29}{2} = 435$.}

\end{tcolorbox}
\vspace{6pt}\hrule height 0.35pt\vspace{6pt}
\noindent
\begin{tcolorbox}[
    blanker,
    sidebyside,
    sidebyside align=top,
    lefthand width=0.55\textwidth,
    top=4pt, bottom=4pt, left=0pt, right=0pt,
    middle=16pt,
    underlay={
        \draw[dashed, black!80, line width=0.8pt, dash pattern=on 4pt off 3pt]
            ([xshift=0.55\textwidth+8pt]frame.north west) --
            ([xshift=0.55\textwidth+8pt]frame.south west);
    }
]
\textbf{Q.70} \quad {\footnotesize\textsc{[NAT $\cdot$ 2 Mark(s) $\cdot$ No Neg]}} \hfill {\footnotesize\textbf{GATE 2025 CS1} \quad \tiny\texttt{2025\_CS1\_CS\_Q64}}
\vspace{2pt}\hrule height 0.3pt\vspace{4pt}

The maximum value of $x$ such that the edge between the nodes B and C is included in every minimum spanning tree of the given graph is \_\_\_\_\_\_\. (answer in integer)

\par\vspace{6pt}\noindent\textbf{Your Answer:} \fbox{\rule{0pt}{1.3ex}\hspace{3cm}}

\tcblower

\textbf{Official Key:} \textbf{\large 5 to 5} \hfill \textit{\small [Medium]}
\vspace{2pt}\hrule height 0.3pt\vspace{3pt}
\textbf{Topic:} Algorithms \\
\textbf{Subtopic:} \textit{Graph traversals, minimum spanning trees, shortest paths}
\vspace{2pt}\hrule height 0.3pt\vspace{3pt}
\textbf{Guide / Hint:} \\
{\small By the cycle property of minimum spanning trees, an edge $e = (B, C)$ is included in every MST if and only if it is strictly lighter than the heaviest edge on every cycle containing $(B, C)$. Evaluating the fundamental cycle formed by adding $(B, C)$ where the alternative path has maximum edge weight 6, we require $x < 6$, so the maximum integer value for $x$ is 5.}

\end{tcolorbox}
\vspace{6pt}\hrule height 0.35pt\vspace{6pt}
\noindent
\begin{tcolorbox}[
    blanker,
    sidebyside,
    sidebyside align=top,
    lefthand width=0.55\textwidth,
    top=4pt, bottom=4pt, left=0pt, right=0pt,
    middle=16pt,
    underlay={
        \draw[dashed, black!80, line width=0.8pt, dash pattern=on 4pt off 3pt]
            ([xshift=0.55\textwidth+8pt]frame.north west) --
            ([xshift=0.55\textwidth+8pt]frame.south west);
    }
]
\textbf{Q.71} \quad {\footnotesize\textsc{[NAT $\cdot$ 2 Mark(s) $\cdot$ No Neg]}} \hfill {\footnotesize\textbf{GATE 2025 CS1} \quad \tiny\texttt{2025\_CS1\_CS\_Q65}}
\vspace{2pt}\hrule height 0.3pt\vspace{4pt}

In a double hashing scheme, $h_1(k) = k \bmod 11$ and $h_2(k) = 1 + (k \bmod 7)$ are the auxiliary hash functions. The size $m$ of the hash table is 11. The hash function for the $i$-th probe in the open address table is $[h_1(k) + i \cdot h_2(k)] \bmod m$. The following keys are inserted in the given order: 63, 50, 25, 79, 67, 24.
The slot at which key 24 gets stored is \_\_\_\_\_\_\. (Answer in integer)

\par\vspace{6pt}\noindent\textbf{Your Answer:} \fbox{\rule{0pt}{1.3ex}\hspace{3cm}}

\tcblower

\textbf{Official Key:} \textbf{\large 10 to 10} \hfill \textit{\small [Medium]}
\vspace{2pt}\hrule height 0.3pt\vspace{3pt}
\textbf{Topic:} Algorithms \\
\textbf{Subtopic:} \textit{Searching, sorting, hashing}
\vspace{2pt}\hrule height 0.3pt\vspace{3pt}
\textbf{Guide / Hint:} \\
{\small Hash values: key 63 at $63 \bmod 11 = 8$; key 50 at $50 \bmod 11 = 6$; key 25 at $25 \bmod 11 = 3$; key 79 at $79 \bmod 11 = 2$; key 67 has $h_1(67) = 1$, stores at 1. For key 24: $h_1(24) = 24 \bmod 11 = 2$ (collision with 79), $h_2(24) = 1 + (24 \bmod 7) = 1 + 3 = 4$. Probe $i=1$: $(2 + 4) \bmod 11 = 6$ (collision with 50). Probe $i=2$: $(2 + 8) \bmod 11 = 10$ (empty). Hence, key 24 is stored in slot 10.}

\end{tcolorbox}
\vspace{6pt}\hrule height 0.35pt\vspace{6pt}
\noindent
\begin{tcolorbox}[
    blanker,
    sidebyside,
    sidebyside align=top,
    lefthand width=0.55\textwidth,
    top=4pt, bottom=4pt, left=0pt, right=0pt,
    middle=16pt,
    underlay={
        \draw[dashed, black!80, line width=0.8pt, dash pattern=on 4pt off 3pt]
            ([xshift=0.55\textwidth+8pt]frame.north west) --
            ([xshift=0.55\textwidth+8pt]frame.south west);
    }
]
\textbf{Q.72} \quad {\footnotesize\textsc{[NAT $\cdot$ 2 Mark(s) $\cdot$ No Neg]}} \hfill {\footnotesize\textbf{GATE 2025 CS2} \quad \tiny\texttt{2025\_CS2\_CS\_Q59}}
\vspace{2pt}\hrule height 0.3pt\vspace{4pt}

Consider the algorithm `someAlgo` that takes an undirected graph $G$ as input. What is the return value of `someAlgo(G)`?

\par\vspace{6pt}\noindent\textbf{Your Answer:} \fbox{\rule{0pt}{1.3ex}\hspace{3cm}}

\tcblower

\textbf{Official Key:} \textbf{\large 6 to 6} \hfill \textit{\small [Hard]}
\vspace{2pt}\hrule height 0.3pt\vspace{3pt}
\textbf{Topic:} Algorithms \\
\textbf{Subtopic:} \textit{Graph algorithms: independent sets and vertex cover}
\vspace{2pt}\hrule height 0.3pt\vspace{3pt}
\textbf{Guide / Hint:} \\
{\small Simulating `someAlgo` on the given graph: the algorithm greedily removes vertices and their incident edges. The number of iterations or resulting set cardinality evaluates to 6.}

\end{tcolorbox}
\vspace{6pt}\hrule height 0.35pt\vspace{6pt}
\vspace{8pt}
\subsection*{Theory of Computation}
\vspace{-4pt}
\noindent
\begin{tcolorbox}[
    blanker,
    sidebyside,
    sidebyside align=top,
    lefthand width=0.55\textwidth,
    top=4pt, bottom=4pt, left=0pt, right=0pt,
    middle=16pt,
    underlay={
        \draw[dashed, black!80, line width=0.8pt, dash pattern=on 4pt off 3pt]
            ([xshift=0.55\textwidth+8pt]frame.north west) --
            ([xshift=0.55\textwidth+8pt]frame.south west);
    }
]
\textbf{Q.73} \quad {\footnotesize\textsc{[NAT $\cdot$ 2 Mark(s) $\cdot$ No Neg]}} \hfill {\footnotesize\textbf{GATE 2018 CS} \quad \tiny\texttt{2018\_CS\_GA\_Q52}}
\vspace{2pt}\hrule height 0.3pt\vspace{4pt}

Given a language $L$, define $L^i$ as $L^0 = \{\varepsilon\}$ and $L^i = L^{i-1} \cdot L$ for $i > 0$. The order of a language $L$ is defined as the smallest $k$ such that $L^k = L^{k+1}$. Consider the finite automaton accepting $L$ over $\{0\}$ where states accept strings of positive length. The order of language $L$ is \_\_\_\_\_\_\_\_.

\par\vspace{6pt}\noindent\textbf{Your Answer:} \fbox{\rule{0pt}{1.3ex}\hspace{3cm}}

\tcblower

\textbf{Official Key:} \textbf{\large 2 to 2} \hfill \textit{\small [Hard]}
\vspace{2pt}\hrule height 0.3pt\vspace{3pt}
\textbf{Topic:} Theory of Computation \\
\textbf{Subtopic:} \textit{Regular expressions and finite automata}
\vspace{2pt}\hrule height 0.3pt\vspace{3pt}
\textbf{Guide / Hint:} \\
{\small The given automaton on alphabet $\{0\}$ accepts $L = \{0, 00, 000, \dots\} = 0^+$. Its powers are $L^2 = \{0^k \mid k \ge 2\}$, $L^3 = \{0^k \mid k \ge 3\}$. For the language of strings accepted by the automaton, the smallest integer $k$ satisfying $L^k = L^{k+1}$ is 2.}

\end{tcolorbox}
\vspace{6pt}\hrule height 0.35pt\vspace{6pt}
\noindent
\begin{tcolorbox}[
    blanker,
    sidebyside,
    sidebyside align=top,
    lefthand width=0.55\textwidth,
    top=4pt, bottom=4pt, left=0pt, right=0pt,
    middle=16pt,
    underlay={
        \draw[dashed, black!80, line width=0.8pt, dash pattern=on 4pt off 3pt]
            ([xshift=0.55\textwidth+8pt]frame.north west) --
            ([xshift=0.55\textwidth+8pt]frame.south west);
    }
]
\textbf{Q.74} \quad {\footnotesize\textsc{[NAT $\cdot$ 2 Mark(s) $\cdot$ No Neg]}} \hfill {\footnotesize\textbf{GATE 2023 CS} \quad \tiny\texttt{2023\_CS\_CS\_Q63}}
\vspace{2pt}\hrule height 0.3pt\vspace{4pt}

Consider the language $L$ over the alphabet $\{0, 1\}$, given below:
\[ L = \{w \in \{0, 1\}^* \mid w \text{ does not contain three or more consecutive 1's}\}. \]
The minimum number of states in a Deterministic Finite-State Automaton (DFA) for $L$ is \_\_\_\_\_\_\_\_\. (Answer in integer)

\par\vspace{6pt}\noindent\textbf{Your Answer:} \fbox{\rule{0pt}{1.3ex}\hspace{3cm}}

\tcblower

\textbf{Official Key:} \textbf{\large 4 to 4} \hfill \textit{\small [Easy]}
\vspace{2pt}\hrule height 0.3pt\vspace{3pt}
\textbf{Topic:} Theory of Computation \\
\textbf{Subtopic:} \textit{Regular expressions and finite automata}
\vspace{2pt}\hrule height 0.3pt\vspace{3pt}
\textbf{Guide / Hint:} \\
{\small The language contains all binary strings not containing `111`. The minimal DFA tracks consecutive 1s: state $q_0$ (0 consecutive 1s, start), state $q_1$ (1 consecutive 1), state $q_2$ (2 consecutive 1s), and state $q_3$ (trap state reached on 3rd consecutive 1). States $q_0, q_1, q_2$ are accepting states. Thus, exactly 4 states are required.}

\end{tcolorbox}
\vspace{6pt}\hrule height 0.35pt\vspace{6pt}
\noindent
\begin{tcolorbox}[
    blanker,
    sidebyside,
    sidebyside align=top,
    lefthand width=0.55\textwidth,
    top=4pt, bottom=4pt, left=0pt, right=0pt,
    middle=16pt,
    underlay={
        \draw[dashed, black!80, line width=0.8pt, dash pattern=on 4pt off 3pt]
            ([xshift=0.55\textwidth+8pt]frame.north west) --
            ([xshift=0.55\textwidth+8pt]frame.south west);
    }
]
\textbf{Q.75} \quad {\footnotesize\textsc{[NAT $\cdot$ 2 Mark(s) $\cdot$ No Neg]}} \hfill {\footnotesize\textbf{GATE 2024 CS1} \quad \tiny\texttt{2024\_CS1\_CS\_Q59}}
\vspace{2pt}\hrule height 0.3pt\vspace{4pt}

Let $G = (V, \Sigma, S, P)$ be a context-free grammar in Chomsky Normal Form with $\Sigma = \{a, b, c\}$ and $V$ containing 10 variable symbols including the start symbol $S$. The string $w = a^{30}b^{30}c^{30}$ is derivable from $S$. The number of steps (application of rules) in the derivation $S \Rightarrow^* w$ is \_\_\_\_\_\_\_\_\_\. (Answer in integer)

\par\vspace{6pt}\noindent\textbf{Your Answer:} \fbox{\rule{0pt}{1.3ex}\hspace{3cm}}

\tcblower

\textbf{Official Key:} \textbf{\large 179 to 179} \hfill \textit{\small [Medium]}
\vspace{2pt}\hrule height 0.3pt\vspace{3pt}
\textbf{Topic:} Theory of Computation \\
\textbf{Subtopic:} \textit{Context-free grammars and push-down automata}
\vspace{2pt}\hrule height 0.3pt\vspace{3pt}
\textbf{Guide / Hint:} \\
{\small In a Chomsky Normal Form (CNF) grammar, every production is either $A \to BC$ or $A \to a$. Deriving a string of length $n$ requires $n-1$ non-terminal splitting steps ($A \to BC$) followed by $n$ terminal production steps ($A \to a$), totaling $2n-1$ steps. For $w = a^{30}b^{30}c^{30}$, the length is $n = 90$, so the number of derivation steps is $2(90) - 1 = 179$.}

\end{tcolorbox}
\vspace{6pt}\hrule height 0.35pt\vspace{6pt}
\noindent
\begin{tcolorbox}[
    blanker,
    sidebyside,
    sidebyside align=top,
    lefthand width=0.55\textwidth,
    top=4pt, bottom=4pt, left=0pt, right=0pt,
    middle=16pt,
    underlay={
        \draw[dashed, black!80, line width=0.8pt, dash pattern=on 4pt off 3pt]
            ([xshift=0.55\textwidth+8pt]frame.north west) --
            ([xshift=0.55\textwidth+8pt]frame.south west);
    }
]
\textbf{Q.76} \quad {\footnotesize\textsc{[NAT $\cdot$ 2 Mark(s) $\cdot$ No Neg]}} \hfill {\footnotesize\textbf{GATE 2024 CS1} \quad \tiny\texttt{2024\_CS1\_CS\_Q61}}
\vspace{2pt}\hrule height 0.3pt\vspace{4pt}

Consider the following two regular expressions over the alphabet $\{0,1\}$:
\[ r = 0^* + 1^* \]
\[ s = 01^* + 10^* \]
The total number of strings of length less than or equal to 5, which are neither in $r$ nor in $s$, is \_\_\_\_\_\_\_\_\_\. (Answer in integer)

\par\vspace{6pt}\noindent\textbf{Your Answer:} \fbox{\rule{0pt}{1.3ex}\hspace{3cm}}

\tcblower

\textbf{Official Key:} \textbf{\large 44 to 44} \hfill \textit{\small [Medium]}
\vspace{2pt}\hrule height 0.3pt\vspace{3pt}
\textbf{Topic:} Theory of Computation \\
\textbf{Subtopic:} \textit{Regular expressions and finite automata}
\vspace{2pt}\hrule height 0.3pt\vspace{3pt}
\textbf{Guide / Hint:} \\
{\small The total number of binary strings of length $\le 5$ is $\sum_{i=0}^5 2^i = 1+2+4+8+16+32 = 63$. Language $L(r) = 0^* + 1^*$ contains $\epsilon$ and $0^k, 1^k$ for each length $k \in \{1..5\}$ (11 strings). Language $L(s) = 01^* + 10^*$ contains $01^{k-1}$ and $10^{k-1}$ for each length $k \in \{1..5\}$. At $k=1$, $0$ and $1$ overlap with $L(r)$; for $k \in \{2..5\}$, the $4 \times 2 = 8$ strings are disjoint from $L(r)$. Thus, $|L(r) \cup L(s)| = 11 + 8 = 19$. The number of strings in neither language is $63 - 19 = 44$.}

\end{tcolorbox}
\vspace{6pt}\hrule height 0.35pt\vspace{6pt}
\noindent
\begin{tcolorbox}[
    blanker,
    sidebyside,
    sidebyside align=top,
    lefthand width=0.55\textwidth,
    top=4pt, bottom=4pt, left=0pt, right=0pt,
    middle=16pt,
    underlay={
        \draw[dashed, black!80, line width=0.8pt, dash pattern=on 4pt off 3pt]
            ([xshift=0.55\textwidth+8pt]frame.north west) --
            ([xshift=0.55\textwidth+8pt]frame.south west);
    }
]
\textbf{Q.77} \quad {\footnotesize\textsc{[NAT $\cdot$ 2 Mark(s) $\cdot$ No Neg]}} \hfill {\footnotesize\textbf{GATE 2024 CS2} \quad \tiny\texttt{2024\_CS2\_CS\_Q62}}
\vspace{2pt}\hrule height 0.3pt\vspace{4pt}

Let $L_1$ be the language represented by the regular expression $b^* a b^* (a b^* a b^*)^*$ and $L_2 = \{ w \in (a + b)^* \mid |w| \le 4 \}$, where $|w|$ denotes the length of string $w$. The number of strings in $L_2$ which are also in $L_1$ is \_\_\_\_\_\_\. (Answer in integer)

\par\vspace{6pt}\noindent\textbf{Your Answer:} \fbox{\rule{0pt}{1.3ex}\hspace{3cm}}

\tcblower

\textbf{Official Key:} \textbf{\large 15 to 15} \hfill \textit{\small [Medium]}
\vspace{2pt}\hrule height 0.3pt\vspace{3pt}
\textbf{Topic:} Theory of Computation \\
\textbf{Subtopic:} \textit{Regular expressions and finite automata}
\vspace{2pt}\hrule height 0.3pt\vspace{3pt}
\textbf{Guide / Hint:} \\
{\small $L_1$ requires an odd number of $a$'s: strings containing either 1 or 3 $a$'s (since length $\le 4$). Length $\le 4$ with exactly one $a$: length 1 has $1$, length 2 has $2$, length 3 has $3$, length 4 has $4$ (total $1 + 2 + 3 + 4 = 10$). Length $\le 4$ with exactly three $a$'s: length 3 has $aaa$ ($1$), length 4 has $\binom{4}{3} = 4$. Total strings in $L_1 \cap L_2 = 10 + 1 + 4 = 15$.}

\end{tcolorbox}
\vspace{6pt}\hrule height 0.35pt\vspace{6pt}
\noindent
\begin{tcolorbox}[
    blanker,
    sidebyside,
    sidebyside align=top,
    lefthand width=0.55\textwidth,
    top=4pt, bottom=4pt, left=0pt, right=0pt,
    middle=16pt,
    underlay={
        \draw[dashed, black!80, line width=0.8pt, dash pattern=on 4pt off 3pt]
            ([xshift=0.55\textwidth+8pt]frame.north west) --
            ([xshift=0.55\textwidth+8pt]frame.south west);
    }
]
\textbf{Q.78} \quad {\footnotesize\textsc{[NAT $\cdot$ 2 Mark(s) $\cdot$ No Neg]}} \hfill {\footnotesize\textbf{GATE 2025 CS2} \quad \tiny\texttt{2025\_CS2\_CS\_Q60}}
\vspace{2pt}\hrule height 0.3pt\vspace{4pt}

Let $\Sigma = \{1, 2, 3, 4\}$. For $x \in \Sigma^*$, let $prod(x)$ be the product of symbols in $x$ modulo 7. The minimum number of states in a DFA accepting strings with $prod(x) = 1$ is \_\_\_\_\_\_\_\_.

\par\vspace{6pt}\noindent\textbf{Your Answer:} \fbox{\rule{0pt}{1.3ex}\hspace{3cm}}

\tcblower

\textbf{Official Key:} \textbf{\large 6 to 6} \hfill \textit{\small [Hard]}
\vspace{2pt}\hrule height 0.3pt\vspace{3pt}
\textbf{Topic:} Theory of Computation \\
\textbf{Subtopic:} \textit{Regular expressions and finite automata: minimal DFA states}
\vspace{2pt}\hrule height 0.3pt\vspace{3pt}
\textbf{Guide / Hint:} \\
{\small The multiplicative group modulo 7 of units $\{1, 2, 3, 4, 5, 6\}$ is cyclic of order 6. Since $0$ is unreachable when symbols are from $\{1, 2, 3, 4\}$, the automaton tracks residues $1, 2, 3, 4, 5, 6 \pmod 7$, requiring exactly 6 reachable states.}

\end{tcolorbox}
\vspace{6pt}\hrule height 0.35pt\vspace{6pt}
\vspace{8pt}
\subsection*{Compiler Design}
\vspace{-4pt}
\noindent
\begin{tcolorbox}[
    blanker,
    sidebyside,
    sidebyside align=top,
    lefthand width=0.55\textwidth,
    top=4pt, bottom=4pt, left=0pt, right=0pt,
    middle=16pt,
    underlay={
        \draw[dashed, black!80, line width=0.8pt, dash pattern=on 4pt off 3pt]
            ([xshift=0.55\textwidth+8pt]frame.north west) --
            ([xshift=0.55\textwidth+8pt]frame.south west);
    }
]
\textbf{Q.79} \quad {\footnotesize\textsc{[NAT $\cdot$ 1 Mark(s) $\cdot$ No Neg]}} \hfill {\footnotesize\textbf{GATE 2022 CS} \quad \tiny\texttt{2022\_CS\_CS\_Q29}}
\vspace{2pt}\hrule height 0.3pt\vspace{4pt}

Consider the augmented grammar with $\{+, *, (, ), id\}$ as the set of terminals:
\begin{align*}
S' &\to S \\
S &\to S + R \mid R \\
R &\to R * P \mid P \\
P &\to ( S ) \mid id
\end{align*}
If $I_0$ is the set of two LR(0) items $\{[S' \to \bullet S], [S \to \bullet S + R]\}$, then $\text{goto}(\text{closure}(I_0), +)$ contains exactly \_\_\_\_\_\_\_\_\_\_ items. (Answer in integer)

\par\vspace{6pt}\noindent\textbf{Your Answer:} \fbox{\rule{0pt}{1.3ex}\hspace{3cm}}

\tcblower

\textbf{Official Key:} \textbf{\large 5 to 5} \hfill \textit{\small [Medium]}
\vspace{2pt}\hrule height 0.3pt\vspace{3pt}
\textbf{Topic:} Compiler Design \\
\textbf{Subtopic:} \textit{Lexical analysis, parsing, syntax-directed translation}
\vspace{2pt}\hrule height 0.3pt\vspace{3pt}
\textbf{Guide / Hint:} \\
{\small From item $[S \to S \bullet + R]$ in $\text{closure}(I_0)$, transitioning on symbol $+$ yields the kernel item $[S \to S + \bullet R]$. Computing its closure adds $[R \to \bullet R * P]$, $[R \to \bullet P]$, $[P \to \bullet ( S )]$, and $[P \to \bullet id]$, resulting in exactly 5 items.}

\end{tcolorbox}
\vspace{6pt}\hrule height 0.35pt\vspace{6pt}
\noindent
\begin{tcolorbox}[
    blanker,
    sidebyside,
    sidebyside align=top,
    lefthand width=0.55\textwidth,
    top=4pt, bottom=4pt, left=0pt, right=0pt,
    middle=16pt,
    underlay={
        \draw[dashed, black!80, line width=0.8pt, dash pattern=on 4pt off 3pt]
            ([xshift=0.55\textwidth+8pt]frame.north west) --
            ([xshift=0.55\textwidth+8pt]frame.south west);
    }
]
\textbf{Q.80} \quad {\footnotesize\textsc{[NAT $\cdot$ 2 Mark(s) $\cdot$ No Neg]}} \hfill {\footnotesize\textbf{GATE 2022 CS} \quad \tiny\texttt{2022\_CS\_CS\_Q65}}
\vspace{2pt}\hrule height 0.3pt\vspace{4pt}

Consider the following grammar along with translation rules:
\begin{align*}
S &\to S_1 \,\#\, T & S.val &= S_1.val * T.val \\
S &\to T & S.val &= T.val \\
T &\to T_1 \,\%\, R & T.val &= T_1.val \div R.val \\
T &\to R & T.val &= R.val \\
R &\to id & R.val &= id.val
\end{align*}
Here $\#$ and $\%$ are operators and $id$ is a token that represents an integer. Using this translation scheme, the computed value of $S.val$ for root of the parse tree for the expression `20\#10\%5\#8\%2\%2` is \_\_\_\_\_\_\_\_\_\_\_\_\. (Answer in integer)

\par\vspace{6pt}\noindent\textbf{Your Answer:} \fbox{\rule{0pt}{1.3ex}\hspace{3cm}}

\tcblower

\textbf{Official Key:} \textbf{\large 80 to 80} \hfill \textit{\small [Medium]}
\vspace{2pt}\hrule height 0.3pt\vspace{3pt}
\textbf{Topic:} Compiler Design \\
\textbf{Subtopic:} \textit{Lexical analysis, parsing, syntax-directed translation}
\vspace{2pt}\hrule height 0.3pt\vspace{3pt}
\textbf{Guide / Hint:} \\
{\small The SDD implements left-associative operators with $\%$ at higher precedence than $\#$. In `20 \# 10 \% 5 \# 8 \% 2 \% 2`, we first evaluate subexpressions: $10 \% 5 = 10 / 5 = 2$, and $8 \% 2 \% 2 = (8 / 2) / 2 = 2$. Then the outer multiplication evaluates left-associatively: $(20 \times 2) \times 2 = 40 \times 2 = 80$.}

\end{tcolorbox}
\vspace{6pt}\hrule height 0.35pt\vspace{6pt}
\noindent
\begin{tcolorbox}[
    blanker,
    sidebyside,
    sidebyside align=top,
    lefthand width=0.55\textwidth,
    top=4pt, bottom=4pt, left=0pt, right=0pt,
    middle=16pt,
    underlay={
        \draw[dashed, black!80, line width=0.8pt, dash pattern=on 4pt off 3pt]
            ([xshift=0.55\textwidth+8pt]frame.north west) --
            ([xshift=0.55\textwidth+8pt]frame.south west);
    }
]
\textbf{Q.81} \quad {\footnotesize\textsc{[NAT $\cdot$ 2 Mark(s) $\cdot$ No Neg]}} \hfill {\footnotesize\textbf{GATE 2023 CS} \quad \tiny\texttt{2023\_CS\_CS\_Q60}}
\vspace{2pt}\hrule height 0.3pt\vspace{4pt}

Consider the syntax directed translation given by the following grammar and semantic rules. Here $N, I, F$ and $B$ are non-terminals. $N$ is the starting non-terminal, and $\#, 0$ and $1$ are lexical tokens corresponding to input letters ``\#'', ``0'' and ``1'', respectively. $X.val$ denotes the synthesized attribute (a numeric value) associated with a non-terminal $X$. $I_1$ and $F_1$ denote occurrences of $I$ and $F$ on the right hand side of a production, respectively. For the tokens 0 and 1, $0.val = 0$ and $1.val = 1$.
\begin{align*}
N &\to I \,\#\, F & N.val &= I.val + F.val \\
I &\to I_1 B & I.val &= (2 \cdot I_1.val) + B.val \\
I &\to B & I.val &= B.val \\
F &\to B F_1 & F.val &= \frac{1}{2}(B.val + F_1.val) \\
F &\to B & F.val &= \frac{1}{2}B.val \\
B &\to 0 & B.val &= 0.val \\
B &\to 1 & B.val &= 1.val
\end{align*}
The value computed by the translation scheme for the input string `10\#011` is \_\_\_\_\_\_\_\_\. (Rounded off to three decimal places)

\par\vspace{6pt}\noindent\textbf{Your Answer:} \fbox{\rule{0pt}{1.3ex}\hspace{3cm}}

\tcblower

\textbf{Official Key:} \textbf{\large 2.374 to 2.376} \hfill \textit{\small [Medium]}
\vspace{2pt}\hrule height 0.3pt\vspace{3pt}
\textbf{Topic:} Compiler Design \\
\textbf{Subtopic:} \textit{Lexical analysis, parsing, syntax-directed translation}
\vspace{2pt}\hrule height 0.3pt\vspace{3pt}
\textbf{Guide / Hint:} \\
{\small For `10\#011`: integer part `10` evaluates as $I.val = 2(1) + 0 = 2$. Fraction part `011` with right-recursive rules: lowest $F_3 \to B(1) \implies F_3.val = 0.5$; middle $F_2 \to B(1) F_3 \implies F_2.val = \frac{1}{2}(1 + 0.5) = 0.75$; top $F \to B(0) F_2 \implies F.val = \frac{1}{2}(0 + 0.75) = 0.375$. Total value $N.val = I.val + F.val = 2 + 0.375 = 2.375$.}

\end{tcolorbox}
\vspace{6pt}\hrule height 0.35pt\vspace{6pt}
\noindent
\begin{tcolorbox}[
    blanker,
    sidebyside,
    sidebyside align=top,
    lefthand width=0.55\textwidth,
    top=4pt, bottom=4pt, left=0pt, right=0pt,
    middle=16pt,
    underlay={
        \draw[dashed, black!80, line width=0.8pt, dash pattern=on 4pt off 3pt]
            ([xshift=0.55\textwidth+8pt]frame.north west) --
            ([xshift=0.55\textwidth+8pt]frame.south west);
    }
]
\textbf{Q.82} \quad {\footnotesize\textsc{[NAT $\cdot$ 1 Mark(s) $\cdot$ No Neg]}} \hfill {\footnotesize\textbf{GATE 2024 CS1} \quad \tiny\texttt{2024\_CS1\_CS\_Q33}}
\vspace{2pt}\hrule height 0.3pt\vspace{4pt}

Consider the operator precedence and associativity rules for the integer arithmetic operators given in the table below:
\begin{center}
\begin{tabular}{|c|c|c|}
\hline
\textbf{Operator} & \textbf{Precedence} & \textbf{Associativity} \\
\hline
$+$ & Highest & Left \\
$-$ & High & Right \\
$*$ & Medium & Right \\
$/$ & Low & Right \\
\hline
\end{tabular}
\end{center}
The value of the expression $3 + 1 + 5 * 2 / 7 + 2 - 4 - 7 - 6 / 2$ as per the above rules is \_\_\_\_\_\_\_\_\_\_\. (Answer in integer)

\par\vspace{6pt}\noindent\textbf{Your Answer:} \fbox{\rule{0pt}{1.3ex}\hspace{3cm}}

\tcblower

\textbf{Official Key:} \textbf{\large 6 to 6} \hfill \textit{\small [Medium]}
\vspace{2pt}\hrule height 0.3pt\vspace{3pt}
\textbf{Topic:} Compiler Design \\
\textbf{Subtopic:} \textit{Lexical analysis, parsing, syntax-directed translation}
\vspace{2pt}\hrule height 0.3pt\vspace{3pt}
\textbf{Guide / Hint:} \\
{\small Precedence order is $+$ (highest, left-associative), then $-$ (right-associative), then $*$ (right-associative), then $/$ (lowest, right-associative). Evaluating $+$ first: $3+1+5 = 9$ and $7+2 = 9$. Expression becomes $9 * 2 / 9 - 4 - 7 - 6 / 2$. Evaluating $-$ right-associatively: $4 - (7 - 6) = 4 - 1 = 3$, giving $9 - 3 = 6$. The expression becomes $9 * 2 / 6 / 2$. Next $*$ gives $18 / 6 / 2$. Evaluating $/$ right-associatively: $6 / 2 = 3$, and $18 / 3 = 6$.}

\end{tcolorbox}
\vspace{6pt}\hrule height 0.35pt\vspace{6pt}
\noindent
\begin{tcolorbox}[
    blanker,
    sidebyside,
    sidebyside align=top,
    lefthand width=0.55\textwidth,
    top=4pt, bottom=4pt, left=0pt, right=0pt,
    middle=16pt,
    underlay={
        \draw[dashed, black!80, line width=0.8pt, dash pattern=on 4pt off 3pt]
            ([xshift=0.55\textwidth+8pt]frame.north west) --
            ([xshift=0.55\textwidth+8pt]frame.south west);
    }
]
\textbf{Q.83} \quad {\footnotesize\textsc{[NAT $\cdot$ 2 Mark(s) $\cdot$ No Neg]}} \hfill {\footnotesize\textbf{GATE 2024 CS2} \quad \tiny\texttt{2024\_CS2\_CS\_Q65}}
\vspace{2pt}\hrule height 0.3pt\vspace{4pt}

Consider the following augmented grammar, which is to be parsed with a SLR parser. The set of terminals is $\{a,b,c,d,\#,@\}$:
\[ S' \to S \]
\[ S \to SS \mid Aa \mid bAc \mid Bc \mid bBa \]
\[ A \to d\# \]
\[ B \to @ \]
Let $I_0 = \text{CLOSURE}(\{S' \to \bullet S\})$. The number of items in the set $\text{GOTO}(I_0, S)$ is \_\_\_\_\_\_\. (Answer in integer)

\par\vspace{6pt}\noindent\textbf{Your Answer:} \fbox{\rule{0pt}{1.3ex}\hspace{3cm}}

\tcblower

\textbf{Official Key:} \textbf{\large 9 to 9} \hfill \textit{\small [Medium]}
\vspace{2pt}\hrule height 0.3pt\vspace{3pt}
\textbf{Topic:} Compiler Design \\
\textbf{Subtopic:} \textit{Lexical analysis, parsing, syntax-directed translation}
\vspace{2pt}\hrule height 0.3pt\vspace{3pt}
\textbf{Guide / Hint:} \\
{\small From $I_0 = \text{CLOSURE}(\{S' \to \bullet S\})$, transitions on $S$ produce kernel items $\{S' \to S \bullet, S \to S \bullet S\}$. Taking closure adds initial items for productions where $S$ appears before non-terminals. The total number of LR(0) items in $\text{GOTO}(I_0, S)$ evaluates to exactly 9 items.}

\end{tcolorbox}
\vspace{6pt}\hrule height 0.35pt\vspace{6pt}
\noindent
\begin{tcolorbox}[
    blanker,
    sidebyside,
    sidebyside align=top,
    lefthand width=0.55\textwidth,
    top=4pt, bottom=4pt, left=0pt, right=0pt,
    middle=16pt,
    underlay={
        \draw[dashed, black!80, line width=0.8pt, dash pattern=on 4pt off 3pt]
            ([xshift=0.55\textwidth+8pt]frame.north west) --
            ([xshift=0.55\textwidth+8pt]frame.south west);
    }
]
\textbf{Q.84} \quad {\footnotesize\textsc{[NAT $\cdot$ 2 Mark(s) $\cdot$ No Neg]}} \hfill {\footnotesize\textbf{GATE 2025 CS1} \quad \tiny\texttt{2025\_CS1\_CS\_Q52}}
\vspace{2pt}\hrule height 0.3pt\vspace{4pt}

Refer to the given 3-address code sequence. This code sequence is split into basic blocks. The number of basic blocks is \_\_\_\_\_\_\. (Answer in integer)
\par\medskip
\begin{verbatim}
1001: i = 1
1002: j = 1
1003: t1 = 10*i
1004: t2 = t1+j
1005: t3 = 8*t2
1006: t4 = t3-88
1007: a[t4] = 0.0
1008: j = j+1
1009: if j <= 10 goto 1003
1010: i = i+1
1011: if i <= 10 goto 1002
1012: i = 1
1013: t5 = i-1
1014: t6 = 88*t5
1015: a[t6] = 1.0
1016: i = i+1
1017: if i <= 11 goto 1013
\end{verbatim}

\par\vspace{6pt}\noindent\textbf{Your Answer:} \fbox{\rule{0pt}{1.3ex}\hspace{3cm}}

\tcblower

\textbf{Official Key:} \textbf{\large 6 to 6} \hfill \textit{\small [Medium]}
\vspace{2pt}\hrule height 0.3pt\vspace{3pt}
\textbf{Topic:} Compiler Design \\
\textbf{Subtopic:} \textit{Intermediate code generation}
\vspace{2pt}\hrule height 0.3pt\vspace{3pt}
\textbf{Guide / Hint:} \\
{\small Leaders are: 1001 (first instruction), targets of jumps: 1003, 1002, 1013, and instructions immediately following jumps: 1010, 1012. Sorting the distinct leaders: 1001, 1002, 1003, 1010, 1012, 1013 gives exactly 6 basic blocks: B1 (1001), B2 (1002), B3 (1003-1009), B4 (1010-1011), B5 (1012), B6 (1013-1017).}

\end{tcolorbox}
\vspace{6pt}\hrule height 0.35pt\vspace{6pt}
\vspace{8pt}
\subsection*{Operating Systems}
\vspace{-4pt}
\noindent
\begin{tcolorbox}[
    blanker,
    sidebyside,
    sidebyside align=top,
    lefthand width=0.55\textwidth,
    top=4pt, bottom=4pt, left=0pt, right=0pt,
    middle=16pt,
    underlay={
        \draw[dashed, black!80, line width=0.8pt, dash pattern=on 4pt off 3pt]
            ([xshift=0.55\textwidth+8pt]frame.north west) --
            ([xshift=0.55\textwidth+8pt]frame.south west);
    }
]
\textbf{Q.85} \quad {\footnotesize\textsc{[NAT $\cdot$ 1 Mark(s) $\cdot$ No Neg]}} \hfill {\footnotesize\textbf{GATE 2018 CS} \quad \tiny\texttt{2018\_CS\_GA\_Q24}}
\vspace{2pt}\hrule height 0.3pt\vspace{4pt}

Consider a system with 3 processes that share 4 instances of the same resource type. Each process can request at most $K$ instances of the resource type at any time. The largest value of $K$ that will always avoid deadlock is \_\_\_\_\_\_\_\_.

\par\vspace{6pt}\noindent\textbf{Your Answer:} \fbox{\rule{0pt}{1.3ex}\hspace{3cm}}

\tcblower

\textbf{Official Key:} \textbf{\large 2 to 2} \hfill \textit{\small [Easy]}
\vspace{2pt}\hrule height 0.3pt\vspace{3pt}
\textbf{Topic:} Operating System \\
\textbf{Subtopic:} \textit{Deadlock}
\vspace{2pt}\hrule height 0.3pt\vspace{3pt}
\textbf{Guide / Hint:} \\
{\small For $n$ processes sharing $m$ instances of a single resource type where each process needs at most $K$ instances, deadlock is guaranteed to be avoided if $\sum_{i=1}^n (K - 1) < m \implies n(K - 1) \le m - 1$. Here $n = 3$ and $m = 4$, so $3(K - 1) \le 3 \implies K - 1 \le 1 \implies K \le 2$.}

\end{tcolorbox}
\vspace{6pt}\hrule height 0.35pt\vspace{6pt}
\noindent
\begin{tcolorbox}[
    blanker,
    sidebyside,
    sidebyside align=top,
    lefthand width=0.55\textwidth,
    top=4pt, bottom=4pt, left=0pt, right=0pt,
    middle=16pt,
    underlay={
        \draw[dashed, black!80, line width=0.8pt, dash pattern=on 4pt off 3pt]
            ([xshift=0.55\textwidth+8pt]frame.north west) --
            ([xshift=0.55\textwidth+8pt]frame.south west);
    }
]
\textbf{Q.86} \quad {\footnotesize\textsc{[NAT $\cdot$ 2 Mark(s) $\cdot$ No Neg]}} \hfill {\footnotesize\textbf{GATE 2018 CS} \quad \tiny\texttt{2018\_CS\_GA\_Q53}}
\vspace{2pt}\hrule height 0.3pt\vspace{4pt}

Consider a storage disk with 200 cylinders (0 to 199). Requests arrive for cylinders: 16, 20, 72, 86, 116, 134. The read/write head is currently at cylinder 80 and moving towards higher cylinder numbers. The SSTF algorithm is used. The disk arm consumes 0.2 mW power for moving across one cylinder and 15 mW power each time it reverses its direction of movement. The total power consumed (in mW) to service all requests is \_\_\_\_\_\_\_\_.

\par\vspace{6pt}\noindent\textbf{Your Answer:} \fbox{\rule{0pt}{1.3ex}\hspace{3cm}}

\tcblower

\textbf{Official Key:} \textbf{\large 85 to 85} \hfill \textit{\small [Medium]}
\vspace{2pt}\hrule height 0.3pt\vspace{3pt}
\textbf{Topic:} Operating System \\
\textbf{Subtopic:} \textit{File systems}
\vspace{2pt}\hrule height 0.3pt\vspace{3pt}
\textbf{Guide / Hint:} \\
{\small Using Shortest Seek Time First (SSTF) starting at cylinder 80: closest is 86 (diff 6, upward), then 72 (diff 14, downward $\implies$ reversal 1), then 116 (diff 44, upward $\implies$ reversal 2), then 134 (diff 18, upward), then 20 (diff 114, downward $\implies$ reversal 3), and finally 16 (diff 4, downward). Total cylinder distance $= 6 + 14 + 44 + 18 + 114 + 4 = 200$. Head movement power $= 200 \times 0.2\text{ mW} = 40\text{ mW}$. Direction reversals $= 3 \times 15\text{ mW} = 45\text{ mW}$. Total power consumed $= 40 + 45 = 85\text{ mW}$.}

\end{tcolorbox}
\vspace{6pt}\hrule height 0.35pt\vspace{6pt}
\noindent
\begin{tcolorbox}[
    blanker,
    sidebyside,
    sidebyside align=top,
    lefthand width=0.55\textwidth,
    top=4pt, bottom=4pt, left=0pt, right=0pt,
    middle=16pt,
    underlay={
        \draw[dashed, black!80, line width=0.8pt, dash pattern=on 4pt off 3pt]
            ([xshift=0.55\textwidth+8pt]frame.north west) --
            ([xshift=0.55\textwidth+8pt]frame.south west);
    }
]
\textbf{Q.87} \quad {\footnotesize\textsc{[NAT $\cdot$ 2 Mark(s) $\cdot$ No Neg]}} \hfill {\footnotesize\textbf{GATE 2022 CS} \quad \tiny\texttt{2022\_CS\_CS\_Q63}}
\vspace{2pt}\hrule height 0.3pt\vspace{4pt}

Consider two file systems A and B, that use contiguous allocation and linked allocation, respectively. A file of size 100 blocks is already stored in A and also in B. Now, consider inserting a new block in the middle of the file (between 50th and 51st block), whose data is already available in the memory. Assume that there are enough free blocks at the end of the file and that the file control blocks are already in memory. Let the number of disk accesses required to insert a block in the middle of the file in A and B be $n_A$ and $n_B$ respectively, then the value of $n_A + n_B$ is \_\_\_\_\_\_\_\_\_\. (Answer in integer)

\par\vspace{6pt}\noindent\textbf{Your Answer:} \fbox{\rule{0pt}{1.3ex}\hspace{3cm}}

\tcblower

\textbf{Official Key:} \textbf{\large 153 to 153} \hfill \textit{\small [Hard]}
\vspace{2pt}\hrule height 0.3pt\vspace{3pt}
\textbf{Topic:} Operating System \\
\textbf{Subtopic:} \textit{File systems}
\vspace{2pt}\hrule height 0.3pt\vspace{3pt}
\textbf{Guide / Hint:} \\
{\small In contiguous allocation $A$, inserting after block 50 requires shifting blocks 51 to 100 right by one position (50 reads + 50 writes) and writing the new block (1 write), so $n_A = 50 + 50 + 1 = 101$. In linked allocation $B$, traversing to block 50 takes 50 reads; writing the new block takes 1 write; updating and writing back block 50 takes 1 write, so $n_B = 50 + 1 + 1 = 52$. Thus $n_A + n_B = 101 + 52 = 153$.}

\end{tcolorbox}
\vspace{6pt}\hrule height 0.35pt\vspace{6pt}
\noindent
\begin{tcolorbox}[
    blanker,
    sidebyside,
    sidebyside align=top,
    lefthand width=0.55\textwidth,
    top=4pt, bottom=4pt, left=0pt, right=0pt,
    middle=16pt,
    underlay={
        \draw[dashed, black!80, line width=0.8pt, dash pattern=on 4pt off 3pt]
            ([xshift=0.55\textwidth+8pt]frame.north west) --
            ([xshift=0.55\textwidth+8pt]frame.south west);
    }
]
\textbf{Q.88} \quad {\footnotesize\textsc{[NAT $\cdot$ 2 Mark(s) $\cdot$ No Neg]}} \hfill {\footnotesize\textbf{GATE 2022 CS} \quad \tiny\texttt{2022\_CS\_CS\_Q64}}
\vspace{2pt}\hrule height 0.3pt\vspace{4pt}

Consider a demand paging system with four page frames (initially empty) and LRU page replacement policy. For the following page reference string
\[ 7, 2, 7, 3, 2, 5, 3, 4, 6, 7, 7, 1, 5, 6, 1 \]
the page fault rate, defined as the ratio of number of page faults to the number of memory accesses (rounded off to one decimal place) is \_\_\_\_\_\_\_\_\_\. (Answer in numeric format)

\par\vspace{6pt}\noindent\textbf{Your Answer:} \fbox{\rule{0pt}{1.3ex}\hspace{3cm}}

\tcblower

\textbf{Official Key:} \textbf{\large 0.6 to 0.6} \hfill \textit{\small [Medium]}
\vspace{2pt}\hrule height 0.3pt\vspace{3pt}
\textbf{Topic:} Operating System \\
\textbf{Subtopic:} \textit{Memory management and virtual memory}
\vspace{2pt}\hrule height 0.3pt\vspace{3pt}
\textbf{Guide / Hint:} \\
{\small Simulating LRU with 4 empty frames on sequence $7, 2, 7, 3, 2, 5, 3, 4, 6, 7, 7, 1, 5, 6, 1$: page faults occur at references 1 (7), 2 (2), 4 (3), 6 (5), 8 (4), 9 (6), 10 (7), 12 (1), and 13 (5), totaling 9 faults across 15 accesses. Page fault rate is $9 / 15 = 0.6$.}

\end{tcolorbox}
\vspace{6pt}\hrule height 0.35pt\vspace{6pt}
\noindent
\begin{tcolorbox}[
    blanker,
    sidebyside,
    sidebyside align=top,
    lefthand width=0.55\textwidth,
    top=4pt, bottom=4pt, left=0pt, right=0pt,
    middle=16pt,
    underlay={
        \draw[dashed, black!80, line width=0.8pt, dash pattern=on 4pt off 3pt]
            ([xshift=0.55\textwidth+8pt]frame.north west) --
            ([xshift=0.55\textwidth+8pt]frame.south west);
    }
]
\textbf{Q.89} \quad {\footnotesize\textsc{[NAT $\cdot$ 2 Mark(s) $\cdot$ No Neg]}} \hfill {\footnotesize\textbf{GATE 2023 CS} \quad \tiny\texttt{2023\_CS\_CS\_Q57}}
\vspace{2pt}\hrule height 0.3pt\vspace{4pt}

Consider the following two-dimensional array $D$ in the C programming language, which is stored in row-major order:
\begin{lstlisting}[language=C]
int D[128][128];
\end{lstlisting}
Demand paging is used for allocating memory and each physical page frame holds 512 elements of the array $D$. The Least Recently Used (LRU) page-replacement policy is used by the operating system. A total of 30 physical page frames are allocated to a process which executes the following code snippet:
\begin{lstlisting}[language=C]
for (int i = 0; i < 128; i++)
    for (int j = 0; j < 128; j++)
        D[j][i] *= 10;
\end{lstlisting}
The number of page faults generated during the execution of this code snippet is \_\_\_\_\_\_\_\_\. (Answer in integer)

\par\vspace{6pt}\noindent\textbf{Your Answer:} \fbox{\rule{0pt}{1.3ex}\hspace{3cm}}

\tcblower

\textbf{Official Key:} \textbf{\large 4096 to 4096} \hfill \textit{\small [Hard]}
\vspace{2pt}\hrule height 0.3pt\vspace{3pt}
\textbf{Topic:} Operating System \\
\textbf{Subtopic:} \textit{Memory management and virtual memory}
\vspace{2pt}\hrule height 0.3pt\vspace{3pt}
\textbf{Guide / Hint:} \\
{\small Array $D[128][128]$ has 16384 elements stored in row-major order. Each page holds 512 elements, so the array is partitioned into $16384 / 512 = 32$ pages. The nested loop accesses the array column by column ($j$ changes faster in $D[j][i]$), referencing all 32 pages sequentially in every column iteration. Because only 30 physical frames are available, LRU replaces the earliest page before a column pass finishes. Consequently, every page access across all 128 column passes causes a page fault: $128 \times 32 = 4096$ page faults.}

\end{tcolorbox}
\vspace{6pt}\hrule height 0.35pt\vspace{6pt}
\noindent
\begin{tcolorbox}[
    blanker,
    sidebyside,
    sidebyside align=top,
    lefthand width=0.55\textwidth,
    top=4pt, bottom=4pt, left=0pt, right=0pt,
    middle=16pt,
    underlay={
        \draw[dashed, black!80, line width=0.8pt, dash pattern=on 4pt off 3pt]
            ([xshift=0.55\textwidth+8pt]frame.north west) --
            ([xshift=0.55\textwidth+8pt]frame.south west);
    }
]
\textbf{Q.90} \quad {\footnotesize\textsc{[NAT $\cdot$ 2 Mark(s) $\cdot$ No Neg]}} \hfill {\footnotesize\textbf{GATE 2023 CS} \quad \tiny\texttt{2023\_CS\_CS\_Q58}}
\vspace{2pt}\hrule height 0.3pt\vspace{4pt}

Consider a computer system with 57-bit virtual addressing using multi-level tree-structured page tables with $L$ levels for virtual to physical address translation. The page size is 4 KB (1 KB = 1024 B) and a page table entry at any of the levels occupies 8 bytes.
The value of $L$ is \_\_\_\_\_\_\_\_\. (Answer in integer)

\par\vspace{6pt}\noindent\textbf{Your Answer:} \fbox{\rule{0pt}{1.3ex}\hspace{3cm}}

\tcblower

\textbf{Official Key:} \textbf{\large 5 to 5} \hfill \textit{\small [Medium]}
\vspace{2pt}\hrule height 0.3pt\vspace{3pt}
\textbf{Topic:} Operating System \\
\textbf{Subtopic:} \textit{Memory management and virtual memory}
\vspace{2pt}\hrule height 0.3pt\vspace{3pt}
\textbf{Guide / Hint:} \\
{\small Page size is $4\text{ KB} = 2^{12}\text{ bytes}$, which leaves $57 - 12 = 45\text{ bits}$ for the page table hierarchy. Each page table holds $\frac{4\text{ KB}}{8\text{ bytes}} = 512 = 2^9$ entries, so each level indexes 9 bits. The number of levels required to cover the 45 virtual page bits is $L = \frac{45}{9} = 5$.}

\end{tcolorbox}
\vspace{6pt}\hrule height 0.35pt\vspace{6pt}
\noindent
\begin{tcolorbox}[
    blanker,
    sidebyside,
    sidebyside align=top,
    lefthand width=0.55\textwidth,
    top=4pt, bottom=4pt, left=0pt, right=0pt,
    middle=16pt,
    underlay={
        \draw[dashed, black!80, line width=0.8pt, dash pattern=on 4pt off 3pt]
            ([xshift=0.55\textwidth+8pt]frame.north west) --
            ([xshift=0.55\textwidth+8pt]frame.south west);
    }
]
\textbf{Q.91} \quad {\footnotesize\textsc{[NAT $\cdot$ 2 Mark(s) $\cdot$ No Neg]}} \hfill {\footnotesize\textbf{GATE 2024 CS1} \quad \tiny\texttt{2024\_CS1\_CS\_Q57}}
\vspace{2pt}\hrule height 0.3pt\vspace{4pt}

Consider the following code snippet using the `fork()` and `wait()` system calls. Assume that the code compiles and runs correctly, and that the system calls run successfully without any errors.
\begin{lstlisting}[language=C]
int x = 3;
while(x > 0) {
    fork();
    printf("hello");
    wait(NULL);
    x--;
}
\end{lstlisting}
The total number of times the `printf` statement is executed is \_\_\_\_\_\_\_\. (Answer in integer)

\par\vspace{6pt}\noindent\textbf{Your Answer:} \fbox{\rule{0pt}{1.3ex}\hspace{3cm}}

\tcblower

\textbf{Official Key:} \textbf{\large 14 to 14} \hfill \textit{\small [Medium]}
\vspace{2pt}\hrule height 0.3pt\vspace{3pt}
\textbf{Topic:} Operating System \\
\textbf{Subtopic:} \textit{System calls, processes, threads, inter-process communication}
\vspace{2pt}\hrule height 0.3pt\vspace{3pt}
\textbf{Guide / Hint:} \\
{\small Initially 1 process with $x=3$. Iteration 1 ($x=3$): `fork()` creates 1 child (total 2 processes); both print "hello" (2 prints) and decrement $x$ to 2. Iteration 2 ($x=2$): each of the 2 processes forks (total 4 processes); all 4 print "hello" (4 prints) and decrement $x$ to 1. Iteration 3 ($x=1$): each of the 4 processes forks (total 8 processes); all 8 print "hello" (8 prints) and decrement $x$ to 0, whereupon the loop ends. Total prints $= 2 + 4 + 8 = 14$.}

\end{tcolorbox}
\vspace{6pt}\hrule height 0.35pt\vspace{6pt}
\noindent
\begin{tcolorbox}[
    blanker,
    sidebyside,
    sidebyside align=top,
    lefthand width=0.55\textwidth,
    top=4pt, bottom=4pt, left=0pt, right=0pt,
    middle=16pt,
    underlay={
        \draw[dashed, black!80, line width=0.8pt, dash pattern=on 4pt off 3pt]
            ([xshift=0.55\textwidth+8pt]frame.north west) --
            ([xshift=0.55\textwidth+8pt]frame.south west);
    }
]
\textbf{Q.92} \quad {\footnotesize\textsc{[NAT $\cdot$ 2 Mark(s) $\cdot$ No Neg]}} \hfill {\footnotesize\textbf{GATE 2024 CS1} \quad \tiny\texttt{2024\_CS1\_CS\_Q62}}
\vspace{2pt}\hrule height 0.3pt\vspace{4pt}

Consider a memory management system that uses a page size of 2 KB. Assume that both the physical and virtual addresses start from 0. Assume that the pages 0, 1, 2, and 3 are stored in the page frames 1, 3, 2, and 0, respectively. The physical address (in decimal format) corresponding to the virtual address 2500 (in decimal format) is \_\_\_\_\_\_\_\_\_\. (Answer in integer)

\par\vspace{6pt}\noindent\textbf{Your Answer:} \fbox{\rule{0pt}{1.3ex}\hspace{3cm}}

\tcblower

\textbf{Official Key:} \textbf{\large 6596 to 6596} \hfill \textit{\small [Easy]}
\vspace{2pt}\hrule height 0.3pt\vspace{3pt}
\textbf{Topic:} Operating System \\
\textbf{Subtopic:} \textit{Memory management and virtual memory}
\vspace{2pt}\hrule height 0.3pt\vspace{3pt}
\textbf{Guide / Hint:} \\
{\small With page size $2\text{ KB} = 2048\text{ bytes}$, virtual address $2500$ has virtual page number $\lfloor 2500 / 2048 \rfloor = 1$ and offset $2500 \pmod{2048} = 452$. Page 1 maps to frame 3. The physical address is $(\text{frame } 3 \times 2048) + \text{offset } 452 = 6144 + 452 = 6596$.}

\end{tcolorbox}
\vspace{6pt}\hrule height 0.35pt\vspace{6pt}
\noindent
\begin{tcolorbox}[
    blanker,
    sidebyside,
    sidebyside align=top,
    lefthand width=0.55\textwidth,
    top=4pt, bottom=4pt, left=0pt, right=0pt,
    middle=16pt,
    underlay={
        \draw[dashed, black!80, line width=0.8pt, dash pattern=on 4pt off 3pt]
            ([xshift=0.55\textwidth+8pt]frame.north west) --
            ([xshift=0.55\textwidth+8pt]frame.south west);
    }
]
\textbf{Q.93} \quad {\footnotesize\textsc{[NAT $\cdot$ 2 Mark(s) $\cdot$ No Neg]}} \hfill {\footnotesize\textbf{GATE 2024 CS2} \quad \tiny\texttt{2024\_CS2\_CS\_Q53}}
\vspace{2pt}\hrule height 0.3pt\vspace{4pt}

Consider a disk with the following specifications: rotation speed of 6000 RPM, average seek time of 5 milliseconds, 500 sectors/track, 512-byte sectors. A file has content stored in 3000 sectors located randomly on the disk. Assuming average rotational latency, the total time (in seconds, rounded off to 2 decimal places) to read the entire file from the disk is \_\_\_\_\_\_\. (Answer in range 29.50 to 30.50)

\par\vspace{6pt}\noindent\textbf{Your Answer:} \fbox{\rule{0pt}{1.3ex}\hspace{3cm}}

\tcblower

\textbf{Official Key:} \textbf{\large 29.50 to 30.50} \hfill \textit{\small [Medium]}
\vspace{2pt}\hrule height 0.3pt\vspace{3pt}
\textbf{Topic:} Operating System \\
\textbf{Subtopic:} \textit{File systems}
\vspace{2pt}\hrule height 0.3pt\vspace{3pt}
\textbf{Guide / Hint:} \\
{\small At 6000 RPM, one rotation takes $\frac{60}{6000} = 10\text{ ms}$. Average rotational latency $= 10/2 = 5\text{ ms}$. Transfer time per sector $= \frac{10\text{ ms}}{500} = 0.02\text{ ms}$. Time to access each random sector $= \text{seek} + \text{rotational latency} + \text{transfer} = 5 + 5 + 0.02 = 10.02\text{ ms}$. Total time for 3000 random sectors $= 3000 \times 10.02\text{ ms} = 30060\text{ ms} \approx 30.06$ seconds.}

\end{tcolorbox}
\vspace{6pt}\hrule height 0.35pt\vspace{6pt}
\noindent
\begin{tcolorbox}[
    blanker,
    sidebyside,
    sidebyside align=top,
    lefthand width=0.55\textwidth,
    top=4pt, bottom=4pt, left=0pt, right=0pt,
    middle=16pt,
    underlay={
        \draw[dashed, black!80, line width=0.8pt, dash pattern=on 4pt off 3pt]
            ([xshift=0.55\textwidth+8pt]frame.north west) --
            ([xshift=0.55\textwidth+8pt]frame.south west);
    }
]
\textbf{Q.94} \quad {\footnotesize\textsc{[NAT $\cdot$ 2 Mark(s) $\cdot$ No Neg]}} \hfill {\footnotesize\textbf{GATE 2024 CS2} \quad \tiny\texttt{2024\_CS2\_CS\_Q64}}
\vspace{2pt}\hrule height 0.3pt\vspace{4pt}

Consider a 32-bit system with 4 KB page size and page table entries of size 4 bytes each. Assume $1\text{ KB} = 2^{10}$ bytes. The OS uses a 2-level page table for memory management, with an outer page directory and an inner page table. The OS allocates 1 page for the outer page directory upon process creation. The OS uses demand paging when allocating memory for the inner page table. An active process accesses 2000 unique pages during its execution. Let $X$ denote the minimum and $Y$ denote the maximum number of pages across the two levels of the page table of the process. The value of $X + Y$ is \_\_\_\_\_\_\. (Answer in integer)

\par\vspace{6pt}\noindent\textbf{Your Answer:} \fbox{\rule{0pt}{1.3ex}\hspace{3cm}}

\tcblower

\textbf{Official Key:} \textbf{\large 1028 to 1028} \hfill \textit{\small [Hard]}
\vspace{2pt}\hrule height 0.3pt\vspace{3pt}
\textbf{Topic:} Operating System \\
\textbf{Subtopic:} \textit{Memory management and virtual memory}
\vspace{2pt}\hrule height 0.3pt\vspace{3pt}
\textbf{Guide / Hint:} \\
{\small Page size is $4\text{ KB} = 2^{12}$ bytes, PTE size is 4 bytes, so each page table holds $2^{12} / 4 = 1024$ entries. The process accesses 2000 unique pages. Minimum inner page tables $X$: $\lceil 2000 / 1024 \rceil = 2$ inner pages plus 1 outer page directory $= 3$. Maximum inner page tables $Y$: 2000 pages distributed across 2000 different inner tables, bounded by total possible entries in directory $= 1024$ inner pages plus 1 outer page directory $= 1025$. Hence, $X + Y = 3 + 1025 = 1028$.}

\end{tcolorbox}
\vspace{6pt}\hrule height 0.35pt\vspace{6pt}
\noindent
\begin{tcolorbox}[
    blanker,
    sidebyside,
    sidebyside align=top,
    lefthand width=0.55\textwidth,
    top=4pt, bottom=4pt, left=0pt, right=0pt,
    middle=16pt,
    underlay={
        \draw[dashed, black!80, line width=0.8pt, dash pattern=on 4pt off 3pt]
            ([xshift=0.55\textwidth+8pt]frame.north west) --
            ([xshift=0.55\textwidth+8pt]frame.south west);
    }
]
\textbf{Q.95} \quad {\footnotesize\textsc{[NAT $\cdot$ 1 Mark(s) $\cdot$ No Neg]}} \hfill {\footnotesize\textbf{GATE 2025 CS1} \quad \tiny\texttt{2025\_CS1\_CS\_Q29}}
\vspace{2pt}\hrule height 0.3pt\vspace{4pt}

Suppose in a multiprogramming environment, the following C program segment is executed. A process goes into I/O queue whenever an I/O related operation is performed. Assume that there will always be a context switch whenever a process requests for an I/O, and also whenever the process returns from an I/O. The number of times the process will enter the ready queue during its lifetime (not counting the time the process enters the ready queue when it is run initially) is \_\_\_\_\_\_\. (Answer in integer)
\par\medskip
\begin{verbatim}
int main()
{
    int x=0, i=0;
    scanf("%d", &x);
    for(i=0; i<20; i++)
    {
        x = x + 20;
        printf("%d\n", x);
    }
    return 0;
}
\end{verbatim}

\par\vspace{6pt}\noindent\textbf{Your Answer:} \fbox{\rule{0pt}{1.3ex}\hspace{3cm}}

\tcblower

\textbf{Official Key:} \textbf{\large 21 to 21} \hfill \textit{\small [Medium]}
\vspace{2pt}\hrule height 0.3pt\vspace{3pt}
\textbf{Topic:} Operating System \\
\textbf{Subtopic:} \textit{CPU and I/O scheduling}
\vspace{2pt}\hrule height 0.3pt\vspace{3pt}
\textbf{Guide / Hint:} \\
{\small There is 1 initial `scanf` and 20 iterations of `printf`, totaling $1 + 20 = 21$ I/O operations. Each time an I/O operation completes, the process transitions from the I/O (waiting) queue back into the ready queue. Excluding the initial program launch into the ready queue, the process enters the ready queue exactly 21 times.}

\end{tcolorbox}
\vspace{6pt}\hrule height 0.35pt\vspace{6pt}
\noindent
\begin{tcolorbox}[
    blanker,
    sidebyside,
    sidebyside align=top,
    lefthand width=0.55\textwidth,
    top=4pt, bottom=4pt, left=0pt, right=0pt,
    middle=16pt,
    underlay={
        \draw[dashed, black!80, line width=0.8pt, dash pattern=on 4pt off 3pt]
            ([xshift=0.55\textwidth+8pt]frame.north west) --
            ([xshift=0.55\textwidth+8pt]frame.south west);
    }
]
\textbf{Q.96} \quad {\footnotesize\textsc{[NAT $\cdot$ 2 Mark(s) $\cdot$ No Neg]}} \hfill {\footnotesize\textbf{GATE 2025 CS1} \quad \tiny\texttt{2025\_CS1\_CS\_Q51}}
\vspace{2pt}\hrule height 0.3pt\vspace{4pt}

A disk of size 512M bytes is divided into blocks of 64K bytes. A file is stored in the disk using linked allocation. In linked allocation, each data block reserves 4 bytes to store the pointer to the next data block. The link part of the last data block contains a NULL pointer (also of 4 bytes). Suppose a file of 1M bytes needs to be stored in the disk. Assume $1\text{K} = 2^{10}$ and $1\text{M} = 2^{20}$. The amount of space in bytes that will be wasted due to internal fragmentation is \_\_\_\_\_\_\. (Answer in integer)

\par\vspace{6pt}\noindent\textbf{Your Answer:} \fbox{\rule{0pt}{1.3ex}\hspace{3cm}}

\tcblower

\textbf{Official Key:} \textbf{\large 65468 to 65468} \hfill \textit{\small [Medium]}
\vspace{2pt}\hrule height 0.3pt\vspace{3pt}
\textbf{Topic:} Operating System \\
\textbf{Subtopic:} \textit{File systems}
\vspace{2pt}\hrule height 0.3pt\vspace{3pt}
\textbf{Guide / Hint:} \\
{\small Block size is $64\text{ KB} = 65536$ bytes, with 4 bytes for pointer, leaving $65536 - 4 = 65532$ bytes of file data per block. File size is $1\text{ MB} = 1048576$ bytes. Number of blocks needed is $\lceil 1048576 / 65532 \rceil = \lceil 16.001037 \rceil = 17$ blocks. The first 16 blocks store $16 \times 65532 = 1048512$ bytes, leaving $1048576 - 1048512 = 64$ bytes stored in the 17th block. Internal fragmentation in the last block is $65532 - 64 = 65468$ bytes.}

\end{tcolorbox}
\vspace{6pt}\hrule height 0.35pt\vspace{6pt}
\noindent
\begin{tcolorbox}[
    blanker,
    sidebyside,
    sidebyside align=top,
    lefthand width=0.55\textwidth,
    top=4pt, bottom=4pt, left=0pt, right=0pt,
    middle=16pt,
    underlay={
        \draw[dashed, black!80, line width=0.8pt, dash pattern=on 4pt off 3pt]
            ([xshift=0.55\textwidth+8pt]frame.north west) --
            ([xshift=0.55\textwidth+8pt]frame.south west);
    }
]
\textbf{Q.97} \quad {\footnotesize\textsc{[NAT $\cdot$ 2 Mark(s) $\cdot$ No Neg]}} \hfill {\footnotesize\textbf{GATE 2025 CS1} \quad \tiny\texttt{2025\_CS1\_CS\_Q54}}
\vspace{2pt}\hrule height 0.3pt\vspace{4pt}

In optimal page replacement algorithm, information about all future page references is available to the operating system (OS). A modification of the optimal page replacement algorithm is as follows:
The OS correctly predicts only up to next 4 page references (including the current page) at the time of allocating a frame to a page.
A process accesses the pages in the following order of page numbers:
1, 3, 2, 4, 2, 3, 1, 2, 4, 3, 1, 4.
If the system has three memory frames that are initially empty, the number of page faults that will occur during execution of the process is \_\_\_\_\_\_\. (Answer in integer)

\par\vspace{6pt}\noindent\textbf{Your Answer:} \fbox{\rule{0pt}{1.3ex}\hspace{3cm}}

\tcblower

\textbf{Official Key:} \textbf{\large 6 to 6} \hfill \textit{\small [Medium]}
\vspace{2pt}\hrule height 0.3pt\vspace{3pt}
\textbf{Topic:} Operating System \\
\textbf{Subtopic:} \textit{Memory management and virtual memory}
\vspace{2pt}\hrule height 0.3pt\vspace{3pt}
\textbf{Guide / Hint:} \\
{\small Simulating the 3 frames with 4-step lookahead: 1 (fault, [1]), 3 (fault, [1, 3]), 2 (fault, [1, 3, 2]). At page 4 (fault): window is (4, 2, 3, 1), farthest is 1, so replace 1 -> [4, 3, 2]. Pages 2, 3 hit. At page 1 (fault): window (1, 2, 4, 3), replace 3 -> [4, 1, 2]. Page 2, 4, 3 (fault): replace 1 -> [4, 3, 2]. Total faults = 3 compulsory + 3 replacement = 6 page faults.}

\end{tcolorbox}
\vspace{6pt}\hrule height 0.35pt\vspace{6pt}
\noindent
\begin{tcolorbox}[
    blanker,
    sidebyside,
    sidebyside align=top,
    lefthand width=0.55\textwidth,
    top=4pt, bottom=4pt, left=0pt, right=0pt,
    middle=16pt,
    underlay={
        \draw[dashed, black!80, line width=0.8pt, dash pattern=on 4pt off 3pt]
            ([xshift=0.55\textwidth+8pt]frame.north west) --
            ([xshift=0.55\textwidth+8pt]frame.south west);
    }
]
\textbf{Q.98} \quad {\footnotesize\textsc{[NAT $\cdot$ 2 Mark(s) $\cdot$ No Neg]}} \hfill {\footnotesize\textbf{GATE 2025 CS2} \quad \tiny\texttt{2025\_CS2\_CS\_Q58}}
\vspace{2pt}\hrule height 0.3pt\vspace{4pt}

A computer system supports a logical address space of $2^{32}$ bytes. It uses two-level hierarchical paging with a page size of 4 KB and page table entry size of 4 bytes. The number of bits in the outer page table index is \_\_\_\_\_\_\_\_.

\par\vspace{6pt}\noindent\textbf{Your Answer:} \fbox{\rule{0pt}{1.3ex}\hspace{3cm}}

\tcblower

\textbf{Official Key:} \textbf{\large 11 to 11} \hfill \textit{\small [Hard]}
\vspace{2pt}\hrule height 0.3pt\vspace{3pt}
\textbf{Topic:} Operating System \\
\textbf{Subtopic:} \textit{Memory management and virtual memory: hierarchical paging}
\vspace{2pt}\hrule height 0.3pt\vspace{3pt}
\textbf{Guide / Hint:} \\
{\small For a $2^{32}$-byte logical address space with two-level hierarchical paging, page offset and page table index bits are derived from page size and entry size, yielding 11 bits for the first-level directory index.}

\end{tcolorbox}
\vspace{6pt}\hrule height 0.35pt\vspace{6pt}
\vspace{8pt}
\subsection*{Databases}
\vspace{-4pt}
\noindent
\begin{tcolorbox}[
    blanker,
    sidebyside,
    sidebyside align=top,
    lefthand width=0.55\textwidth,
    top=4pt, bottom=4pt, left=0pt, right=0pt,
    middle=16pt,
    underlay={
        \draw[dashed, black!80, line width=0.8pt, dash pattern=on 4pt off 3pt]
            ([xshift=0.55\textwidth+8pt]frame.north west) --
            ([xshift=0.55\textwidth+8pt]frame.south west);
    }
]
\textbf{Q.99} \quad {\footnotesize\textsc{[NAT $\cdot$ 1 Mark(s) $\cdot$ No Neg]}} \hfill {\footnotesize\textbf{GATE 2022 CS} \quad \tiny\texttt{2022\_CS\_CS\_Q31}}
\vspace{2pt}\hrule height 0.3pt\vspace{4pt}

Consider a relation $R(A, B, C, D, E)$ with the following three functional dependencies:
\[ AB \to C; \quad BC \to D; \quad C \to E; \]
The number of superkeys in the relation $R$ is \_\_\_\_\_\_\_\_\_\_\_\. (Answer in integer)

\par\vspace{6pt}\noindent\textbf{Your Answer:} \fbox{\rule{0pt}{1.3ex}\hspace{3cm}}

\tcblower

\textbf{Official Key:} \textbf{\large 8 to 8} \hfill \textit{\small [Easy]}
\vspace{2pt}\hrule height 0.3pt\vspace{3pt}
\textbf{Topic:} Databases \\
\textbf{Subtopic:} \textit{Integrity constraints, normal forms}
\vspace{2pt}\hrule height 0.3pt\vspace{3pt}
\textbf{Guide / Hint:} \\
{\small Computing attribute closures: $(AB)^+ = \{A, B, C, D, E\}$. Attributes $A$ and $B$ never appear on any RHS, so $AB$ is the unique candidate key. Every superkey must contain $AB$, and can include any subset of the remaining 3 attributes $\{C, D, E\}$, giving $2^3 = 8$ superkeys.}

\end{tcolorbox}
\vspace{6pt}\hrule height 0.35pt\vspace{6pt}
\noindent
\begin{tcolorbox}[
    blanker,
    sidebyside,
    sidebyside align=top,
    lefthand width=0.55\textwidth,
    top=4pt, bottom=4pt, left=0pt, right=0pt,
    middle=16pt,
    underlay={
        \draw[dashed, black!80, line width=0.8pt, dash pattern=on 4pt off 3pt]
            ([xshift=0.55\textwidth+8pt]frame.north west) --
            ([xshift=0.55\textwidth+8pt]frame.south west);
    }
]
\textbf{Q.100} \quad {\footnotesize\textsc{[NAT $\cdot$ 2 Mark(s) $\cdot$ No Neg]}} \hfill {\footnotesize\textbf{GATE 2022 CS} \quad \tiny\texttt{2022\_CS\_CS\_Q56}}
\vspace{2pt}\hrule height 0.3pt\vspace{4pt}

Consider the relational database with the following four schemas and their respective instances:
\[ \text{Student}(sNo, sName, dNo), \quad \text{Dept}(dNo, dName), \quad \text{Course}(cNo, cName, dNo), \quad \text{Register}(sNo, cNo) \]
\begin{lstlisting}[language=SQL]
SELECT * FROM Student AS S WHERE NOT EXISTS
    (SELECT cNo FROM Course WHERE dNo = 'D01'
     EXCEPT
     SELECT cNo FROM Register WHERE sNo = S.sNo)
\end{lstlisting}
The number of rows returned by the above SQL query is \_\_\_\_\_\_\_\_\_\_\_\. (Answer in integer)

\par\vspace{6pt}\noindent\textbf{Your Answer:} \fbox{\rule{0pt}{1.3ex}\hspace{3cm}}

\tcblower

\textbf{Official Key:} \textbf{\large 2 to 2} \hfill \textit{\small [Medium]}
\vspace{2pt}\hrule height 0.3pt\vspace{3pt}
\textbf{Topic:} Databases \\
\textbf{Subtopic:} \textit{Relational model: relational algebra, tuple calculus, SQL}
\vspace{2pt}\hrule height 0.3pt\vspace{3pt}
\textbf{Guide / Hint:} \\
{\small The nested subquery with `EXCEPT` and `NOT EXISTS` expresses relational division: it finds students who have registered for every course offered by department 'D01'. In the database instance, exactly 2 students satisfy this condition.}

\end{tcolorbox}
\vspace{6pt}\hrule height 0.35pt\vspace{6pt}
\noindent
\begin{tcolorbox}[
    blanker,
    sidebyside,
    sidebyside align=top,
    lefthand width=0.55\textwidth,
    top=4pt, bottom=4pt, left=0pt, right=0pt,
    middle=16pt,
    underlay={
        \draw[dashed, black!80, line width=0.8pt, dash pattern=on 4pt off 3pt]
            ([xshift=0.55\textwidth+8pt]frame.north west) --
            ([xshift=0.55\textwidth+8pt]frame.south west);
    }
]
\textbf{Q.101} \quad {\footnotesize\textsc{[NAT $\cdot$ 2 Mark(s) $\cdot$ No Neg]}} \hfill {\footnotesize\textbf{GATE 2023 CS} \quad \tiny\texttt{2023\_CS\_CS\_Q61}}
\vspace{2pt}\hrule height 0.3pt\vspace{4pt}

Consider the following table named `Student` in a relational database. The primary key of this table is `rollNum`.
\begin{center}
\begin{tabular}{|c|c|c|c|}
\hline
\textbf{rollNum} & \textbf{name} & \textbf{gender} & \textbf{marks} \\
\hline
1 & Naman & M & 62 \\
2 & Aliya & F & 70 \\
3 & Aliya & F & 80 \\
4 & James & M & 82 \\
5 & Swati & F & 65 \\
\hline
\end{tabular}
\end{center}
The SQL query below is executed on this database:
\begin{lstlisting}[language=SQL]
SELECT *
FROM Student
WHERE gender = 'F' AND marks > 65;
\end{lstlisting}
The number of rows returned by the query is \_\_\_\_\_\_\_\_\. (Answer in integer)

\par\vspace{6pt}\noindent\textbf{Your Answer:} \fbox{\rule{0pt}{1.3ex}\hspace{3cm}}

\tcblower

\textbf{Official Key:} \textbf{\large 2 to 2} \hfill \textit{\small [Easy]}
\vspace{2pt}\hrule height 0.3pt\vspace{3pt}
\textbf{Topic:} Databases \\
\textbf{Subtopic:} \textit{Relational model: relational algebra, tuple calculus, SQL}
\vspace{2pt}\hrule height 0.3pt\vspace{3pt}
\textbf{Guide / Hint:} \\
{\small The query evaluates the WHERE clause condition `gender = 'F' AND marks > 65`. Among female students: Aliya (rollNum 2, marks 70) and Aliya (rollNum 3, marks 80) satisfy both conditions, whereas Swati (rollNum 5, marks 65) fails the strict inequality. Thus, exactly 2 rows are returned.}

\end{tcolorbox}
\vspace{6pt}\hrule height 0.35pt\vspace{6pt}
\noindent
\begin{tcolorbox}[
    blanker,
    sidebyside,
    sidebyside align=top,
    lefthand width=0.55\textwidth,
    top=4pt, bottom=4pt, left=0pt, right=0pt,
    middle=16pt,
    underlay={
        \draw[dashed, black!80, line width=0.8pt, dash pattern=on 4pt off 3pt]
            ([xshift=0.55\textwidth+8pt]frame.north west) --
            ([xshift=0.55\textwidth+8pt]frame.south west);
    }
]
\textbf{Q.102} \quad {\footnotesize\textsc{[NAT $\cdot$ 1 Mark(s) $\cdot$ No Neg]}} \hfill {\footnotesize\textbf{GATE 2024 CS1} \quad \tiny\texttt{2024\_CS1\_CS\_Q35}}
\vspace{2pt}\hrule height 0.3pt\vspace{4pt}

Consider the following two relations, $R(A,B)$ and $S(A,C)$:
\begin{center}
\begin{tabular}{|c|c|}
\multicolumn{2}{c}{$R$} \\
\hline
$A$ & $B$ \\
\hline
10 & 20 \\
20 & 30 \\
30 & 40 \\
30 & 50 \\
\hline
\end{tabular}
\qquad
\begin{tabular}{|c|c|}
\multicolumn{2}{c}{$S$} \\
\hline
$A$ & $C$ \\
\hline
10 & 90 \\
30 & 45 \\
40 & 80 \\
50 & 95 \\
\hline
\end{tabular}
\end{center}
The total number of tuples obtained by evaluating the following expression
\[ \sigma\_{B < C}(R \bowtie\_{R.A = S.A} S) \]
is \_\_\_\_\_\_\_\_\_\. (Answer in integer)

\par\vspace{6pt}\noindent\textbf{Your Answer:} \fbox{\rule{0pt}{1.3ex}\hspace{3cm}}

\tcblower

\textbf{Official Key:} \textbf{\large 2 to 2} \hfill \textit{\small [Easy]}
\vspace{2pt}\hrule height 0.3pt\vspace{3pt}
\textbf{Topic:} Databases \\
\textbf{Subtopic:} \textit{Relational model: relational algebra, tuple calculus, SQL}
\vspace{2pt}\hrule height 0.3pt\vspace{3pt}
\textbf{Guide / Hint:} \\
{\small Joining $R(A,B)$ and $S(A,C)$ on $R.A = S.A$ yields: $(10, 20, 90)$ for $A=10$, and $(30, 40, 45)$, $(30, 50, 45)$ for $A=30$. Applying selection $\sigma_{B < C}$: $20 < 90$ is true (tuple $(10, 20, 90)$), $40 < 45$ is true (tuple $(30, 40, 45)$), and $50 < 45$ is false. Thus, exactly 2 tuples satisfy the condition.}

\end{tcolorbox}
\vspace{6pt}\hrule height 0.35pt\vspace{6pt}
\noindent
\begin{tcolorbox}[
    blanker,
    sidebyside,
    sidebyside align=top,
    lefthand width=0.55\textwidth,
    top=4pt, bottom=4pt, left=0pt, right=0pt,
    middle=16pt,
    underlay={
        \draw[dashed, black!80, line width=0.8pt, dash pattern=on 4pt off 3pt]
            ([xshift=0.55\textwidth+8pt]frame.north west) --
            ([xshift=0.55\textwidth+8pt]frame.south west);
    }
]
\textbf{Q.103} \quad {\footnotesize\textsc{[NAT $\cdot$ 2 Mark(s) $\cdot$ No Neg]}} \hfill {\footnotesize\textbf{GATE 2024 CS2} \quad \tiny\texttt{2024\_CS2\_CS\_Q56}}
\vspace{2pt}\hrule height 0.3pt\vspace{4pt}

A functional dependency $F: X \to Y$ is termed as a useful functional dependency if and only if it satisfies all the following three conditions:
\begin{itemize}
\item $X$ is not the empty set.
\item $Y$ is not the empty set.
\item Intersection of $X$ and $Y$ is the empty set.
\end{itemize}
For a relation $R$ with 4 attributes, the total number of possible useful functional dependencies is \_\_\_\_\_\_\. (Answer in integer)

\par\vspace{6pt}\noindent\textbf{Your Answer:} \fbox{\rule{0pt}{1.3ex}\hspace{3cm}}

\tcblower

\textbf{Official Key:} \textbf{\large 50 to 50} \hfill \textit{\small [Medium]}
\vspace{2pt}\hrule height 0.3pt\vspace{3pt}
\textbf{Topic:} Databases \\
\textbf{Subtopic:} \textit{Integrity constraints, normal forms}
\vspace{2pt}\hrule height 0.3pt\vspace{3pt}
\textbf{Guide / Hint:} \\
{\small For each attribute $A$ among the 4 attributes, there are 3 mutually exclusive choices: $A \in X$, $A \in Y$, or $A \notin (X \cup Y)$, giving $3^4 = 81$ pairs of disjoint sets $(X, Y)$. We subtract cases where $X = \emptyset$ ($2^4 = 16$), where $Y = \emptyset$ ($2^4 = 16$), and add back the double-subtracted case where $X = Y = \emptyset$ ($1$). Total useful FDs $= 81 - 16 - 16 + 1 = 50$.}

\end{tcolorbox}
\vspace{6pt}\hrule height 0.35pt\vspace{6pt}
\noindent
\begin{tcolorbox}[
    blanker,
    sidebyside,
    sidebyside align=top,
    lefthand width=0.55\textwidth,
    top=4pt, bottom=4pt, left=0pt, right=0pt,
    middle=16pt,
    underlay={
        \draw[dashed, black!80, line width=0.8pt, dash pattern=on 4pt off 3pt]
            ([xshift=0.55\textwidth+8pt]frame.north west) --
            ([xshift=0.55\textwidth+8pt]frame.south west);
    }
]
\textbf{Q.104} \quad {\footnotesize\textsc{[NAT $\cdot$ 2 Mark(s) $\cdot$ No Neg]}} \hfill {\footnotesize\textbf{GATE 2025 CS2} \quad \tiny\texttt{2025\_CS2\_CS\_Q57}}
\vspace{2pt}\hrule height 0.3pt\vspace{4pt}

In a B+- tree where each node can hold at most four key values, a root to leaf path consists of nodes containing keys. What is the maximum number of keys in the tree under the given path constraints?

\par\vspace{6pt}\noindent\textbf{Your Answer:} \fbox{\rule{0pt}{1.3ex}\hspace{3cm}}

\tcblower

\textbf{Official Key:} \textbf{\large 33 to 33} \hfill \textit{\small [Hard]}
\vspace{2pt}\hrule height 0.3pt\vspace{3pt}
\textbf{Topic:} Databases \\
\textbf{Subtopic:} \textit{File organization, indexing: B and B+ trees}
\vspace{2pt}\hrule height 0.3pt\vspace{3pt}
\textbf{Guide / Hint:} \\
{\small In a B+ tree of order 5 (max 4 keys, max 5 children per internal node), min keys per non-root internal node is $\lceil 5/2 \rceil - 1 = 2$. Calculating the maximum or minimum number of keys/records traversed along root-to-leaf paths yields 33.}

\end{tcolorbox}
\vspace{6pt}\hrule height 0.35pt\vspace{6pt}
\vspace{8pt}
\subsection*{Computer Networks}
\vspace{-4pt}
\noindent
\begin{tcolorbox}[
    blanker,
    sidebyside,
    sidebyside align=top,
    lefthand width=0.55\textwidth,
    top=4pt, bottom=4pt, left=0pt, right=0pt,
    middle=16pt,
    underlay={
        \draw[dashed, black!80, line width=0.8pt, dash pattern=on 4pt off 3pt]
            ([xshift=0.55\textwidth+8pt]frame.north west) --
            ([xshift=0.55\textwidth+8pt]frame.south west);
    }
]
\textbf{Q.105} \quad {\footnotesize\textsc{[NAT $\cdot$ 1 Mark(s) $\cdot$ No Neg]}} \hfill {\footnotesize\textbf{GATE 2018 CS} \quad \tiny\texttt{2018\_CS\_GA\_Q25}}
\vspace{2pt}\hrule height 0.3pt\vspace{4pt}

Consider a 1 Gbps network with an MTU of 1500 bytes and a 32-bit sequence number space. The minimum time (in seconds, rounded to the nearest integer) after which the sequence numbers can wrap around is \_\_\_\_\_\_\_\_.

\par\vspace{6pt}\noindent\textbf{Your Answer:} \fbox{\rule{0pt}{1.3ex}\hspace{3cm}}

\tcblower

\textbf{Official Key:} \textbf{\large 34 to 34} \hfill \textit{\small [Medium]}
\vspace{2pt}\hrule height 0.3pt\vspace{3pt}
\textbf{Topic:} Computer Networks \\
\textbf{Subtopic:} \textit{Transport layer: flow control and congestion control, UDP, TCP, sockets}
\vspace{2pt}\hrule height 0.3pt\vspace{3pt}
\textbf{Guide / Hint:} \\
{\small A 32-bit sequence number space contains $2^{32} = 4{,}294{,}967{,}296$ bytes. At a transmission rate of $1\text{ Gbps} = \frac{10^9}{8} = 1.25 \times 10^8\text{ bytes/s}$, the wrap-around time is $\frac{2^{32}}{1.25 \times 10^8} = \frac{4294967296}{125000000} \approx 34.36$ seconds $\approx 34$ seconds.}

\end{tcolorbox}
\vspace{6pt}\hrule height 0.35pt\vspace{6pt}
\noindent
\begin{tcolorbox}[
    blanker,
    sidebyside,
    sidebyside align=top,
    lefthand width=0.55\textwidth,
    top=4pt, bottom=4pt, left=0pt, right=0pt,
    middle=16pt,
    underlay={
        \draw[dashed, black!80, line width=0.8pt, dash pattern=on 4pt off 3pt]
            ([xshift=0.55\textwidth+8pt]frame.north west) --
            ([xshift=0.55\textwidth+8pt]frame.south west);
    }
]
\textbf{Q.106} \quad {\footnotesize\textsc{[NAT $\cdot$ 2 Mark(s) $\cdot$ No Neg]}} \hfill {\footnotesize\textbf{GATE 2018 CS} \quad \tiny\texttt{2018\_CS\_GA\_Q54}}
\vspace{2pt}\hrule height 0.3pt\vspace{4pt}

Consider an IP packet with a length of 4,500 bytes that includes a 20-byte IPv4 header and a 40-byte TCP header. The packet is forwarded across an IP network with an MTU of 600 bytes. The value of the fragment offset field in the third fragment (in 8-byte units) is \_\_\_\_\_\_\_\_.

\par\vspace{6pt}\noindent\textbf{Your Answer:} \fbox{\rule{0pt}{1.3ex}\hspace{3cm}}

\tcblower

\textbf{Official Key:} \textbf{\large 144 to 144} \hfill \textit{\small [Medium]}
\vspace{2pt}\hrule height 0.3pt\vspace{3pt}
\textbf{Topic:} Computer Networks \\
\textbf{Subtopic:} \textit{Fragmentation and IP addressing, IPv4, CIDR notation}
\vspace{2pt}\hrule height 0.3pt\vspace{3pt}
\textbf{Guide / Hint:} \\
{\small Total payload to fragment $= 4500 - 20 = 4480$ bytes. With $\text{MTU} = 600$, maximum payload per fragment must be a multiple of 8: $\lfloor (600 - 20) / 8 \rfloor \times 8 = 72 \times 8 = 576$ bytes. Fragment 1 carries bytes 0 to 575 (offset 0). Fragment 2 carries bytes 576 to 1151 (offset $576/8 = 72$). Fragment 3 carries bytes 1152 onwards, so its fragment offset field is $\frac{1152}{8} = 144$.}

\end{tcolorbox}
\vspace{6pt}\hrule height 0.35pt\vspace{6pt}
\noindent
\begin{tcolorbox}[
    blanker,
    sidebyside,
    sidebyside align=top,
    lefthand width=0.55\textwidth,
    top=4pt, bottom=4pt, left=0pt, right=0pt,
    middle=16pt,
    underlay={
        \draw[dashed, black!80, line width=0.8pt, dash pattern=on 4pt off 3pt]
            ([xshift=0.55\textwidth+8pt]frame.north west) --
            ([xshift=0.55\textwidth+8pt]frame.south west);
    }
]
\textbf{Q.107} \quad {\footnotesize\textsc{[NAT $\cdot$ 1 Mark(s) $\cdot$ No Neg]}} \hfill {\footnotesize\textbf{GATE 2022 CS} \quad \tiny\texttt{2022\_CS\_CS\_Q35}}
\vspace{2pt}\hrule height 0.3pt\vspace{4pt}

Consider the resolution of the domain name `www.gate.org.in` by a DNS resolver. Assume that no resource records are cached anywhere across the DNS servers and that iterative query mechanism is used in the resolution. The number of DNS query-response pairs involved in completely resolving the domain name is \_\_\_\_\_\_\_\_\_\_\_\_\. (Answer in integer)

\par\vspace{6pt}\noindent\textbf{Your Answer:} \fbox{\rule{0pt}{1.3ex}\hspace{3cm}}

\tcblower

\textbf{Official Key:} \textbf{\large 4 to 4} \hfill \textit{\small [Easy]}
\vspace{2pt}\hrule height 0.3pt\vspace{3pt}
\textbf{Topic:} Computer Networks \\
\textbf{Subtopic:} \textit{Application layer protocols: DNS, SMTP, HTTP, FTP, Email}
\vspace{2pt}\hrule height 0.3pt\vspace{3pt}
\textbf{Guide / Hint:} \\
{\small In an un-cached iterative resolution for `www.gate.org.in`: the resolver queries Root DNS server (1), `.in` TLD server (2), `.org.in` SLD server (3), and `gate.org.in` Authoritative server (4). Each query yields a response, totaling 4 query-response pairs.}

\end{tcolorbox}
\vspace{6pt}\hrule height 0.35pt\vspace{6pt}
\noindent
\begin{tcolorbox}[
    blanker,
    sidebyside,
    sidebyside align=top,
    lefthand width=0.55\textwidth,
    top=4pt, bottom=4pt, left=0pt, right=0pt,
    middle=16pt,
    underlay={
        \draw[dashed, black!80, line width=0.8pt, dash pattern=on 4pt off 3pt]
            ([xshift=0.55\textwidth+8pt]frame.north west) --
            ([xshift=0.55\textwidth+8pt]frame.south west);
    }
]
\textbf{Q.108} \quad {\footnotesize\textsc{[NAT $\cdot$ 2 Mark(s) $\cdot$ No Neg]}} \hfill {\footnotesize\textbf{GATE 2022 CS} \quad \tiny\texttt{2022\_CS\_CS\_Q59}}
\vspace{2pt}\hrule height 0.3pt\vspace{4pt}

Consider a 100 Mbps link between an earth station (sender) and a satellite (receiver) at an altitude of 2100 km. The signal propagates at a speed of $3 \times 10^8\text{ m/s}$. The time taken (in milliseconds, rounded off to two decimal places) for the receiver to completely receive a packet of 1000 bytes transmitted by the sender is \_\_\_\_\_\_\_\_\_\. (Answer in numeric format)

\par\vspace{6pt}\noindent\textbf{Your Answer:} \fbox{\rule{0pt}{1.3ex}\hspace{3cm}}

\tcblower

\textbf{Official Key:} \textbf{\large 7.07 to 7.09} \hfill \textit{\small [Easy]}
\vspace{2pt}\hrule height 0.3pt\vspace{3pt}
\textbf{Topic:} Computer Networks \\
\textbf{Subtopic:} \textit{Data link layer: framing, error detection, Medium Access Control, Ethernet bridging}
\vspace{2pt}\hrule height 0.3pt\vspace{3pt}
\textbf{Guide / Hint:} \\
{\small Propagation time is $T_{\text{prop}} = \frac{2100 \times 10^3\text{ m}}{3 \times 10^8\text{ m/s}} = 7.00\text{ ms}$. Transmission time is $T_{\text{trans}} = \frac{1000 \times 8\text{ bits}}{100 \times 10^6\text{ bps}} = 0.08\text{ ms}$. Total time to completely receive the packet is $T_{\text{prop}} + T_{\text{trans}} = 7.00 + 0.08 = 7.08\text{ ms}$.}

\end{tcolorbox}
\vspace{6pt}\hrule height 0.35pt\vspace{6pt}
\noindent
\begin{tcolorbox}[
    blanker,
    sidebyside,
    sidebyside align=top,
    lefthand width=0.55\textwidth,
    top=4pt, bottom=4pt, left=0pt, right=0pt,
    middle=16pt,
    underlay={
        \draw[dashed, black!80, line width=0.8pt, dash pattern=on 4pt off 3pt]
            ([xshift=0.55\textwidth+8pt]frame.north west) --
            ([xshift=0.55\textwidth+8pt]frame.south west);
    }
]
\textbf{Q.109} \quad {\footnotesize\textsc{[NAT $\cdot$ 2 Mark(s) $\cdot$ No Neg]}} \hfill {\footnotesize\textbf{GATE 2022 CS} \quad \tiny\texttt{2022\_CS\_CS\_Q60}}
\vspace{2pt}\hrule height 0.3pt\vspace{4pt}

Consider the data transfer using TCP over a 1 Gbps link. Assuming that the maximum segment lifetime (MSL) is set to 60 seconds, the minimum number of bits required for the sequence number field of the TCP header, to prevent the sequence number space from wrapping around during the MSL is \_\_\_\_\_\_\_\_\_\_\_\_\. (Answer in integer)

\par\vspace{6pt}\noindent\textbf{Your Answer:} \fbox{\rule{0pt}{1.3ex}\hspace{3cm}}

\tcblower

\textbf{Official Key:} \textbf{\large 33 to 33} \hfill \textit{\small [Medium]}
\vspace{2pt}\hrule height 0.3pt\vspace{3pt}
\textbf{Topic:} Computer Networks \\
\textbf{Subtopic:} \textit{Transport layer: flow control and congestion control, UDP, TCP, sockets}
\vspace{2pt}\hrule height 0.3pt\vspace{3pt}
\textbf{Guide / Hint:} \\
{\small Data transmitted at 1 Gbps line rate in $\text{MSL} = 60\text{ s}$ is $\frac{10^9\text{ bps}}{8} \times 60\text{ s} = 7.5 \times 10^9\text{ bytes}$. The sequence number space must satisfy $2^k \ge 7.5 \times 10^9$. Since $2^{32} \approx 4.29 \times 10^9 < 7.5 \times 10^9$ and $2^{33} \approx 8.59 \times 10^9 > 7.5 \times 10^9$, at least 33 bits are required.}

\end{tcolorbox}
\vspace{6pt}\hrule height 0.35pt\vspace{6pt}
\noindent
\begin{tcolorbox}[
    blanker,
    sidebyside,
    sidebyside align=top,
    lefthand width=0.55\textwidth,
    top=4pt, bottom=4pt, left=0pt, right=0pt,
    middle=16pt,
    underlay={
        \draw[dashed, black!80, line width=0.8pt, dash pattern=on 4pt off 3pt]
            ([xshift=0.55\textwidth+8pt]frame.north west) --
            ([xshift=0.55\textwidth+8pt]frame.south west);
    }
]
\textbf{Q.110} \quad {\footnotesize\textsc{[NAT $\cdot$ 2 Mark(s) $\cdot$ No Neg]}} \hfill {\footnotesize\textbf{GATE 2023 CS} \quad \tiny\texttt{2023\_CS\_CS\_Q65}}
\vspace{2pt}\hrule height 0.3pt\vspace{4pt}

The forwarding table of a router is shown below:
\begin{center}
\begin{tabular}{|l|l|c|}
\hline
\textbf{Subnet Number} & \textbf{Subnet Mask} & \textbf{Interface ID} \\
\hline
200.150.0.0 & 255.255.0.0 & 1 \\
200.150.64.0 & 255.255.224.0 & 2 \\
200.150.68.0 & 255.255.255.0 & 3 \\
200.150.68.64 & 255.255.255.224 & 4 \\
Default & & 0 \\
\hline
\end{tabular}
\end{center}
A packet addressed to a destination address 200.150.68.118 arrives at the router. It will be forwarded to the interface with ID \_\_\_\_\_\_\_\_\. (Answer in integer)

\par\vspace{6pt}\noindent\textbf{Your Answer:} \fbox{\rule{0pt}{1.3ex}\hspace{3cm}}

\tcblower

\textbf{Official Key:} \textbf{\large 3 to 3} \hfill \textit{\small [Medium]}
\vspace{2pt}\hrule height 0.3pt\vspace{3pt}
\textbf{Topic:} Computer Networks \\
\textbf{Subtopic:} \textit{Routing protocols: shortest path, flooding, distance vector and link state routing}
\vspace{2pt}\hrule height 0.3pt\vspace{3pt}
\textbf{Guide / Hint:} \\
{\small For destination 200.150.68.118: subnet /16 (prefix length 16) matches. Subnet 200.150.64.0/19 (mask 255.255.224.0) matches since $68 \& 224 = 64$. Subnet 200.150.68.0/24 (mask 255.255.255.0) matches since $68 = 68$. Subnet 200.150.68.64/27 (mask 255.255.255.224) does not match since $118 \& 224 = 96 \neq 64$. By the Longest Prefix Match rule, the /24 route (interface 3) is selected.}

\end{tcolorbox}
\vspace{6pt}\hrule height 0.35pt\vspace{6pt}
\noindent
\begin{tcolorbox}[
    blanker,
    sidebyside,
    sidebyside align=top,
    lefthand width=0.55\textwidth,
    top=4pt, bottom=4pt, left=0pt, right=0pt,
    middle=16pt,
    underlay={
        \draw[dashed, black!80, line width=0.8pt, dash pattern=on 4pt off 3pt]
            ([xshift=0.55\textwidth+8pt]frame.north west) --
            ([xshift=0.55\textwidth+8pt]frame.south west);
    }
]
\textbf{Q.111} \quad {\footnotesize\textsc{[NAT $\cdot$ 2 Mark(s) $\cdot$ No Neg]}} \hfill {\footnotesize\textbf{GATE 2024 CS1} \quad \tiny\texttt{2024\_CS1\_CS\_Q58}}
\vspace{2pt}\hrule height 0.3pt\vspace{4pt}

Consider the entries shown below in the forwarding table of an IP router. Each entry consists of an IP prefix and the corresponding next hop router for packets whose destination IP address matches the prefix. The notation ``/N'' in a prefix indicates a subnet mask with the most significant $N$ bits set to 1.
\begin{center}
\begin{tabular}{|l|c|}
\hline
\textbf{Prefix} & \textbf{Next hop router} \\
\hline
10.1.1.0/24 & R1 \\
10.1.1.128/25 & R2 \\
10.1.1.64/26 & R3 \\
10.1.1.192/26 & R4 \\
\hline
\end{tabular}
\end{center}
This router forwards 20 packets each to 5 hosts. The IP addresses of the hosts are 10.1.1.16, 10.1.1.72, 10.1.1.132, 10.1.1.191, and 10.1.1.205. The number of packets forwarded via the next hop router R2 is \_\_\_\_\_\_\_\. (Answer in integer)

\par\vspace{6pt}\noindent\textbf{Your Answer:} \fbox{\rule{0pt}{1.3ex}\hspace{3cm}}

\tcblower

\textbf{Official Key:} \textbf{\large 40 to 40} \hfill \textit{\small [Medium]}
\vspace{2pt}\hrule height 0.3pt\vspace{3pt}
\textbf{Topic:} Computer Networks \\
\textbf{Subtopic:} \textit{Routing protocols: shortest path, flooding, distance vector and link state routing}
\vspace{2pt}\hrule height 0.3pt\vspace{3pt}
\textbf{Guide / Hint:} \\
{\small Using longest prefix match for each host address: 10.1.1.16 matches /24 (R1); 10.1.1.72 matches /24 and 10.1.1.64/26 (R3); 10.1.1.132 matches /24 and 10.1.1.128/25 (R2); 10.1.1.191 matches /24 and 10.1.1.128/25 (R2); 10.1.1.205 matches /24, /25, and 10.1.1.192/26 (R4). Thus, 2 hosts (10.1.1.132 and 10.1.1.191) match router R2. Total packets forwarded via R2 $= 2 \times 20 = 40$.}

\end{tcolorbox}
\vspace{6pt}\hrule height 0.35pt\vspace{6pt}
\noindent
\begin{tcolorbox}[
    blanker,
    sidebyside,
    sidebyside align=top,
    lefthand width=0.55\textwidth,
    top=4pt, bottom=4pt, left=0pt, right=0pt,
    middle=16pt,
    underlay={
        \draw[dashed, black!80, line width=0.8pt, dash pattern=on 4pt off 3pt]
            ([xshift=0.55\textwidth+8pt]frame.north west) --
            ([xshift=0.55\textwidth+8pt]frame.south west);
    }
]
\textbf{Q.112} \quad {\footnotesize\textsc{[NAT $\cdot$ 2 Mark(s) $\cdot$ No Neg]}} \hfill {\footnotesize\textbf{GATE 2024 CS1} \quad \tiny\texttt{2024\_CS1\_CS\_Q65}}
\vspace{2pt}\hrule height 0.3pt\vspace{4pt}

Consider sending an IP datagram of size 1420 bytes (including 20 bytes of IP header) from a sender to a receiver over a path of two links with a router between them. The first link (sender to router) has an MTU (Maximum Transmission Unit) size of 542 bytes, while the second link (router to receiver) has an MTU size of 360 bytes. The number of fragments that would be delivered at the receiver is \_\_\_\_\_\_\_\_\. (Answer in integer)

\par\vspace{6pt}\noindent\textbf{Your Answer:} \fbox{\rule{0pt}{1.3ex}\hspace{3cm}}

\tcblower

\textbf{Official Key:} \textbf{\large 6 to 6} \hfill \textit{\small [Medium]}
\vspace{2pt}\hrule height 0.3pt\vspace{3pt}
\textbf{Topic:} Computer Networks \\
\textbf{Subtopic:} \textit{Fragmentation and IP addressing, IPv4, CIDR notation}
\vspace{2pt}\hrule height 0.3pt\vspace{3pt}
\textbf{Guide / Hint:} \\
{\small The datagram has 20 bytes IP header and 1400 bytes data. Link 1 has $\text{MTU} = 542$ bytes (max data per fragment is $\lfloor (542-20)/8 \rfloor \times 8 = 520$ bytes), producing 3 fragments with data sizes 520, 520, and 360 bytes. Link 2 has $\text{MTU} = 360$ bytes (max data is $\lfloor (360-20)/8 \rfloor \times 8 = 336$ bytes). At router Q, fragment 1 (520 bytes) is fragmented into 2 ($336 + 184$), fragment 2 (520 bytes) is fragmented into 2 ($336 + 184$), and fragment 3 (360 bytes) is fragmented into 2 ($336 + 24$). Total fragments delivered at the receiver $= 2 + 2 + 2 = 6$.}

\end{tcolorbox}
\vspace{6pt}\hrule height 0.35pt\vspace{6pt}
\noindent
\begin{tcolorbox}[
    blanker,
    sidebyside,
    sidebyside align=top,
    lefthand width=0.55\textwidth,
    top=4pt, bottom=4pt, left=0pt, right=0pt,
    middle=16pt,
    underlay={
        \draw[dashed, black!80, line width=0.8pt, dash pattern=on 4pt off 3pt]
            ([xshift=0.55\textwidth+8pt]frame.north west) --
            ([xshift=0.55\textwidth+8pt]frame.south west);
    }
]
\textbf{Q.113} \quad {\footnotesize\textsc{[NAT $\cdot$ 2 Mark(s) $\cdot$ No Neg]}} \hfill {\footnotesize\textbf{GATE 2024 CS2} \quad \tiny\texttt{2024\_CS2\_CS\_Q54}}
\vspace{2pt}\hrule height 0.3pt\vspace{4pt}

Consider a TCP connection operating at a point of time with the congestion window of size 12 MSS (Maximum Segment Size), when a timeout occurs due to packet loss. Assuming that all the segments transmitted in the next two RTTs (Round Trip Time) are acknowledged correctly, the congestion window size (in MSS) during the third RTT will be \_\_\_\_\_\_\. (Answer in integer)

\par\vspace{6pt}\noindent\textbf{Your Answer:} \fbox{\rule{0pt}{1.3ex}\hspace{3cm}}

\tcblower

\textbf{Official Key:} \textbf{\large 4 to 4} \hfill \textit{\small [Medium]}
\vspace{2pt}\hrule height 0.3pt\vspace{3pt}
\textbf{Topic:} Computer Networks \\
\textbf{Subtopic:} \textit{Transport layer: flow control and congestion control, UDP, TCP, sockets}
\vspace{2pt}\hrule height 0.3pt\vspace{3pt}
\textbf{Guide / Hint:} \\
{\small When timeout occurs at $\text{cwnd} = 12\text{ MSS}$, slow start threshold is set to $\text{ssthresh} = \lfloor 12 / 2 \rfloor = 6\text{ MSS}$, and $\text{cwnd}$ resets to $1\text{ MSS}$ for the 1st RTT. After receiving ACKs in the 1st RTT, slow start doubles $\text{cwnd}$ to $2\text{ MSS}$ during the 2nd RTT. After receiving ACKs in the 2nd RTT, $\text{cwnd}$ doubles again to $4\text{ MSS}$ during the 3rd RTT.}

\end{tcolorbox}
\vspace{6pt}\hrule height 0.35pt\vspace{6pt}
\noindent
\begin{tcolorbox}[
    blanker,
    sidebyside,
    sidebyside align=top,
    lefthand width=0.55\textwidth,
    top=4pt, bottom=4pt, left=0pt, right=0pt,
    middle=16pt,
    underlay={
        \draw[dashed, black!80, line width=0.8pt, dash pattern=on 4pt off 3pt]
            ([xshift=0.55\textwidth+8pt]frame.north west) --
            ([xshift=0.55\textwidth+8pt]frame.south west);
    }
]
\textbf{Q.114} \quad {\footnotesize\textsc{[NAT $\cdot$ 2 Mark(s) $\cdot$ No Neg]}} \hfill {\footnotesize\textbf{GATE 2024 CS2} \quad \tiny\texttt{2024\_CS2\_CS\_Q55}}
\vspace{2pt}\hrule height 0.3pt\vspace{4pt}

Consider an Ethernet segment with a transmission speed of $10^8$ bits/sec and a maximum segment length of 500 meters. If the speed of propagation of the signal in the medium is $2 \times 10^8$ meters/sec, then the minimum frame size (in bits) required for collision detection is \_\_\_\_\_\_\. (Answer in integer)

\par\vspace{6pt}\noindent\textbf{Your Answer:} \fbox{\rule{0pt}{1.3ex}\hspace{3cm}}

\tcblower

\textbf{Official Key:} \textbf{\large 500 to 500} \hfill \textit{\small [Easy]}
\vspace{2pt}\hrule height 0.3pt\vspace{3pt}
\textbf{Topic:} Computer Networks \\
\textbf{Subtopic:} \textit{Data link layer: framing, error detection, Medium Access Control, Ethernet bridging}
\vspace{2pt}\hrule height 0.3pt\vspace{3pt}
\textbf{Guide / Hint:} \\
{\small Propagation delay $T_p = \frac{\text{Distance}}{\text{Velocity}} = \frac{500}{2 \times 10^8} = 2.5 \times 10^{-6}\text{ s} = 2.5\,\mu\text{s}$. For collision detection in CSMA/CD, transmission time must satisfy $T_t \ge 2 T_p = 5\,\mu\text{s}$. Minimum frame size $L = 2 T_p \times B = (5 \times 10^{-6}\text{ s}) \times 10^8\text{ bps} = 500$ bits.}

\end{tcolorbox}
\vspace{6pt}\hrule height 0.35pt\vspace{6pt}
\noindent
\begin{tcolorbox}[
    blanker,
    sidebyside,
    sidebyside align=top,
    lefthand width=0.55\textwidth,
    top=4pt, bottom=4pt, left=0pt, right=0pt,
    middle=16pt,
    underlay={
        \draw[dashed, black!80, line width=0.8pt, dash pattern=on 4pt off 3pt]
            ([xshift=0.55\textwidth+8pt]frame.north west) --
            ([xshift=0.55\textwidth+8pt]frame.south west);
    }
]
\textbf{Q.115} \quad {\footnotesize\textsc{[NAT $\cdot$ 2 Mark(s) $\cdot$ No Neg]}} \hfill {\footnotesize\textbf{GATE 2025 CS1} \quad \tiny\texttt{2025\_CS1\_CS\_Q56}}
\vspace{2pt}\hrule height 0.3pt\vspace{4pt}

Suppose a 5-bit message is transmitted from a source to a destination through a noisy channel. The probability that a bit of the message gets flipped during transmission is 0.01. Flipping of each bit is independent of one another. The probability that the message is delivered error-free to the destination is \_\_\_\_\_\_\. (rounded off to three decimal places)

\par\vspace{6pt}\noindent\textbf{Your Answer:} \fbox{\rule{0pt}{1.3ex}\hspace{3cm}}

\tcblower

\textbf{Official Key:} \textbf{\large 0.949 to 0.952} \hfill \textit{\small [Easy]}
\vspace{2pt}\hrule height 0.3pt\vspace{3pt}
\textbf{Topic:} Computer Networks \\
\textbf{Subtopic:} \textit{Data link layer: framing, error detection, Medium Access Control, Ethernet bridging}
\vspace{2pt}\hrule height 0.3pt\vspace{3pt}
\textbf{Guide / Hint:} \\
{\small For each bit, the probability of correct transmission without flipping is $1 - 0.01 = 0.99$. Since all 5 bit transmissions are mutually independent, the probability that all 5 bits are transmitted error-free is $(0.99)^5 = 0.95099 \approx 0.951$.}

\end{tcolorbox}
\vspace{6pt}\hrule height 0.35pt\vspace{6pt}
\noindent
\begin{tcolorbox}[
    blanker,
    sidebyside,
    sidebyside align=top,
    lefthand width=0.55\textwidth,
    top=4pt, bottom=4pt, left=0pt, right=0pt,
    middle=16pt,
    underlay={
        \draw[dashed, black!80, line width=0.8pt, dash pattern=on 4pt off 3pt]
            ([xshift=0.55\textwidth+8pt]frame.north west) --
            ([xshift=0.55\textwidth+8pt]frame.south west);
    }
]
\textbf{Q.116} \quad {\footnotesize\textsc{[NAT $\cdot$ 2 Mark(s) $\cdot$ No Neg]}} \hfill {\footnotesize\textbf{GATE 2025 CS1} \quad \tiny\texttt{2025\_CS1\_CS\_Q57}}
\vspace{2pt}\hrule height 0.3pt\vspace{4pt}

Suppose a message of size 15000 bytes is transmitted from a source to a destination using IPv4 protocol via two routers as shown: Source $\to$ Router-1 ($\text{MTU} = 5000$ bytes) $\to$ Router-2 ($\text{MTU} = 3000$ bytes) $\to$ Destination. Each router has a defined maximum transmission unit (MTU) including the 20-byte IP header. The number of fragments that will be delivered to the destination is \_\_\_\_\_\_\. (Answer in integer)

\par\vspace{6pt}\noindent\textbf{Your Answer:} \fbox{\rule{0pt}{1.3ex}\hspace{3cm}}

\tcblower

\textbf{Official Key:} \textbf{\large 7 to 7} \hfill \textit{\small [Hard]}
\vspace{2pt}\hrule height 0.3pt\vspace{3pt}
\textbf{Topic:} Computer Networks \\
\textbf{Subtopic:} \textit{Fragmentation and IP addressing, IPv4, CIDR notation}
\vspace{2pt}\hrule height 0.3pt\vspace{3pt}
\textbf{Guide / Hint:} \\
{\small At Router-1, max payload per fragment is $\lfloor (5000 - 20) / 8 \rfloor \times 8 = 4976$ bytes. 15000 bytes payload is fragmented into payloads of 4976, 4976, 4976, and 72 bytes (4 fragments). At Router-2 ($\text{MTU} = 3000$, max payload $\lfloor 2980/8 \rfloor \times 8 = 2976$ bytes), each of the three 4976-byte fragments splits into 2 fragments ($2976 + 2000$), yielding $3 \times 2 = 6$ fragments. The 72-byte fragment passes through intact. Total delivered fragments = $6 + 1 = 7$.}

\end{tcolorbox}
\vspace{6pt}\hrule height 0.35pt\vspace{6pt}
\end{document}